\section*{Lecture 7 - Look-ahead Methods}
\addcontentsline{toc}{section}{Lecture 7 - Look-ahead Methods}

\clearpage
\SlideHeader{Lecture 7 - Look-ahead Methods}{001}{Data Driven Decision Making}
\begin{minipage}[t]{0.405\textwidth}
\SlideImage{1}{../Materials/2026/3DM_07_Lookahead-Methods_SS26.pdf}
\begin{adjustbox}{max totalsize={\textwidth}{0.43\textheight},center}
\begin{minipage}{\textwidth}\scriptsize
\NoteSection{日本語訳}\begin{itemize}
\item データ駆動型意思決定
\item 講義 7 – 先読み Methods
\end{itemize}\NoteSection{試験対策ポイント}\begin{itemize}
\item Policy Function Approximation, Rolling Horizon, CFA, Look-ahead, VFAを何を近似するかで区別する。
\item 各手法の強みと弱みを1セットで説明する。
\end{itemize}
\end{minipage}
\end{adjustbox}
\end{minipage}\hfill
\begin{minipage}[t]{0.58\textwidth}
\begin{adjustbox}{max totalsize={\textwidth}{0.89\textheight},center}
\begin{minipage}{\textwidth}\scriptsize
\NoteSection{解説}このページでは「Data Driven Decision Making」を理解する。スライド上の表現は短いが、試験では定義、具体例、モデル上の意味、手法の限界まで説明できる必要がある。\par
Policy Function Approximation は、経験則やルールを直接方策として使う方法である。例えば『在庫が閾値を下回ったら注文する』『最も近いドライバーへ割り当てる』などである。計算は速く解釈しやすいが、複雑な将来影響を十分に扱えない可能性がある。\par
Rolling Horizon Re-Optimization は、現在時点で見えている情報を使って一定期間の最適化問題を解き、最初の一部だけ実行し、次時点でまた解き直す方法である。現実変化へ適応できるが、見えている範囲だけを最適化するため近視眼的になることがある。\par
Cost Function Approximation は、現在の目的関数や制約に補正項を加えて、将来に良い影響を持つ決定を誘導する方法である。例えば、将来需要が多い地域から車両を離さないようにペナルティを入れる。設計は問題依存で、補正項が適切でないと逆効果になる。\par
Look-ahead Methods は、将来シナリオを明示的に考え、現在決定の良し悪しを未来までシミュレーションまたは最適化して評価する方法である。将来を直接見るため直感的だが、計算負荷が大きく、サンプリングやホライズンの選び方に依存する。\par
Value Function Approximation は、状態または後決定状態の将来価値を特徴量とモデルで近似する方法である。学習した価値を使えば、各決定について遠い未来まで毎回シミュレーションしなくても将来影響を考慮できる。しかし良い特徴量と学習データが必要である。\par\NoteSection{確認問題と解答}\begin{enumerate}\item \textbf{L07-S001: 「Data Driven Decision Making」の概念を、定義・具体例・モデル上の意味の順に説明せよ。}\\定義としては、Policy Function Approximation, Rolling Horizon, CFA, Look-ahead, VFAを何を近似するかで区別する。。具体例では、配送・在庫・フードデリバリーのような現実問題を用い、状態、決定、外生情報、報酬、または情報モデル/意思決定モデルのどこに対応するかを明示する。モデル上の意味として、この概念が目的関数、制約、状態表現、将来価値、または方策選択にどのような影響を与えるかを書く。試験では、用語だけを列挙せず、入力データから意思決定、さらに現実の実装へつながる流れを文章で示すことが重要である。\item \textbf{L07-S001: このスライドの考え方を、講義全体のDDDMプロセスの中に位置づけよ。}\\DDDMプロセスでは、現実世界のデータを集約して情報モデルを作り、意思決定モデルまたは方策で行動を選び、実装後に結果を観測して次の状態へ進む。このスライドの概念は、そのどこか一部だけではなく、データ表現、意思決定、将来不確実性、計算可能性のいずれかに関わる。答案では、まずこの概念がどの段階に属するかを述べ、次に他の段階へどう影響するかを書く。例えば価値関数なら将来価値を現在決定へ戻す役割、FIFOなら旅行時間情報モデルの整合性を保証する役割、CFAなら目的関数を調整して将来に良い行動を誘導する役割である。\end{enumerate}
\end{minipage}
\end{adjustbox}
\end{minipage}


\clearpage
\SlideHeader{Lecture 7 - Look-ahead Methods}{002}{Policy Classes (Compare Powell 2011)}
\begin{minipage}[t]{0.405\textwidth}
\SlideImage{2}{../Materials/2026/3DM_07_Lookahead-Methods_SS26.pdf}
\begin{adjustbox}{max totalsize={\textwidth}{0.43\textheight},center}
\begin{minipage}{\textwidth}\scriptsize
\NoteSection{日本語訳}\begin{itemize}
\item 方策 Classes (Compare Powell 2011)
\item 方策関数近似
\item Rule-Based（この用語は講義スライド上の専門語として扱う）
\item ローリングホライズン再最適化
\item 費用関数近似
\item 先読み Models
\item 価値関数近似 Optimization-Based
\end{itemize}\NoteSection{試験対策ポイント}\begin{itemize}
\item Policy Function Approximation, Rolling Horizon, CFA, Look-ahead, VFAを何を近似するかで区別する。
\item 各手法の強みと弱みを1セットで説明する。
\end{itemize}
\end{minipage}
\end{adjustbox}
\end{minipage}\hfill
\begin{minipage}[t]{0.58\textwidth}
\begin{adjustbox}{max totalsize={\textwidth}{0.89\textheight},center}
\begin{minipage}{\textwidth}\scriptsize
\NoteSection{解説}このページでは「Policy Classes (Compare Powell 2011)」を理解する。スライド上の表現は短いが、試験では定義、具体例、モデル上の意味、手法の限界まで説明できる必要がある。\par
Policy Function Approximation は、経験則やルールを直接方策として使う方法である。例えば『在庫が閾値を下回ったら注文する』『最も近いドライバーへ割り当てる』などである。計算は速く解釈しやすいが、複雑な将来影響を十分に扱えない可能性がある。\par
Rolling Horizon Re-Optimization は、現在時点で見えている情報を使って一定期間の最適化問題を解き、最初の一部だけ実行し、次時点でまた解き直す方法である。現実変化へ適応できるが、見えている範囲だけを最適化するため近視眼的になることがある。\par
Cost Function Approximation は、現在の目的関数や制約に補正項を加えて、将来に良い影響を持つ決定を誘導する方法である。例えば、将来需要が多い地域から車両を離さないようにペナルティを入れる。設計は問題依存で、補正項が適切でないと逆効果になる。\par
Look-ahead Methods は、将来シナリオを明示的に考え、現在決定の良し悪しを未来までシミュレーションまたは最適化して評価する方法である。将来を直接見るため直感的だが、計算負荷が大きく、サンプリングやホライズンの選び方に依存する。\par
Value Function Approximation は、状態または後決定状態の将来価値を特徴量とモデルで近似する方法である。学習した価値を使えば、各決定について遠い未来まで毎回シミュレーションしなくても将来影響を考慮できる。しかし良い特徴量と学習データが必要である。\par\NoteSection{確認問題と解答}\begin{enumerate}\item \textbf{L07-S002: 「Policy Classes (Compare Powell 2011)」の概念を、定義・具体例・モデル上の意味の順に説明せよ。}\\定義としては、Policy Function Approximation, Rolling Horizon, CFA, Look-ahead, VFAを何を近似するかで区別する。。具体例では、配送・在庫・フードデリバリーのような現実問題を用い、状態、決定、外生情報、報酬、または情報モデル/意思決定モデルのどこに対応するかを明示する。モデル上の意味として、この概念が目的関数、制約、状態表現、将来価値、または方策選択にどのような影響を与えるかを書く。試験では、用語だけを列挙せず、入力データから意思決定、さらに現実の実装へつながる流れを文章で示すことが重要である。\item \textbf{L07-S002: このスライドの考え方を、講義全体のDDDMプロセスの中に位置づけよ。}\\DDDMプロセスでは、現実世界のデータを集約して情報モデルを作り、意思決定モデルまたは方策で行動を選び、実装後に結果を観測して次の状態へ進む。このスライドの概念は、そのどこか一部だけではなく、データ表現、意思決定、将来不確実性、計算可能性のいずれかに関わる。答案では、まずこの概念がどの段階に属するかを述べ、次に他の段階へどう影響するかを書く。例えば価値関数なら将来価値を現在決定へ戻す役割、FIFOなら旅行時間情報モデルの整合性を保証する役割、CFAなら目的関数を調整して将来に良い行動を誘導する役割である。\end{enumerate}
\end{minipage}
\end{adjustbox}
\end{minipage}


\clearpage
\SlideHeader{Lecture 7 - Look-ahead Methods}{003}{Policy Classes}
\begin{minipage}[t]{0.405\textwidth}
\SlideImage{3}{../Materials/2026/3DM_07_Lookahead-Methods_SS26.pdf}
\begin{adjustbox}{max totalsize={\textwidth}{0.43\textheight},center}
\begin{minipage}{\textwidth}\scriptsize
\NoteSection{日本語訳}\begin{itemize}
\item 方策クラス
\item 方策 function approximation:
\item Transfer a rule-of-thumb into a 方策
\item Rolling-Horizon Re-Optimization:（この用語は講義スライド上の専門語として扱う）
\item Determine 決定 based on current
\item information (myopic)（この用語は講義スライド上の専門語として扱う）
\item Cost-Function Approximation:（この用語は講義スライド上の専門語として扱う）
\item Manipulate 報酬 function or constraints to
\item incentivize (/avoid) (in)flexible 決定s
\end{itemize}\NoteSection{試験対策ポイント}\begin{itemize}
\item Policy Function Approximation, Rolling Horizon, CFA, Look-ahead, VFAを何を近似するかで区別する。
\item 各手法の強みと弱みを1セットで説明する。
\end{itemize}
\end{minipage}
\end{adjustbox}
\end{minipage}\hfill
\begin{minipage}[t]{0.58\textwidth}
\begin{adjustbox}{max totalsize={\textwidth}{0.89\textheight},center}
\begin{minipage}{\textwidth}\scriptsize
\NoteSection{解説}このページでは「Policy Classes」を理解する。スライド上の表現は短いが、試験では定義、具体例、モデル上の意味、手法の限界まで説明できる必要がある。\par
Policy Function Approximation は、経験則やルールを直接方策として使う方法である。例えば『在庫が閾値を下回ったら注文する』『最も近いドライバーへ割り当てる』などである。計算は速く解釈しやすいが、複雑な将来影響を十分に扱えない可能性がある。\par
Rolling Horizon Re-Optimization は、現在時点で見えている情報を使って一定期間の最適化問題を解き、最初の一部だけ実行し、次時点でまた解き直す方法である。現実変化へ適応できるが、見えている範囲だけを最適化するため近視眼的になることがある。\par
Cost Function Approximation は、現在の目的関数や制約に補正項を加えて、将来に良い影響を持つ決定を誘導する方法である。例えば、将来需要が多い地域から車両を離さないようにペナルティを入れる。設計は問題依存で、補正項が適切でないと逆効果になる。\par
Look-ahead Methods は、将来シナリオを明示的に考え、現在決定の良し悪しを未来までシミュレーションまたは最適化して評価する方法である。将来を直接見るため直感的だが、計算負荷が大きく、サンプリングやホライズンの選び方に依存する。\par
Value Function Approximation は、状態または後決定状態の将来価値を特徴量とモデルで近似する方法である。学習した価値を使えば、各決定について遠い未来まで毎回シミュレーションしなくても将来影響を考慮できる。しかし良い特徴量と学習データが必要である。\par\NoteSection{確認問題と解答}\begin{enumerate}\item \textbf{L07-S003: 「Policy Classes」の概念を、定義・具体例・モデル上の意味の順に説明せよ。}\\定義としては、Policy Function Approximation, Rolling Horizon, CFA, Look-ahead, VFAを何を近似するかで区別する。。具体例では、配送・在庫・フードデリバリーのような現実問題を用い、状態、決定、外生情報、報酬、または情報モデル/意思決定モデルのどこに対応するかを明示する。モデル上の意味として、この概念が目的関数、制約、状態表現、将来価値、または方策選択にどのような影響を与えるかを書く。試験では、用語だけを列挙せず、入力データから意思決定、さらに現実の実装へつながる流れを文章で示すことが重要である。\item \textbf{L07-S003: このスライドの考え方を、講義全体のDDDMプロセスの中に位置づけよ。}\\DDDMプロセスでは、現実世界のデータを集約して情報モデルを作り、意思決定モデルまたは方策で行動を選び、実装後に結果を観測して次の状態へ進む。このスライドの概念は、そのどこか一部だけではなく、データ表現、意思決定、将来不確実性、計算可能性のいずれかに関わる。答案では、まずこの概念がどの段階に属するかを述べ、次に他の段階へどう影響するかを書く。例えば価値関数なら将来価値を現在決定へ戻す役割、FIFOなら旅行時間情報モデルの整合性を保証する役割、CFAなら目的関数を調整して将来に良い行動を誘導する役割である。\end{enumerate}
\end{minipage}
\end{adjustbox}
\end{minipage}


\clearpage
\SlideHeader{Lecture 7 - Look-ahead Methods}{004}{Agenda}
\begin{minipage}[t]{0.405\textwidth}
\SlideImage{4}{../Materials/2026/3DM_07_Lookahead-Methods_SS26.pdf}
\begin{adjustbox}{max totalsize={\textwidth}{0.43\textheight},center}
\begin{minipage}{\textwidth}\scriptsize
\NoteSection{日本語訳}\begin{itemize}
\item アジェンダ
\item 複数シナリオアプローチ
\item 後決定状態ロールアウト
\end{itemize}
\end{minipage}
\end{adjustbox}
\end{minipage}\hfill
\begin{minipage}[t]{0.58\textwidth}
\begin{adjustbox}{max totalsize={\textwidth}{0.89\textheight},center}
\begin{minipage}{\textwidth}\scriptsize
\NoteSection{注記}このページは表紙・目次・事務情報・導入画像が中心であるため、無理に確認問題や過去問を付けない。
\end{minipage}
\end{adjustbox}
\end{minipage}


\clearpage
\SlideHeader{Lecture 7 - Look-ahead Methods}{005}{Sampling + Mixed Integer Programming}
\begin{minipage}[t]{0.405\textwidth}
\SlideImage{5}{../Materials/2026/3DM_07_Lookahead-Methods_SS26.pdf}
\begin{adjustbox}{max totalsize={\textwidth}{0.43\textheight},center}
\begin{minipage}{\textwidth}\scriptsize
\NoteSection{日本語訳}\begin{itemize}
\item サンプリング + 混合整数計画
\item 考え方: Instead of 状態 遷移s, sample stream of 確率的 information (for example,
\item customer requests, production orders, demand, etc.)（この用語は講義スライド上の専門語として扱う）
\item Extend 状態 information with sampled information
\item Solve mixed-integer program (最大化する 報酬)
\item Remove parts of solution that consider sampled information（この用語は講義スライド上の専門語として扱う）
\item Implement (partial) solution（この用語は講義スライド上の専門語として扱う）
\item 利点: Conventional MIP-solvers can be used
\item 欠点: Only one sample used for 決定 making (overfitting)
\end{itemize}\NoteSection{試験対策ポイント}\begin{itemize}
\item MSAは複数シナリオを解き、解の共通性から現在決定を選ぶ。
\item Consensus functionは問題構造に依存して設計する。
\end{itemize}
\end{minipage}
\end{adjustbox}
\end{minipage}\hfill
\begin{minipage}[t]{0.58\textwidth}
\begin{adjustbox}{max totalsize={\textwidth}{0.89\textheight},center}
\begin{minipage}{\textwidth}\scriptsize
\NoteSection{解説}このページでは「Sampling + Mixed Integer Programming」を理解する。スライド上の表現は短いが、試験では定義、具体例、モデル上の意味、手法の限界まで説明できる必要がある。\par
Look-ahead は、現在の決定を評価するために将来を仮想的に見る方法である。単一シナリオだけを見ると、そのシナリオに過剰適合する危険がある。そこで Multiple Scenario Approach では、複数の将来シナリオをサンプリングし、それぞれを確定的な問題として解く。\par
MSAの基本手順は、複数の将来情報列を生成し、各シナリオでMIPなどの確定的最適化を解き、実際にはまだ見えていない将来情報に依存する部分を取り除き、現在時点で実行可能な部分を比較する、という流れである。\par
Consensus function は、複数シナリオで得られた解の中から代表的な解を選ぶための類似度尺度である。生産スケジューリングなら次に処理するジョブ、VRPなら同じ顧客ペアが同じルートにあるか、TSPならアークや隣接関係が一致するか、施設配置なら選ばれた施設集合が一致するかで測る。\par
MSAの利点は、通常のMIPソルバーを使えることと、複数シナリオに対して極端に悪くない現在決定を選べることである。欠点は、stochasticityを解選択で間接的に扱うだけで、動的な再意思決定そのものを完全にモデル化しているわけではない点である。\par
試験では、MSAを『複数シナリオの平均を取る』だけと書かない。各シナリオの解を作り、その解の共通性や代表性を測り、現在実装すべき決定を選ぶ、という構造を説明する。Consensus functionは問題ごとに設計する必要がある。\par\NoteSection{確認問題と解答}\begin{enumerate}\item \textbf{L07-S005: 「Sampling + Mixed Integer Programming」の概念を、定義・具体例・モデル上の意味の順に説明せよ。}\\定義としては、MSAは複数シナリオを解き、解の共通性から現在決定を選ぶ。。具体例では、配送・在庫・フードデリバリーのような現実問題を用い、状態、決定、外生情報、報酬、または情報モデル/意思決定モデルのどこに対応するかを明示する。モデル上の意味として、この概念が目的関数、制約、状態表現、将来価値、または方策選択にどのような影響を与えるかを書く。試験では、用語だけを列挙せず、入力データから意思決定、さらに現実の実装へつながる流れを文章で示すことが重要である。\item \textbf{L07-S005: このスライドの考え方を、講義全体のDDDMプロセスの中に位置づけよ。}\\DDDMプロセスでは、現実世界のデータを集約して情報モデルを作り、意思決定モデルまたは方策で行動を選び、実装後に結果を観測して次の状態へ進む。このスライドの概念は、そのどこか一部だけではなく、データ表現、意思決定、将来不確実性、計算可能性のいずれかに関わる。答案では、まずこの概念がどの段階に属するかを述べ、次に他の段階へどう影響するかを書く。例えば価値関数なら将来価値を現在決定へ戻す役割、FIFOなら旅行時間情報モデルの整合性を保証する役割、CFAなら目的関数を調整して将来に良い行動を誘導する役割である。\end{enumerate}
\end{minipage}
\end{adjustbox}
\end{minipage}


\clearpage
\SlideHeader{Lecture 7 - Look-ahead Methods}{006}{Example: Production Planning}
\begin{minipage}[t]{0.405\textwidth}
\SlideImage{6}{../Materials/2026/3DM_07_Lookahead-Methods_SS26.pdf}
\begin{adjustbox}{max totalsize={\textwidth}{0.43\textheight},center}
\begin{minipage}{\textwidth}\scriptsize
\NoteSection{日本語訳}\begin{itemize}
\item 例: Production Planning
\item 状態
\item Sampling（この用語は講義スライド上の専門語として扱う）
\item A
\item Solve（この用語は講義スライド上の専門語として扱う）
\item B
\item Implement B（この用語は講義スライド上の専門語として扱う）
\end{itemize}\NoteSection{試験対策ポイント}\begin{itemize}
\item Policy Function Approximation, Rolling Horizon, CFA, Look-ahead, VFAを何を近似するかで区別する。
\item 各手法の強みと弱みを1セットで説明する。
\end{itemize}
\end{minipage}
\end{adjustbox}
\end{minipage}\hfill
\begin{minipage}[t]{0.58\textwidth}
\begin{adjustbox}{max totalsize={\textwidth}{0.89\textheight},center}
\begin{minipage}{\textwidth}\scriptsize
\NoteSection{解説}このページでは「Example: Production Planning」を理解する。スライド上の表現は短いが、試験では定義、具体例、モデル上の意味、手法の限界まで説明できる必要がある。\par
Policy Function Approximation は、経験則やルールを直接方策として使う方法である。例えば『在庫が閾値を下回ったら注文する』『最も近いドライバーへ割り当てる』などである。計算は速く解釈しやすいが、複雑な将来影響を十分に扱えない可能性がある。\par
Rolling Horizon Re-Optimization は、現在時点で見えている情報を使って一定期間の最適化問題を解き、最初の一部だけ実行し、次時点でまた解き直す方法である。現実変化へ適応できるが、見えている範囲だけを最適化するため近視眼的になることがある。\par
Cost Function Approximation は、現在の目的関数や制約に補正項を加えて、将来に良い影響を持つ決定を誘導する方法である。例えば、将来需要が多い地域から車両を離さないようにペナルティを入れる。設計は問題依存で、補正項が適切でないと逆効果になる。\par
Look-ahead Methods は、将来シナリオを明示的に考え、現在決定の良し悪しを未来までシミュレーションまたは最適化して評価する方法である。将来を直接見るため直感的だが、計算負荷が大きく、サンプリングやホライズンの選び方に依存する。\par
Value Function Approximation は、状態または後決定状態の将来価値を特徴量とモデルで近似する方法である。学習した価値を使えば、各決定について遠い未来まで毎回シミュレーションしなくても将来影響を考慮できる。しかし良い特徴量と学習データが必要である。\par\NoteSection{確認問題と解答}\begin{enumerate}\item \textbf{L07-S006: 「Example: Production Planning」の概念を、定義・具体例・モデル上の意味の順に説明せよ。}\\定義としては、Policy Function Approximation, Rolling Horizon, CFA, Look-ahead, VFAを何を近似するかで区別する。。具体例では、配送・在庫・フードデリバリーのような現実問題を用い、状態、決定、外生情報、報酬、または情報モデル/意思決定モデルのどこに対応するかを明示する。モデル上の意味として、この概念が目的関数、制約、状態表現、将来価値、または方策選択にどのような影響を与えるかを書く。試験では、用語だけを列挙せず、入力データから意思決定、さらに現実の実装へつながる流れを文章で示すことが重要である。\item \textbf{L07-S006: このスライドの考え方を、講義全体のDDDMプロセスの中に位置づけよ。}\\DDDMプロセスでは、現実世界のデータを集約して情報モデルを作り、意思決定モデルまたは方策で行動を選び、実装後に結果を観測して次の状態へ進む。このスライドの概念は、そのどこか一部だけではなく、データ表現、意思決定、将来不確実性、計算可能性のいずれかに関わる。答案では、まずこの概念がどの段階に属するかを述べ、次に他の段階へどう影響するかを書く。例えば価値関数なら将来価値を現在決定へ戻す役割、FIFOなら旅行時間情報モデルの整合性を保証する役割、CFAなら目的関数を調整して将来に良い行動を誘導する役割である。\end{enumerate}
\end{minipage}
\end{adjustbox}
\end{minipage}


\clearpage
\SlideHeader{Lecture 7 - Look-ahead Methods}{007}{Multiple Scenario Approach (MSA)}
\begin{minipage}[t]{0.405\textwidth}
\SlideImage{7}{../Materials/2026/3DM_07_Lookahead-Methods_SS26.pdf}
\begin{adjustbox}{max totalsize={\textwidth}{0.43\textheight},center}
\begin{minipage}{\textwidth}\scriptsize
\NoteSection{日本語訳}\begin{itemize}
\item Multiple Scenario Approach（MSA）
\item Bent, Van Hentenryck 2004（この用語は講義スライド上の専門語として扱う）
\item Two-stage 確率的 programming “light”
\item Sample set of scenarios, determine solution that fits all scenarios best（この用語は講義スライド上の専門語として扱う）
\item Sample several (streams of) information model realizations (e.g., sets of requests)（この用語は講義スライド上の専門語として扱う）
\item Solve deterministic MIPs individually（この用語は講義スライド上の専門語として扱う）
\item Clean solutions from augmented information (and repair, if necessary)（この用語は講義スライド上の専門語として扱う）
\item Compare solutions and select “most similar” solution（この用語は講義スライド上の専門語として扱う）
\item 利点: Conventional MIP-solver can be used
\end{itemize}\NoteSection{試験対策ポイント}\begin{itemize}
\item MSAは複数シナリオを解き、解の共通性から現在決定を選ぶ。
\item Consensus functionは問題構造に依存して設計する。
\end{itemize}
\end{minipage}
\end{adjustbox}
\end{minipage}\hfill
\begin{minipage}[t]{0.58\textwidth}
\begin{adjustbox}{max totalsize={\textwidth}{0.89\textheight},center}
\begin{minipage}{\textwidth}\scriptsize
\NoteSection{解説}このページでは「Multiple Scenario Approach (MSA)」を理解する。スライド上の表現は短いが、試験では定義、具体例、モデル上の意味、手法の限界まで説明できる必要がある。\par
Look-ahead は、現在の決定を評価するために将来を仮想的に見る方法である。単一シナリオだけを見ると、そのシナリオに過剰適合する危険がある。そこで Multiple Scenario Approach では、複数の将来シナリオをサンプリングし、それぞれを確定的な問題として解く。\par
MSAの基本手順は、複数の将来情報列を生成し、各シナリオでMIPなどの確定的最適化を解き、実際にはまだ見えていない将来情報に依存する部分を取り除き、現在時点で実行可能な部分を比較する、という流れである。\par
Consensus function は、複数シナリオで得られた解の中から代表的な解を選ぶための類似度尺度である。生産スケジューリングなら次に処理するジョブ、VRPなら同じ顧客ペアが同じルートにあるか、TSPならアークや隣接関係が一致するか、施設配置なら選ばれた施設集合が一致するかで測る。\par
MSAの利点は、通常のMIPソルバーを使えることと、複数シナリオに対して極端に悪くない現在決定を選べることである。欠点は、stochasticityを解選択で間接的に扱うだけで、動的な再意思決定そのものを完全にモデル化しているわけではない点である。\par
試験では、MSAを『複数シナリオの平均を取る』だけと書かない。各シナリオの解を作り、その解の共通性や代表性を測り、現在実装すべき決定を選ぶ、という構造を説明する。Consensus functionは問題ごとに設計する必要がある。\par\NoteSection{確認問題と解答}\begin{enumerate}\item \textbf{L07-S007: 「Multiple Scenario Approach (MSA)」の概念を、定義・具体例・モデル上の意味の順に説明せよ。}\\定義としては、MSAは複数シナリオを解き、解の共通性から現在決定を選ぶ。。具体例では、配送・在庫・フードデリバリーのような現実問題を用い、状態、決定、外生情報、報酬、または情報モデル/意思決定モデルのどこに対応するかを明示する。モデル上の意味として、この概念が目的関数、制約、状態表現、将来価値、または方策選択にどのような影響を与えるかを書く。試験では、用語だけを列挙せず、入力データから意思決定、さらに現実の実装へつながる流れを文章で示すことが重要である。\item \textbf{L07-S007: このスライドの考え方を、講義全体のDDDMプロセスの中に位置づけよ。}\\DDDMプロセスでは、現実世界のデータを集約して情報モデルを作り、意思決定モデルまたは方策で行動を選び、実装後に結果を観測して次の状態へ進む。このスライドの概念は、そのどこか一部だけではなく、データ表現、意思決定、将来不確実性、計算可能性のいずれかに関わる。答案では、まずこの概念がどの段階に属するかを述べ、次に他の段階へどう影響するかを書く。例えば価値関数なら将来価値を現在決定へ戻す役割、FIFOなら旅行時間情報モデルの整合性を保証する役割、CFAなら目的関数を調整して将来に良い行動を誘導する役割である。\end{enumerate}
\end{minipage}
\end{adjustbox}
\end{minipage}


\clearpage
\SlideHeader{Lecture 7 - Look-ahead Methods}{008}{MSA: Consensus Function}
\begin{minipage}[t]{0.405\textwidth}
\SlideImage{8}{../Materials/2026/3DM_07_Lookahead-Methods_SS26.pdf}
\begin{adjustbox}{max totalsize={\textwidth}{0.43\textheight},center}
\begin{minipage}{\textwidth}\scriptsize
\NoteSection{日本語訳}\begin{itemize}
\item MSA: 合意関数
\item The consensus function determines the “most similar” solution（この用語は講義スライド上の専門語として扱う）
\item Nice idea, but highly problem dependent（この用語は講義スライド上の専門語として扱う）
\item Consensus function for production planning:（この用語は講義スライド上の専門語として扱う）
\item Solution: Sequence of jobs per production line（この用語は講義スライド上の専門語として扱う）
\item Consider only next job each line produces（この用語は講義スライド上の専門語として扱う）
\item Similarity measure between two solutions: number of identical jobs（この用語は講義スライド上の専門語として扱う）
\item Select solution with maximum similarity（この用語は講義スライド上の専門語として扱う）
\item 例: 3 lines, A, B, C, six jobs 1,2,3,4,5,6
\end{itemize}\NoteSection{試験対策ポイント}\begin{itemize}
\item MSAは複数シナリオを解き、解の共通性から現在決定を選ぶ。
\item Consensus functionは問題構造に依存して設計する。
\end{itemize}
\end{minipage}
\end{adjustbox}
\end{minipage}\hfill
\begin{minipage}[t]{0.58\textwidth}
\begin{adjustbox}{max totalsize={\textwidth}{0.89\textheight},center}
\begin{minipage}{\textwidth}\scriptsize
\NoteSection{解説}このページでは「MSA: Consensus Function」を理解する。スライド上の表現は短いが、試験では定義、具体例、モデル上の意味、手法の限界まで説明できる必要がある。\par
Look-ahead は、現在の決定を評価するために将来を仮想的に見る方法である。単一シナリオだけを見ると、そのシナリオに過剰適合する危険がある。そこで Multiple Scenario Approach では、複数の将来シナリオをサンプリングし、それぞれを確定的な問題として解く。\par
MSAの基本手順は、複数の将来情報列を生成し、各シナリオでMIPなどの確定的最適化を解き、実際にはまだ見えていない将来情報に依存する部分を取り除き、現在時点で実行可能な部分を比較する、という流れである。\par
Consensus function は、複数シナリオで得られた解の中から代表的な解を選ぶための類似度尺度である。生産スケジューリングなら次に処理するジョブ、VRPなら同じ顧客ペアが同じルートにあるか、TSPならアークや隣接関係が一致するか、施設配置なら選ばれた施設集合が一致するかで測る。\par
MSAの利点は、通常のMIPソルバーを使えることと、複数シナリオに対して極端に悪くない現在決定を選べることである。欠点は、stochasticityを解選択で間接的に扱うだけで、動的な再意思決定そのものを完全にモデル化しているわけではない点である。\par
試験では、MSAを『複数シナリオの平均を取る』だけと書かない。各シナリオの解を作り、その解の共通性や代表性を測り、現在実装すべき決定を選ぶ、という構造を説明する。Consensus functionは問題ごとに設計する必要がある。\par\NoteSection{確認問題と解答}\begin{enumerate}\item \textbf{L07-S008: 「MSA: Consensus Function」の概念を、定義・具体例・モデル上の意味の順に説明せよ。}\\定義としては、MSAは複数シナリオを解き、解の共通性から現在決定を選ぶ。。具体例では、配送・在庫・フードデリバリーのような現実問題を用い、状態、決定、外生情報、報酬、または情報モデル/意思決定モデルのどこに対応するかを明示する。モデル上の意味として、この概念が目的関数、制約、状態表現、将来価値、または方策選択にどのような影響を与えるかを書く。試験では、用語だけを列挙せず、入力データから意思決定、さらに現実の実装へつながる流れを文章で示すことが重要である。\item \textbf{L07-S008: このスライドの考え方を、講義全体のDDDMプロセスの中に位置づけよ。}\\DDDMプロセスでは、現実世界のデータを集約して情報モデルを作り、意思決定モデルまたは方策で行動を選び、実装後に結果を観測して次の状態へ進む。このスライドの概念は、そのどこか一部だけではなく、データ表現、意思決定、将来不確実性、計算可能性のいずれかに関わる。答案では、まずこの概念がどの段階に属するかを述べ、次に他の段階へどう影響するかを書く。例えば価値関数なら将来価値を現在決定へ戻す役割、FIFOなら旅行時間情報モデルの整合性を保証する役割、CFAなら目的関数を調整して将来に良い行動を誘導する役割である。\end{enumerate}
\end{minipage}
\end{adjustbox}
\end{minipage}


\clearpage
\SlideHeader{Lecture 7 - Look-ahead Methods}{009}{C Bierwirth, DC Mattfeld, H Kopfer: On permutation}
\begin{minipage}[t]{0.405\textwidth}
\SlideImage{9}{../Materials/2026/3DM_07_Lookahead-Methods_SS26.pdf}
\begin{adjustbox}{max totalsize={\textwidth}{0.43\textheight},center}
\begin{minipage}{\textwidth}\scriptsize
\NoteSection{日本語訳}\begin{itemize}
\item C Bierwirth, DC Mattfeld, H Kopfer: On permutation（この用語は講義スライド上の専門語として扱う）
\item Excursion: consensus for schedules representations for scheduling problems, International（この用語は講義スライド上の専門語として扱う）
\item conference on parallel problem solving from Nature, 1996（この用語は講義スライド上の専門語として扱う）
\item In a job shop, three jobs (1,2,3), (4,5,6) and (7,8,9) are to processed by three machines（この用語は講義スライド上の専門語として扱う）
\item The three machines are depicted in light grey, grey and dark grey.（この用語は講義スライド上の専門語として扱う）
\item We aim at measuring similarity among schedules. Therefore we consider precedence.（この用語は講義スライド上の専門語として扱う）
\item This schedule is represented（この用語は講義スライド上の専門語として扱う）
\item as binary vector, where oper-（この用語は講義スライド上の専門語として扱う）
\item ations to be processed on the（この用語は講義スライド上の専門語として扱う）
\end{itemize}\NoteSection{試験対策ポイント}\begin{itemize}
\item MSAは複数シナリオを解き、解の共通性から現在決定を選ぶ。
\item Consensus functionは問題構造に依存して設計する。
\end{itemize}
\end{minipage}
\end{adjustbox}
\end{minipage}\hfill
\begin{minipage}[t]{0.58\textwidth}
\begin{adjustbox}{max totalsize={\textwidth}{0.89\textheight},center}
\begin{minipage}{\textwidth}\scriptsize
\NoteSection{解説}このページでは「C Bierwirth, DC Mattfeld, H Kopfer: On permutation」を理解する。スライド上の表現は短いが、試験では定義、具体例、モデル上の意味、手法の限界まで説明できる必要がある。\par
Look-ahead は、現在の決定を評価するために将来を仮想的に見る方法である。単一シナリオだけを見ると、そのシナリオに過剰適合する危険がある。そこで Multiple Scenario Approach では、複数の将来シナリオをサンプリングし、それぞれを確定的な問題として解く。\par
MSAの基本手順は、複数の将来情報列を生成し、各シナリオでMIPなどの確定的最適化を解き、実際にはまだ見えていない将来情報に依存する部分を取り除き、現在時点で実行可能な部分を比較する、という流れである。\par
Consensus function は、複数シナリオで得られた解の中から代表的な解を選ぶための類似度尺度である。生産スケジューリングなら次に処理するジョブ、VRPなら同じ顧客ペアが同じルートにあるか、TSPならアークや隣接関係が一致するか、施設配置なら選ばれた施設集合が一致するかで測る。\par
MSAの利点は、通常のMIPソルバーを使えることと、複数シナリオに対して極端に悪くない現在決定を選べることである。欠点は、stochasticityを解選択で間接的に扱うだけで、動的な再意思決定そのものを完全にモデル化しているわけではない点である。\par
試験では、MSAを『複数シナリオの平均を取る』だけと書かない。各シナリオの解を作り、その解の共通性や代表性を測り、現在実装すべき決定を選ぶ、という構造を説明する。Consensus functionは問題ごとに設計する必要がある。\par\NoteSection{確認問題と解答}\begin{enumerate}\item \textbf{L07-S009: 「C Bierwirth, DC Mattfeld, H Kopfer: On permutation」の概念を、定義・具体例・モデル上の意味の順に説明せよ。}\\定義としては、MSAは複数シナリオを解き、解の共通性から現在決定を選ぶ。。具体例では、配送・在庫・フードデリバリーのような現実問題を用い、状態、決定、外生情報、報酬、または情報モデル/意思決定モデルのどこに対応するかを明示する。モデル上の意味として、この概念が目的関数、制約、状態表現、将来価値、または方策選択にどのような影響を与えるかを書く。試験では、用語だけを列挙せず、入力データから意思決定、さらに現実の実装へつながる流れを文章で示すことが重要である。\item \textbf{L07-S009: このスライドの考え方を、講義全体のDDDMプロセスの中に位置づけよ。}\\DDDMプロセスでは、現実世界のデータを集約して情報モデルを作り、意思決定モデルまたは方策で行動を選び、実装後に結果を観測して次の状態へ進む。このスライドの概念は、そのどこか一部だけではなく、データ表現、意思決定、将来不確実性、計算可能性のいずれかに関わる。答案では、まずこの概念がどの段階に属するかを述べ、次に他の段階へどう影響するかを書く。例えば価値関数なら将来価値を現在決定へ戻す役割、FIFOなら旅行時間情報モデルの整合性を保証する役割、CFAなら目的関数を調整して将来に良い行動を誘導する役割である。\end{enumerate}
\end{minipage}
\end{adjustbox}
\end{minipage}


\clearpage
\SlideHeader{Lecture 7 - Look-ahead Methods}{010}{Excursion: consensus for facility location / travelling salesman}
\begin{minipage}[t]{0.405\textwidth}
\SlideImage{10}{../Materials/2026/3DM_07_Lookahead-Methods_SS26.pdf}
\begin{adjustbox}{max totalsize={\textwidth}{0.43\textheight},center}
\begin{minipage}{\textwidth}\scriptsize
\NoteSection{日本語訳}\begin{itemize}
\item Excursion: consensus for facility location / travelling salesman（この用語は講義スライド上の専門語として扱う）
\item The facility location problem picks a subset of locations from 1（この用語は講義スライド上の専門語として扱う）
\item a potentially larger set of possible locatitons（この用語は講義スライド上の専門語として扱う）
\item 1 2 3 4 5 6 7 8
\item Representation of solution: 6 8（この用語は講義スライド上の専門語として扱う）
\item 0 1 1 0 1 0 1 0
\item 5 4
\item The symmetric travelling salesman problem constructs a tour（この用語は講義スライド上の専門語として扱う）
\item among all locations by selecting arcs that connect nodes（この用語は講義スライド上の専門語として扱う）
\end{itemize}\NoteSection{試験対策ポイント}\begin{itemize}
\item MSAは複数シナリオを解き、解の共通性から現在決定を選ぶ。
\item Consensus functionは問題構造に依存して設計する。
\end{itemize}
\end{minipage}
\end{adjustbox}
\end{minipage}\hfill
\begin{minipage}[t]{0.58\textwidth}
\begin{adjustbox}{max totalsize={\textwidth}{0.89\textheight},center}
\begin{minipage}{\textwidth}\scriptsize
\NoteSection{解説}このページでは「Excursion: consensus for facility location / travelling salesman」を理解する。スライド上の表現は短いが、試験では定義、具体例、モデル上の意味、手法の限界まで説明できる必要がある。\par
Look-ahead は、現在の決定を評価するために将来を仮想的に見る方法である。単一シナリオだけを見ると、そのシナリオに過剰適合する危険がある。そこで Multiple Scenario Approach では、複数の将来シナリオをサンプリングし、それぞれを確定的な問題として解く。\par
MSAの基本手順は、複数の将来情報列を生成し、各シナリオでMIPなどの確定的最適化を解き、実際にはまだ見えていない将来情報に依存する部分を取り除き、現在時点で実行可能な部分を比較する、という流れである。\par
Consensus function は、複数シナリオで得られた解の中から代表的な解を選ぶための類似度尺度である。生産スケジューリングなら次に処理するジョブ、VRPなら同じ顧客ペアが同じルートにあるか、TSPならアークや隣接関係が一致するか、施設配置なら選ばれた施設集合が一致するかで測る。\par
MSAの利点は、通常のMIPソルバーを使えることと、複数シナリオに対して極端に悪くない現在決定を選べることである。欠点は、stochasticityを解選択で間接的に扱うだけで、動的な再意思決定そのものを完全にモデル化しているわけではない点である。\par
試験では、MSAを『複数シナリオの平均を取る』だけと書かない。各シナリオの解を作り、その解の共通性や代表性を測り、現在実装すべき決定を選ぶ、という構造を説明する。Consensus functionは問題ごとに設計する必要がある。\par\NoteSection{確認問題と解答}\begin{enumerate}\item \textbf{L07-S010: 「Excursion: consensus for facility location / travelling salesman」の概念を、定義・具体例・モデル上の意味の順に説明せよ。}\\定義としては、MSAは複数シナリオを解き、解の共通性から現在決定を選ぶ。。具体例では、配送・在庫・フードデリバリーのような現実問題を用い、状態、決定、外生情報、報酬、または情報モデル/意思決定モデルのどこに対応するかを明示する。モデル上の意味として、この概念が目的関数、制約、状態表現、将来価値、または方策選択にどのような影響を与えるかを書く。試験では、用語だけを列挙せず、入力データから意思決定、さらに現実の実装へつながる流れを文章で示すことが重要である。\item \textbf{L07-S010: このスライドの考え方を、講義全体のDDDMプロセスの中に位置づけよ。}\\DDDMプロセスでは、現実世界のデータを集約して情報モデルを作り、意思決定モデルまたは方策で行動を選び、実装後に結果を観測して次の状態へ進む。このスライドの概念は、そのどこか一部だけではなく、データ表現、意思決定、将来不確実性、計算可能性のいずれかに関わる。答案では、まずこの概念がどの段階に属するかを述べ、次に他の段階へどう影響するかを書く。例えば価値関数なら将来価値を現在決定へ戻す役割、FIFOなら旅行時間情報モデルの整合性を保証する役割、CFAなら目的関数を調整して将来に良い行動を誘導する役割である。\end{enumerate}
\end{minipage}
\end{adjustbox}
\end{minipage}


\clearpage
\SlideHeader{Lecture 7 - Look-ahead Methods}{011}{Case study: Autonomy of Crowdsourced Suppliers}
\begin{minipage}[t]{0.405\textwidth}
\SlideImage{11}{../Materials/2026/3DM_07_Lookahead-Methods_SS26.pdf}
\begin{adjustbox}{max totalsize={\textwidth}{0.43\textheight},center}
\begin{minipage}{\textwidth}\scriptsize
\NoteSection{日本語訳}\begin{itemize}
\item ケーススタディ: Autonomy of Crowdsourced Suppliers
\item In a gig-economy, suppliers work for multiple clients（この用語は講義スライド上の専門語として扱う）
\item Suppliers therefore can reject offered assignments（この用語は講義スライド上の専門語として扱う）
\item Reasons: travel duration, location, payment, etc.（この用語は講義スライド上の専門語として扱う）
\item Problems arise for the requesting client（この用語は講義スライド上の専門語として扱う）
\item Less money for the platform（この用語は講義スライド上の専門語として扱う）
\item Delays, customer cancelations（この用語は講義スライド上の専門語として扱う）
\item Frustration among drivers who arrive late（この用語は講義スライド上の専門語として扱う）
\item Hedge against rejections with so called menus（この用語は講義スライド上の専門語として扱う）
\end{itemize}\NoteSection{試験対策ポイント}\begin{itemize}
\item Policy Function Approximation, Rolling Horizon, CFA, Look-ahead, VFAを何を近似するかで区別する。
\item 各手法の強みと弱みを1セットで説明する。
\end{itemize}
\end{minipage}
\end{adjustbox}
\end{minipage}\hfill
\begin{minipage}[t]{0.58\textwidth}
\begin{adjustbox}{max totalsize={\textwidth}{0.89\textheight},center}
\begin{minipage}{\textwidth}\scriptsize
\NoteSection{解説}このページでは「Case study: Autonomy of Crowdsourced Suppliers」を理解する。スライド上の表現は短いが、試験では定義、具体例、モデル上の意味、手法の限界まで説明できる必要がある。\par
Policy Function Approximation は、経験則やルールを直接方策として使う方法である。例えば『在庫が閾値を下回ったら注文する』『最も近いドライバーへ割り当てる』などである。計算は速く解釈しやすいが、複雑な将来影響を十分に扱えない可能性がある。\par
Rolling Horizon Re-Optimization は、現在時点で見えている情報を使って一定期間の最適化問題を解き、最初の一部だけ実行し、次時点でまた解き直す方法である。現実変化へ適応できるが、見えている範囲だけを最適化するため近視眼的になることがある。\par
Cost Function Approximation は、現在の目的関数や制約に補正項を加えて、将来に良い影響を持つ決定を誘導する方法である。例えば、将来需要が多い地域から車両を離さないようにペナルティを入れる。設計は問題依存で、補正項が適切でないと逆効果になる。\par
Look-ahead Methods は、将来シナリオを明示的に考え、現在決定の良し悪しを未来までシミュレーションまたは最適化して評価する方法である。将来を直接見るため直感的だが、計算負荷が大きく、サンプリングやホライズンの選び方に依存する。\par
Value Function Approximation は、状態または後決定状態の将来価値を特徴量とモデルで近似する方法である。学習した価値を使えば、各決定について遠い未来まで毎回シミュレーションしなくても将来影響を考慮できる。しかし良い特徴量と学習データが必要である。\par\NoteSection{確認問題と解答}\begin{enumerate}\item \textbf{L07-S011: 「Case study: Autonomy of Crowdsourced Suppliers」の概念を、定義・具体例・モデル上の意味の順に説明せよ。}\\定義としては、Policy Function Approximation, Rolling Horizon, CFA, Look-ahead, VFAを何を近似するかで区別する。。具体例では、配送・在庫・フードデリバリーのような現実問題を用い、状態、決定、外生情報、報酬、または情報モデル/意思決定モデルのどこに対応するかを明示する。モデル上の意味として、この概念が目的関数、制約、状態表現、将来価値、または方策選択にどのような影響を与えるかを書く。試験では、用語だけを列挙せず、入力データから意思決定、さらに現実の実装へつながる流れを文章で示すことが重要である。\item \textbf{L07-S011: このスライドの考え方を、講義全体のDDDMプロセスの中に位置づけよ。}\\DDDMプロセスでは、現実世界のデータを集約して情報モデルを作り、意思決定モデルまたは方策で行動を選び、実装後に結果を観測して次の状態へ進む。このスライドの概念は、そのどこか一部だけではなく、データ表現、意思決定、将来不確実性、計算可能性のいずれかに関わる。答案では、まずこの概念がどの段階に属するかを述べ、次に他の段階へどう影響するかを書く。例えば価値関数なら将来価値を現在決定へ戻す役割、FIFOなら旅行時間情報モデルの整合性を保証する役割、CFAなら目的関数を調整して将来に良い行動を誘導する役割である。\end{enumerate}
\end{minipage}
\end{adjustbox}
\end{minipage}


\clearpage
\SlideHeader{Lecture 7 - Look-ahead Methods}{012}{Menus}
\begin{minipage}[t]{0.405\textwidth}
\SlideImage{12}{../Materials/2026/3DM_07_Lookahead-Methods_SS26.pdf}
\begin{adjustbox}{max totalsize={\textwidth}{0.43\textheight},center}
\begin{minipage}{\textwidth}\scriptsize
\NoteSection{日本語訳}\begin{itemize}
\item Menus（この用語は講義スライド上の専門語として扱う）
\item For every supplier a menu of potential customers is offered（この用語は講義スライド上の専門語として扱う）
\item Instead of asking suppliers about favorite deliveries,（この用語は講義スライド上の専門語として扱う）
\item a choice model depicts possible preferences 1（この用語は講義スライド上の専門語として扱う）
\item Every supplier selects her/his “best” choice（この用語は講義スライド上の専門語として扱う）
\item from the menu of suggested deliveries A C D（この用語は講義スライド上の専門語として扱う）
\item B
\item Platform finally decides on assignment（この用語は講義スライド上の専門語として扱う）
\item Multiple offers increase the probability of selection 3（この用語は講義スライド上の専門語として扱う）
\end{itemize}\NoteSection{試験対策ポイント}\begin{itemize}
\item Policy Function Approximation, Rolling Horizon, CFA, Look-ahead, VFAを何を近似するかで区別する。
\item 各手法の強みと弱みを1セットで説明する。
\end{itemize}
\end{minipage}
\end{adjustbox}
\end{minipage}\hfill
\begin{minipage}[t]{0.58\textwidth}
\begin{adjustbox}{max totalsize={\textwidth}{0.89\textheight},center}
\begin{minipage}{\textwidth}\scriptsize
\NoteSection{解説}このページでは「Menus」を理解する。スライド上の表現は短いが、試験では定義、具体例、モデル上の意味、手法の限界まで説明できる必要がある。\par
Policy Function Approximation は、経験則やルールを直接方策として使う方法である。例えば『在庫が閾値を下回ったら注文する』『最も近いドライバーへ割り当てる』などである。計算は速く解釈しやすいが、複雑な将来影響を十分に扱えない可能性がある。\par
Rolling Horizon Re-Optimization は、現在時点で見えている情報を使って一定期間の最適化問題を解き、最初の一部だけ実行し、次時点でまた解き直す方法である。現実変化へ適応できるが、見えている範囲だけを最適化するため近視眼的になることがある。\par
Cost Function Approximation は、現在の目的関数や制約に補正項を加えて、将来に良い影響を持つ決定を誘導する方法である。例えば、将来需要が多い地域から車両を離さないようにペナルティを入れる。設計は問題依存で、補正項が適切でないと逆効果になる。\par
Look-ahead Methods は、将来シナリオを明示的に考え、現在決定の良し悪しを未来までシミュレーションまたは最適化して評価する方法である。将来を直接見るため直感的だが、計算負荷が大きく、サンプリングやホライズンの選び方に依存する。\par
Value Function Approximation は、状態または後決定状態の将来価値を特徴量とモデルで近似する方法である。学習した価値を使えば、各決定について遠い未来まで毎回シミュレーションしなくても将来影響を考慮できる。しかし良い特徴量と学習データが必要である。\par\NoteSection{確認問題と解答}\begin{enumerate}\item \textbf{L07-S012: 「Menus」の概念を、定義・具体例・モデル上の意味の順に説明せよ。}\\定義としては、Policy Function Approximation, Rolling Horizon, CFA, Look-ahead, VFAを何を近似するかで区別する。。具体例では、配送・在庫・フードデリバリーのような現実問題を用い、状態、決定、外生情報、報酬、または情報モデル/意思決定モデルのどこに対応するかを明示する。モデル上の意味として、この概念が目的関数、制約、状態表現、将来価値、または方策選択にどのような影響を与えるかを書く。試験では、用語だけを列挙せず、入力データから意思決定、さらに現実の実装へつながる流れを文章で示すことが重要である。\item \textbf{L07-S012: このスライドの考え方を、講義全体のDDDMプロセスの中に位置づけよ。}\\DDDMプロセスでは、現実世界のデータを集約して情報モデルを作り、意思決定モデルまたは方策で行動を選び、実装後に結果を観測して次の状態へ進む。このスライドの概念は、そのどこか一部だけではなく、データ表現、意思決定、将来不確実性、計算可能性のいずれかに関わる。答案では、まずこの概念がどの段階に属するかを述べ、次に他の段階へどう影響するかを書く。例えば価値関数なら将来価値を現在決定へ戻す役割、FIFOなら旅行時間情報モデルの整合性を保証する役割、CFAなら目的関数を調整して将来に良い行動を誘導する役割である。\end{enumerate}
\end{minipage}
\end{adjustbox}
\end{minipage}


\clearpage
\SlideHeader{Lecture 7 - Look-ahead Methods}{013}{Decision sequence}
\begin{minipage}[t]{0.405\textwidth}
\SlideImage{13}{../Materials/2026/3DM_07_Lookahead-Methods_SS26.pdf}
\begin{adjustbox}{max totalsize={\textwidth}{0.43\textheight},center}
\begin{minipage}{\textwidth}\scriptsize
\NoteSection{日本語訳}\begin{itemize}
\item 決定 sequence
\item At time t=5 two suppliers and three request are observed（この用語は講義スライド上の専門語として扱う）
\item Request A and B are offered to supplier 1, B and C are offered to supplier 2（この用語は講義スライド上の専門語として扱う）
\item Both suppliers opt for request B, and eventually B is assigned to supplier 1 → implement（この用語は講義スライド上の専門語として扱う）
\item At time t=6 A has disappeared, C remains, D is new and supplier 3 has come into play（この用語は講義スライド上の専門語として扱う）
\end{itemize}\NoteSection{試験対策ポイント}\begin{itemize}
\item Policy Function Approximation, Rolling Horizon, CFA, Look-ahead, VFAを何を近似するかで区別する。
\item 各手法の強みと弱みを1セットで説明する。
\end{itemize}
\end{minipage}
\end{adjustbox}
\end{minipage}\hfill
\begin{minipage}[t]{0.58\textwidth}
\begin{adjustbox}{max totalsize={\textwidth}{0.89\textheight},center}
\begin{minipage}{\textwidth}\scriptsize
\NoteSection{解説}このページでは「Decision sequence」を理解する。スライド上の表現は短いが、試験では定義、具体例、モデル上の意味、手法の限界まで説明できる必要がある。\par
Policy Function Approximation は、経験則やルールを直接方策として使う方法である。例えば『在庫が閾値を下回ったら注文する』『最も近いドライバーへ割り当てる』などである。計算は速く解釈しやすいが、複雑な将来影響を十分に扱えない可能性がある。\par
Rolling Horizon Re-Optimization は、現在時点で見えている情報を使って一定期間の最適化問題を解き、最初の一部だけ実行し、次時点でまた解き直す方法である。現実変化へ適応できるが、見えている範囲だけを最適化するため近視眼的になることがある。\par
Cost Function Approximation は、現在の目的関数や制約に補正項を加えて、将来に良い影響を持つ決定を誘導する方法である。例えば、将来需要が多い地域から車両を離さないようにペナルティを入れる。設計は問題依存で、補正項が適切でないと逆効果になる。\par
Look-ahead Methods は、将来シナリオを明示的に考え、現在決定の良し悪しを未来までシミュレーションまたは最適化して評価する方法である。将来を直接見るため直感的だが、計算負荷が大きく、サンプリングやホライズンの選び方に依存する。\par
Value Function Approximation は、状態または後決定状態の将来価値を特徴量とモデルで近似する方法である。学習した価値を使えば、各決定について遠い未来まで毎回シミュレーションしなくても将来影響を考慮できる。しかし良い特徴量と学習データが必要である。\par\NoteSection{確認問題と解答}\begin{enumerate}\item \textbf{L07-S013: 「Decision sequence」の概念を、定義・具体例・モデル上の意味の順に説明せよ。}\\定義としては、Policy Function Approximation, Rolling Horizon, CFA, Look-ahead, VFAを何を近似するかで区別する。。具体例では、配送・在庫・フードデリバリーのような現実問題を用い、状態、決定、外生情報、報酬、または情報モデル/意思決定モデルのどこに対応するかを明示する。モデル上の意味として、この概念が目的関数、制約、状態表現、将来価値、または方策選択にどのような影響を与えるかを書く。試験では、用語だけを列挙せず、入力データから意思決定、さらに現実の実装へつながる流れを文章で示すことが重要である。\item \textbf{L07-S013: このスライドの考え方を、講義全体のDDDMプロセスの中に位置づけよ。}\\DDDMプロセスでは、現実世界のデータを集約して情報モデルを作り、意思決定モデルまたは方策で行動を選び、実装後に結果を観測して次の状態へ進む。このスライドの概念は、そのどこか一部だけではなく、データ表現、意思決定、将来不確実性、計算可能性のいずれかに関わる。答案では、まずこの概念がどの段階に属するかを述べ、次に他の段階へどう影響するかを書く。例えば価値関数なら将来価値を現在決定へ戻す役割、FIFOなら旅行時間情報モデルの整合性を保証する役割、CFAなら目的関数を調整して将来に良い行動を誘導する役割である。\end{enumerate}
\end{minipage}
\end{adjustbox}
\end{minipage}


\clearpage
\SlideHeader{Lecture 7 - Look-ahead Methods}{014}{Multiple Scenario Approach}
\begin{minipage}[t]{0.405\textwidth}
\SlideImage{14}{../Materials/2026/3DM_07_Lookahead-Methods_SS26.pdf}
\begin{adjustbox}{max totalsize={\textwidth}{0.43\textheight},center}
\begin{minipage}{\textwidth}\scriptsize
\NoteSection{日本語訳}\begin{itemize}
\item 複数シナリオアプローチ
\item In each 状態
\item Generate a set of scenarios（この用語は講義スライド上の専門語として扱う）
\item Scenario is a known ranking of customers（この用語は講義スライド上の専門語として扱う）
\item for each supplier（この用語は講義スライド上の専門語として扱う）
\item Solve individual scenarios with some IP-magic（この用語は講義スライド上の専門語として扱う）
\item Derive “consensus” solution from scenario solutions（この用語は講義スライド上の専門語として扱う）
\item Observe customer selection of supplier（この用語は講義スライド上の専門語として扱う）
\item Perform final assignment of customers to supplier（この用語は講義スライド上の専門語として扱う）
\end{itemize}\NoteSection{試験対策ポイント}\begin{itemize}
\item MSAは複数シナリオを解き、解の共通性から現在決定を選ぶ。
\item Consensus functionは問題構造に依存して設計する。
\end{itemize}
\end{minipage}
\end{adjustbox}
\end{minipage}\hfill
\begin{minipage}[t]{0.58\textwidth}
\begin{adjustbox}{max totalsize={\textwidth}{0.89\textheight},center}
\begin{minipage}{\textwidth}\scriptsize
\NoteSection{解説}このページでは「Multiple Scenario Approach」を理解する。スライド上の表現は短いが、試験では定義、具体例、モデル上の意味、手法の限界まで説明できる必要がある。\par
Look-ahead は、現在の決定を評価するために将来を仮想的に見る方法である。単一シナリオだけを見ると、そのシナリオに過剰適合する危険がある。そこで Multiple Scenario Approach では、複数の将来シナリオをサンプリングし、それぞれを確定的な問題として解く。\par
MSAの基本手順は、複数の将来情報列を生成し、各シナリオでMIPなどの確定的最適化を解き、実際にはまだ見えていない将来情報に依存する部分を取り除き、現在時点で実行可能な部分を比較する、という流れである。\par
Consensus function は、複数シナリオで得られた解の中から代表的な解を選ぶための類似度尺度である。生産スケジューリングなら次に処理するジョブ、VRPなら同じ顧客ペアが同じルートにあるか、TSPならアークや隣接関係が一致するか、施設配置なら選ばれた施設集合が一致するかで測る。\par
MSAの利点は、通常のMIPソルバーを使えることと、複数シナリオに対して極端に悪くない現在決定を選べることである。欠点は、stochasticityを解選択で間接的に扱うだけで、動的な再意思決定そのものを完全にモデル化しているわけではない点である。\par
試験では、MSAを『複数シナリオの平均を取る』だけと書かない。各シナリオの解を作り、その解の共通性や代表性を測り、現在実装すべき決定を選ぶ、という構造を説明する。Consensus functionは問題ごとに設計する必要がある。\par\NoteSection{確認問題と解答}\begin{enumerate}\item \textbf{L07-S014: 「Multiple Scenario Approach」の概念を、定義・具体例・モデル上の意味の順に説明せよ。}\\定義としては、MSAは複数シナリオを解き、解の共通性から現在決定を選ぶ。。具体例では、配送・在庫・フードデリバリーのような現実問題を用い、状態、決定、外生情報、報酬、または情報モデル/意思決定モデルのどこに対応するかを明示する。モデル上の意味として、この概念が目的関数、制約、状態表現、将来価値、または方策選択にどのような影響を与えるかを書く。試験では、用語だけを列挙せず、入力データから意思決定、さらに現実の実装へつながる流れを文章で示すことが重要である。\item \textbf{L07-S014: このスライドの考え方を、講義全体のDDDMプロセスの中に位置づけよ。}\\DDDMプロセスでは、現実世界のデータを集約して情報モデルを作り、意思決定モデルまたは方策で行動を選び、実装後に結果を観測して次の状態へ進む。このスライドの概念は、そのどこか一部だけではなく、データ表現、意思決定、将来不確実性、計算可能性のいずれかに関わる。答案では、まずこの概念がどの段階に属するかを述べ、次に他の段階へどう影響するかを書く。例えば価値関数なら将来価値を現在決定へ戻す役割、FIFOなら旅行時間情報モデルの整合性を保証する役割、CFAなら目的関数を調整して将来に良い行動を誘導する役割である。\end{enumerate}
\end{minipage}
\end{adjustbox}
\end{minipage}


\clearpage
\SlideHeader{Lecture 7 - Look-ahead Methods}{015}{Step 1: Generating scenarios}
\begin{minipage}[t]{0.405\textwidth}
\SlideImage{15}{../Materials/2026/3DM_07_Lookahead-Methods_SS26.pdf}
\begin{adjustbox}{max totalsize={\textwidth}{0.43\textheight},center}
\begin{minipage}{\textwidth}\scriptsize
\NoteSection{日本語訳}\begin{itemize}
\item Step 1: Generating scenarios（この用語は講義スライド上の専門語として扱う）
\item Preferences of suppliers with regard to attributes of the current 状態 are not known
\item Distance to the customer and 旅行時間 are identified as meaningful attributes
\item The weights of attributes are modified for each supplier at random（この用語は講義スライド上の専門語として扱う）
\item Solving the respective MIP leads to a ranking/rejection of requests for suppliers（この用語は講義スライド上の専門語として扱う）
\end{itemize}\NoteSection{試験対策ポイント}\begin{itemize}
\item Policy Function Approximation, Rolling Horizon, CFA, Look-ahead, VFAを何を近似するかで区別する。
\item 各手法の強みと弱みを1セットで説明する。
\end{itemize}
\end{minipage}
\end{adjustbox}
\end{minipage}\hfill
\begin{minipage}[t]{0.58\textwidth}
\begin{adjustbox}{max totalsize={\textwidth}{0.89\textheight},center}
\begin{minipage}{\textwidth}\scriptsize
\NoteSection{解説}このページでは「Step 1: Generating scenarios」を理解する。スライド上の表現は短いが、試験では定義、具体例、モデル上の意味、手法の限界まで説明できる必要がある。\par
Policy Function Approximation は、経験則やルールを直接方策として使う方法である。例えば『在庫が閾値を下回ったら注文する』『最も近いドライバーへ割り当てる』などである。計算は速く解釈しやすいが、複雑な将来影響を十分に扱えない可能性がある。\par
Rolling Horizon Re-Optimization は、現在時点で見えている情報を使って一定期間の最適化問題を解き、最初の一部だけ実行し、次時点でまた解き直す方法である。現実変化へ適応できるが、見えている範囲だけを最適化するため近視眼的になることがある。\par
Cost Function Approximation は、現在の目的関数や制約に補正項を加えて、将来に良い影響を持つ決定を誘導する方法である。例えば、将来需要が多い地域から車両を離さないようにペナルティを入れる。設計は問題依存で、補正項が適切でないと逆効果になる。\par
Look-ahead Methods は、将来シナリオを明示的に考え、現在決定の良し悪しを未来までシミュレーションまたは最適化して評価する方法である。将来を直接見るため直感的だが、計算負荷が大きく、サンプリングやホライズンの選び方に依存する。\par
Value Function Approximation は、状態または後決定状態の将来価値を特徴量とモデルで近似する方法である。学習した価値を使えば、各決定について遠い未来まで毎回シミュレーションしなくても将来影響を考慮できる。しかし良い特徴量と学習データが必要である。\par\NoteSection{確認問題と解答}\begin{enumerate}\item \textbf{L07-S015: 「Step 1: Generating scenarios」の概念を、定義・具体例・モデル上の意味の順に説明せよ。}\\定義としては、Policy Function Approximation, Rolling Horizon, CFA, Look-ahead, VFAを何を近似するかで区別する。。具体例では、配送・在庫・フードデリバリーのような現実問題を用い、状態、決定、外生情報、報酬、または情報モデル/意思決定モデルのどこに対応するかを明示する。モデル上の意味として、この概念が目的関数、制約、状態表現、将来価値、または方策選択にどのような影響を与えるかを書く。試験では、用語だけを列挙せず、入力データから意思決定、さらに現実の実装へつながる流れを文章で示すことが重要である。\item \textbf{L07-S015: このスライドの考え方を、講義全体のDDDMプロセスの中に位置づけよ。}\\DDDMプロセスでは、現実世界のデータを集約して情報モデルを作り、意思決定モデルまたは方策で行動を選び、実装後に結果を観測して次の状態へ進む。このスライドの概念は、そのどこか一部だけではなく、データ表現、意思決定、将来不確実性、計算可能性のいずれかに関わる。答案では、まずこの概念がどの段階に属するかを述べ、次に他の段階へどう影響するかを書く。例えば価値関数なら将来価値を現在決定へ戻す役割、FIFOなら旅行時間情報モデルの整合性を保証する役割、CFAなら目的関数を調整して将来に良い行動を誘導する役割である。\end{enumerate}
\end{minipage}
\end{adjustbox}
\end{minipage}


\clearpage
\SlideHeader{Lecture 7 - Look-ahead Methods}{016}{Step 2: Solving an bi-level Integer Program}
\begin{minipage}[t]{0.405\textwidth}
\SlideImage{16}{../Materials/2026/3DM_07_Lookahead-Methods_SS26.pdf}
\begin{adjustbox}{max totalsize={\textwidth}{0.43\textheight},center}
\begin{minipage}{\textwidth}\scriptsize
\NoteSection{日本語訳}\begin{itemize}
\item Step 2: Solving an bi-level Integer Program（この用語は講義スライド上の専門語として扱う）
\item First level: Assignment of customer requests to supplier in a scenario（この用語は講義スライド上の専門語として扱う）
\item Second level: Supplier (artificially) select requests in accordance to rank（この用語は講義スライド上の専門語として扱う）
\item Conversion of second level objective function in constraint → single-level IP（この用語は講義スライド上の専門語として扱う）
\item Solve the resulting assignment problem for 𝑥𝑖𝑗𝑡𝜎 supplier 𝑖, customer 𝑗, 状態 𝑡, scenario 𝜎
\item The assignment problem is variant of the transportation problem → modi method（この用語は講義スライド上の専門語として扱う）
\end{itemize}\NoteSection{試験対策ポイント}\begin{itemize}
\item Policy Function Approximation, Rolling Horizon, CFA, Look-ahead, VFAを何を近似するかで区別する。
\item 各手法の強みと弱みを1セットで説明する。
\end{itemize}
\end{minipage}
\end{adjustbox}
\end{minipage}\hfill
\begin{minipage}[t]{0.58\textwidth}
\begin{adjustbox}{max totalsize={\textwidth}{0.89\textheight},center}
\begin{minipage}{\textwidth}\scriptsize
\NoteSection{解説}このページでは「Step 2: Solving an bi-level Integer Program」を理解する。スライド上の表現は短いが、試験では定義、具体例、モデル上の意味、手法の限界まで説明できる必要がある。\par
Policy Function Approximation は、経験則やルールを直接方策として使う方法である。例えば『在庫が閾値を下回ったら注文する』『最も近いドライバーへ割り当てる』などである。計算は速く解釈しやすいが、複雑な将来影響を十分に扱えない可能性がある。\par
Rolling Horizon Re-Optimization は、現在時点で見えている情報を使って一定期間の最適化問題を解き、最初の一部だけ実行し、次時点でまた解き直す方法である。現実変化へ適応できるが、見えている範囲だけを最適化するため近視眼的になることがある。\par
Cost Function Approximation は、現在の目的関数や制約に補正項を加えて、将来に良い影響を持つ決定を誘導する方法である。例えば、将来需要が多い地域から車両を離さないようにペナルティを入れる。設計は問題依存で、補正項が適切でないと逆効果になる。\par
Look-ahead Methods は、将来シナリオを明示的に考え、現在決定の良し悪しを未来までシミュレーションまたは最適化して評価する方法である。将来を直接見るため直感的だが、計算負荷が大きく、サンプリングやホライズンの選び方に依存する。\par
Value Function Approximation は、状態または後決定状態の将来価値を特徴量とモデルで近似する方法である。学習した価値を使えば、各決定について遠い未来まで毎回シミュレーションしなくても将来影響を考慮できる。しかし良い特徴量と学習データが必要である。\par\NoteSection{確認問題と解答}\begin{enumerate}\item \textbf{L07-S016: 「Step 2: Solving an bi-level Integer Program」の概念を、定義・具体例・モデル上の意味の順に説明せよ。}\\定義としては、Policy Function Approximation, Rolling Horizon, CFA, Look-ahead, VFAを何を近似するかで区別する。。具体例では、配送・在庫・フードデリバリーのような現実問題を用い、状態、決定、外生情報、報酬、または情報モデル/意思決定モデルのどこに対応するかを明示する。モデル上の意味として、この概念が目的関数、制約、状態表現、将来価値、または方策選択にどのような影響を与えるかを書く。試験では、用語だけを列挙せず、入力データから意思決定、さらに現実の実装へつながる流れを文章で示すことが重要である。\item \textbf{L07-S016: このスライドの考え方を、講義全体のDDDMプロセスの中に位置づけよ。}\\DDDMプロセスでは、現実世界のデータを集約して情報モデルを作り、意思決定モデルまたは方策で行動を選び、実装後に結果を観測して次の状態へ進む。このスライドの概念は、そのどこか一部だけではなく、データ表現、意思決定、将来不確実性、計算可能性のいずれかに関わる。答案では、まずこの概念がどの段階に属するかを述べ、次に他の段階へどう影響するかを書く。例えば価値関数なら将来価値を現在決定へ戻す役割、FIFOなら旅行時間情報モデルの整合性を保証する役割、CFAなら目的関数を調整して将来に良い行動を誘導する役割である。\end{enumerate}
\end{minipage}
\end{adjustbox}
\end{minipage}


\clearpage
\SlideHeader{Lecture 7 - Look-ahead Methods}{017}{Step 3: Consensus solution}
\begin{minipage}[t]{0.405\textwidth}
\SlideImage{17}{../Materials/2026/3DM_07_Lookahead-Methods_SS26.pdf}
\begin{adjustbox}{max totalsize={\textwidth}{0.43\textheight},center}
\begin{minipage}{\textwidth}\scriptsize
\NoteSection{日本語訳}\begin{itemize}
\item Step 3: Consensus solution（この用語は講義スライド上の専門語として扱う）
\item Reflects the structure discovered by the scenario solutions by means of a majority vote（この用語は講義スライド上の専門語として扱う）
\item Compute frequency fijt = σ𝜎 𝑥𝑖𝑗𝑡𝜎 over scenarios 𝜎 and 最大化する σ𝑖 σ𝑗 𝑓𝑖𝑗𝑡 𝑥𝑖𝑗𝑡
\item The consensus 𝑥𝑖𝑗𝑡 considers frequently observed requests in one feasible solution（この用語は講義スライド上の専門語として扱う）
\end{itemize}\NoteSection{試験対策ポイント}\begin{itemize}
\item MSAは複数シナリオを解き、解の共通性から現在決定を選ぶ。
\item Consensus functionは問題構造に依存して設計する。
\end{itemize}
\end{minipage}
\end{adjustbox}
\end{minipage}\hfill
\begin{minipage}[t]{0.58\textwidth}
\begin{adjustbox}{max totalsize={\textwidth}{0.89\textheight},center}
\begin{minipage}{\textwidth}\scriptsize
\NoteSection{解説}このページでは「Step 3: Consensus solution」を理解する。スライド上の表現は短いが、試験では定義、具体例、モデル上の意味、手法の限界まで説明できる必要がある。\par
Look-ahead は、現在の決定を評価するために将来を仮想的に見る方法である。単一シナリオだけを見ると、そのシナリオに過剰適合する危険がある。そこで Multiple Scenario Approach では、複数の将来シナリオをサンプリングし、それぞれを確定的な問題として解く。\par
MSAの基本手順は、複数の将来情報列を生成し、各シナリオでMIPなどの確定的最適化を解き、実際にはまだ見えていない将来情報に依存する部分を取り除き、現在時点で実行可能な部分を比較する、という流れである。\par
Consensus function は、複数シナリオで得られた解の中から代表的な解を選ぶための類似度尺度である。生産スケジューリングなら次に処理するジョブ、VRPなら同じ顧客ペアが同じルートにあるか、TSPならアークや隣接関係が一致するか、施設配置なら選ばれた施設集合が一致するかで測る。\par
MSAの利点は、通常のMIPソルバーを使えることと、複数シナリオに対して極端に悪くない現在決定を選べることである。欠点は、stochasticityを解選択で間接的に扱うだけで、動的な再意思決定そのものを完全にモデル化しているわけではない点である。\par
試験では、MSAを『複数シナリオの平均を取る』だけと書かない。各シナリオの解を作り、その解の共通性や代表性を測り、現在実装すべき決定を選ぶ、という構造を説明する。Consensus functionは問題ごとに設計する必要がある。\par\NoteSection{確認問題と解答}\begin{enumerate}\item \textbf{L07-S017: 「Step 3: Consensus solution」の概念を、定義・具体例・モデル上の意味の順に説明せよ。}\\定義としては、MSAは複数シナリオを解き、解の共通性から現在決定を選ぶ。。具体例では、配送・在庫・フードデリバリーのような現実問題を用い、状態、決定、外生情報、報酬、または情報モデル/意思決定モデルのどこに対応するかを明示する。モデル上の意味として、この概念が目的関数、制約、状態表現、将来価値、または方策選択にどのような影響を与えるかを書く。試験では、用語だけを列挙せず、入力データから意思決定、さらに現実の実装へつながる流れを文章で示すことが重要である。\item \textbf{L07-S017: このスライドの考え方を、講義全体のDDDMプロセスの中に位置づけよ。}\\DDDMプロセスでは、現実世界のデータを集約して情報モデルを作り、意思決定モデルまたは方策で行動を選び、実装後に結果を観測して次の状態へ進む。このスライドの概念は、そのどこか一部だけではなく、データ表現、意思決定、将来不確実性、計算可能性のいずれかに関わる。答案では、まずこの概念がどの段階に属するかを述べ、次に他の段階へどう影響するかを書く。例えば価値関数なら将来価値を現在決定へ戻す役割、FIFOなら旅行時間情報モデルの整合性を保証する役割、CFAなら目的関数を調整して将来に良い行動を誘導する役割である。\end{enumerate}
\end{minipage}
\end{adjustbox}
\end{minipage}


\clearpage
\SlideHeader{Lecture 7 - Look-ahead Methods}{018}{Step 4 and 5: Supplier selection and final assignment}
\begin{minipage}[t]{0.405\textwidth}
\SlideImage{18}{../Materials/2026/3DM_07_Lookahead-Methods_SS26.pdf}
\begin{adjustbox}{max totalsize={\textwidth}{0.43\textheight},center}
\begin{minipage}{\textwidth}\scriptsize
\NoteSection{日本語訳}\begin{itemize}
\item Step 4 and 5: Supplier selection and final assignment（この用語は講義スライド上の専門語として扱う）
\item The consensus solution is transferred to all suppliers（この用語は講義スライド上の専門語として扱う）
\item The suppliers individually select favorable requests（この用語は講義スライド上の専門語として扱う）
\item Finally, a non-overlapping assignment is derived that is（この用語は講義スライド上の専門語として扱う）
\item hopefully acknowledged by the suppliers（この用語は講義スライド上の専門語として扱う）
\item Results:（この用語は講義スライド上の専門語として扱う）
\item MSA: Our Approach（この用語は講義スライド上の専門語として扱う）
\item EAA: mean value approach（この用語は講義スライド上の専門語として扱う）
\item CR: Offer n closest customers（この用語は講義スライド上の専門語として扱う）
\end{itemize}\NoteSection{試験対策ポイント}\begin{itemize}
\item MSAは複数シナリオを解き、解の共通性から現在決定を選ぶ。
\item Consensus functionは問題構造に依存して設計する。
\end{itemize}
\end{minipage}
\end{adjustbox}
\end{minipage}\hfill
\begin{minipage}[t]{0.58\textwidth}
\begin{adjustbox}{max totalsize={\textwidth}{0.89\textheight},center}
\begin{minipage}{\textwidth}\scriptsize
\NoteSection{解説}このページでは「Step 4 and 5: Supplier selection and final assignment」を理解する。スライド上の表現は短いが、試験では定義、具体例、モデル上の意味、手法の限界まで説明できる必要がある。\par
Look-ahead は、現在の決定を評価するために将来を仮想的に見る方法である。単一シナリオだけを見ると、そのシナリオに過剰適合する危険がある。そこで Multiple Scenario Approach では、複数の将来シナリオをサンプリングし、それぞれを確定的な問題として解く。\par
MSAの基本手順は、複数の将来情報列を生成し、各シナリオでMIPなどの確定的最適化を解き、実際にはまだ見えていない将来情報に依存する部分を取り除き、現在時点で実行可能な部分を比較する、という流れである。\par
Consensus function は、複数シナリオで得られた解の中から代表的な解を選ぶための類似度尺度である。生産スケジューリングなら次に処理するジョブ、VRPなら同じ顧客ペアが同じルートにあるか、TSPならアークや隣接関係が一致するか、施設配置なら選ばれた施設集合が一致するかで測る。\par
MSAの利点は、通常のMIPソルバーを使えることと、複数シナリオに対して極端に悪くない現在決定を選べることである。欠点は、stochasticityを解選択で間接的に扱うだけで、動的な再意思決定そのものを完全にモデル化しているわけではない点である。\par
試験では、MSAを『複数シナリオの平均を取る』だけと書かない。各シナリオの解を作り、その解の共通性や代表性を測り、現在実装すべき決定を選ぶ、という構造を説明する。Consensus functionは問題ごとに設計する必要がある。\par\NoteSection{確認問題と解答}\begin{enumerate}\item \textbf{L07-S018: 「Step 4 and 5: Supplier selection and final assignment」の概念を、定義・具体例・モデル上の意味の順に説明せよ。}\\定義としては、MSAは複数シナリオを解き、解の共通性から現在決定を選ぶ。。具体例では、配送・在庫・フードデリバリーのような現実問題を用い、状態、決定、外生情報、報酬、または情報モデル/意思決定モデルのどこに対応するかを明示する。モデル上の意味として、この概念が目的関数、制約、状態表現、将来価値、または方策選択にどのような影響を与えるかを書く。試験では、用語だけを列挙せず、入力データから意思決定、さらに現実の実装へつながる流れを文章で示すことが重要である。\item \textbf{L07-S018: このスライドの考え方を、講義全体のDDDMプロセスの中に位置づけよ。}\\DDDMプロセスでは、現実世界のデータを集約して情報モデルを作り、意思決定モデルまたは方策で行動を選び、実装後に結果を観測して次の状態へ進む。このスライドの概念は、そのどこか一部だけではなく、データ表現、意思決定、将来不確実性、計算可能性のいずれかに関わる。答案では、まずこの概念がどの段階に属するかを述べ、次に他の段階へどう影響するかを書く。例えば価値関数なら将来価値を現在決定へ戻す役割、FIFOなら旅行時間情報モデルの整合性を保証する役割、CFAなら目的関数を調整して将来に良い行動を誘導する役割である。\end{enumerate}
\end{minipage}
\end{adjustbox}
\end{minipage}


\clearpage
\SlideHeader{Lecture 7 - Look-ahead Methods}{019}{Summary}
\begin{minipage}[t]{0.405\textwidth}
\SlideImage{19}{../Materials/2026/3DM_07_Lookahead-Methods_SS26.pdf}
\begin{adjustbox}{max totalsize={\textwidth}{0.43\textheight},center}
\begin{minipage}{\textwidth}\scriptsize
\NoteSection{日本語訳}\begin{itemize}
\item まとめ
\item Menus can reduce waiting times for customers and suppliers（この用語は講義スライド上の専門語として扱う）
\item Menus are particularly important if supplier behavior is uncertain（この用語は講義スライド上の専門語として扱う）
\item In practice, multiple attributes describing individual drivers are taken into account（この用語は講義スライド上の専門語として扱う）
\item Drawback: requests unfavorable to all suppliers have to be assigned as well（この用語は講義スライド上の専門語として扱う）
\item Frustration → ask for a request selection first and assign other requests later on?（この用語は講義スライド上の専門語として扱う）
\item This matter is subject to investigation in lecture 11 → フードデリバリー
\end{itemize}\NoteSection{試験対策ポイント}\begin{itemize}
\item Policy Function Approximation, Rolling Horizon, CFA, Look-ahead, VFAを何を近似するかで区別する。
\item 各手法の強みと弱みを1セットで説明する。
\end{itemize}
\end{minipage}
\end{adjustbox}
\end{minipage}\hfill
\begin{minipage}[t]{0.58\textwidth}
\begin{adjustbox}{max totalsize={\textwidth}{0.89\textheight},center}
\begin{minipage}{\textwidth}\scriptsize
\NoteSection{解説}このページでは「Summary」を理解する。スライド上の表現は短いが、試験では定義、具体例、モデル上の意味、手法の限界まで説明できる必要がある。\par
Policy Function Approximation は、経験則やルールを直接方策として使う方法である。例えば『在庫が閾値を下回ったら注文する』『最も近いドライバーへ割り当てる』などである。計算は速く解釈しやすいが、複雑な将来影響を十分に扱えない可能性がある。\par
Rolling Horizon Re-Optimization は、現在時点で見えている情報を使って一定期間の最適化問題を解き、最初の一部だけ実行し、次時点でまた解き直す方法である。現実変化へ適応できるが、見えている範囲だけを最適化するため近視眼的になることがある。\par
Cost Function Approximation は、現在の目的関数や制約に補正項を加えて、将来に良い影響を持つ決定を誘導する方法である。例えば、将来需要が多い地域から車両を離さないようにペナルティを入れる。設計は問題依存で、補正項が適切でないと逆効果になる。\par
Look-ahead Methods は、将来シナリオを明示的に考え、現在決定の良し悪しを未来までシミュレーションまたは最適化して評価する方法である。将来を直接見るため直感的だが、計算負荷が大きく、サンプリングやホライズンの選び方に依存する。\par
Value Function Approximation は、状態または後決定状態の将来価値を特徴量とモデルで近似する方法である。学習した価値を使えば、各決定について遠い未来まで毎回シミュレーションしなくても将来影響を考慮できる。しかし良い特徴量と学習データが必要である。\par\NoteSection{確認問題と解答}\begin{enumerate}\item \textbf{L07-S019: 「Summary」の概念を、定義・具体例・モデル上の意味の順に説明せよ。}\\定義としては、Policy Function Approximation, Rolling Horizon, CFA, Look-ahead, VFAを何を近似するかで区別する。。具体例では、配送・在庫・フードデリバリーのような現実問題を用い、状態、決定、外生情報、報酬、または情報モデル/意思決定モデルのどこに対応するかを明示する。モデル上の意味として、この概念が目的関数、制約、状態表現、将来価値、または方策選択にどのような影響を与えるかを書く。試験では、用語だけを列挙せず、入力データから意思決定、さらに現実の実装へつながる流れを文章で示すことが重要である。\item \textbf{L07-S019: このスライドの考え方を、講義全体のDDDMプロセスの中に位置づけよ。}\\DDDMプロセスでは、現実世界のデータを集約して情報モデルを作り、意思決定モデルまたは方策で行動を選び、実装後に結果を観測して次の状態へ進む。このスライドの概念は、そのどこか一部だけではなく、データ表現、意思決定、将来不確実性、計算可能性のいずれかに関わる。答案では、まずこの概念がどの段階に属するかを述べ、次に他の段階へどう影響するかを書く。例えば価値関数なら将来価値を現在決定へ戻す役割、FIFOなら旅行時間情報モデルの整合性を保証する役割、CFAなら目的関数を調整して将来に良い行動を誘導する役割である。\end{enumerate}
\end{minipage}
\end{adjustbox}
\end{minipage}


\clearpage
\SlideHeader{Lecture 7 - Look-ahead Methods}{020}{MSA Summary}
\begin{minipage}[t]{0.405\textwidth}
\SlideImage{20}{../Materials/2026/3DM_07_Lookahead-Methods_SS26.pdf}
\begin{adjustbox}{max totalsize={\textwidth}{0.43\textheight},center}
\begin{minipage}{\textwidth}\scriptsize
\NoteSection{日本語訳}\begin{itemize}
\item MSA Summary（この用語は講義スライド上の専門語として扱う）
\item Multiple scenario approach considers the entire 決定 space.
\item MSA does not capture dynamism explicitly.（この用語は講義スライド上の専門語として扱う）
\item 決定s are derived at once and not subsequently.
\item Controllable 確率的ity:
\item Time windows（この用語は講義スライド上の専門語として扱う）
\item Start times, etc.（この用語は講義スライド上の専門語として扱う）
\item Capture dynamism implicitly（この用語は講義スライド上の専門語として扱う）
\item Static solutions → always on the edge of infeasibility（この用語は講義スライド上の専門語として扱う）
\end{itemize}\NoteSection{試験対策ポイント}\begin{itemize}
\item MSAは複数シナリオを解き、解の共通性から現在決定を選ぶ。
\item Consensus functionは問題構造に依存して設計する。
\end{itemize}
\end{minipage}
\end{adjustbox}
\end{minipage}\hfill
\begin{minipage}[t]{0.58\textwidth}
\begin{adjustbox}{max totalsize={\textwidth}{0.89\textheight},center}
\begin{minipage}{\textwidth}\scriptsize
\NoteSection{解説}このページでは「MSA Summary」を理解する。スライド上の表現は短いが、試験では定義、具体例、モデル上の意味、手法の限界まで説明できる必要がある。\par
Look-ahead は、現在の決定を評価するために将来を仮想的に見る方法である。単一シナリオだけを見ると、そのシナリオに過剰適合する危険がある。そこで Multiple Scenario Approach では、複数の将来シナリオをサンプリングし、それぞれを確定的な問題として解く。\par
MSAの基本手順は、複数の将来情報列を生成し、各シナリオでMIPなどの確定的最適化を解き、実際にはまだ見えていない将来情報に依存する部分を取り除き、現在時点で実行可能な部分を比較する、という流れである。\par
Consensus function は、複数シナリオで得られた解の中から代表的な解を選ぶための類似度尺度である。生産スケジューリングなら次に処理するジョブ、VRPなら同じ顧客ペアが同じルートにあるか、TSPならアークや隣接関係が一致するか、施設配置なら選ばれた施設集合が一致するかで測る。\par
MSAの利点は、通常のMIPソルバーを使えることと、複数シナリオに対して極端に悪くない現在決定を選べることである。欠点は、stochasticityを解選択で間接的に扱うだけで、動的な再意思決定そのものを完全にモデル化しているわけではない点である。\par
試験では、MSAを『複数シナリオの平均を取る』だけと書かない。各シナリオの解を作り、その解の共通性や代表性を測り、現在実装すべき決定を選ぶ、という構造を説明する。Consensus functionは問題ごとに設計する必要がある。\par\NoteSection{確認問題と解答}\begin{enumerate}\item \textbf{L07-S020: 「MSA Summary」の概念を、定義・具体例・モデル上の意味の順に説明せよ。}\\定義としては、MSAは複数シナリオを解き、解の共通性から現在決定を選ぶ。。具体例では、配送・在庫・フードデリバリーのような現実問題を用い、状態、決定、外生情報、報酬、または情報モデル/意思決定モデルのどこに対応するかを明示する。モデル上の意味として、この概念が目的関数、制約、状態表現、将来価値、または方策選択にどのような影響を与えるかを書く。試験では、用語だけを列挙せず、入力データから意思決定、さらに現実の実装へつながる流れを文章で示すことが重要である。\item \textbf{L07-S020: このスライドの考え方を、講義全体のDDDMプロセスの中に位置づけよ。}\\DDDMプロセスでは、現実世界のデータを集約して情報モデルを作り、意思決定モデルまたは方策で行動を選び、実装後に結果を観測して次の状態へ進む。このスライドの概念は、そのどこか一部だけではなく、データ表現、意思決定、将来不確実性、計算可能性のいずれかに関わる。答案では、まずこの概念がどの段階に属するかを述べ、次に他の段階へどう影響するかを書く。例えば価値関数なら将来価値を現在決定へ戻す役割、FIFOなら旅行時間情報モデルの整合性を保証する役割、CFAなら目的関数を調整して将来に良い行動を誘導する役割である。\end{enumerate}
\end{minipage}
\end{adjustbox}
\end{minipage}


\clearpage
\SlideHeader{Lecture 7 - Look-ahead Methods}{021}{Agenda}
\begin{minipage}[t]{0.405\textwidth}
\SlideImage{21}{../Materials/2026/3DM_07_Lookahead-Methods_SS26.pdf}
\begin{adjustbox}{max totalsize={\textwidth}{0.43\textheight},center}
\begin{minipage}{\textwidth}\scriptsize
\NoteSection{日本語訳}\begin{itemize}
\item アジェンダ
\item 複数シナリオアプローチ
\item 後決定状態ロールアウト
\end{itemize}
\end{minipage}
\end{adjustbox}
\end{minipage}\hfill
\begin{minipage}[t]{0.58\textwidth}
\begin{adjustbox}{max totalsize={\textwidth}{0.89\textheight},center}
\begin{minipage}{\textwidth}\scriptsize
\NoteSection{注記}このページは表紙・目次・事務情報・導入画像が中心であるため、無理に確認問題や過去問を付けない。
\end{minipage}
\end{adjustbox}
\end{minipage}


\clearpage
\SlideHeader{Lecture 7 - Look-ahead Methods}{022}{Look-ahead Methods}
\begin{minipage}[t]{0.405\textwidth}
\SlideImage{22}{../Materials/2026/3DM_07_Lookahead-Methods_SS26.pdf}
\begin{adjustbox}{max totalsize={\textwidth}{0.43\textheight},center}
\begin{minipage}{\textwidth}\scriptsize
\NoteSection{日本語訳}\begin{itemize}
\item 先読み Methods
\item 考え方: Build and solve part of the MDP
\item Components（この用語は講義スライド上の専門語として扱う）
\item Information model: Sampling from information model 分布
\item 決定 model: None, Mixed Integer Program or given Base 方策
\item 決定 Points: Limited 先読み Horizon
\item Many different 先読み methods
\item Sampling（この用語は講義スライド上の専門語として扱う）
\item Sampling + ロールアウト of 方策
\end{itemize}\NoteSection{試験対策ポイント}\begin{itemize}
\item MSAは複数シナリオを解き、解の共通性から現在決定を選ぶ。
\item Consensus functionは問題構造に依存して設計する。
\end{itemize}
\end{minipage}
\end{adjustbox}
\end{minipage}\hfill
\begin{minipage}[t]{0.58\textwidth}
\begin{adjustbox}{max totalsize={\textwidth}{0.89\textheight},center}
\begin{minipage}{\textwidth}\scriptsize
\NoteSection{解説}このページでは「Look-ahead Methods」を理解する。スライド上の表現は短いが、試験では定義、具体例、モデル上の意味、手法の限界まで説明できる必要がある。\par
Look-ahead は、現在の決定を評価するために将来を仮想的に見る方法である。単一シナリオだけを見ると、そのシナリオに過剰適合する危険がある。そこで Multiple Scenario Approach では、複数の将来シナリオをサンプリングし、それぞれを確定的な問題として解く。\par
MSAの基本手順は、複数の将来情報列を生成し、各シナリオでMIPなどの確定的最適化を解き、実際にはまだ見えていない将来情報に依存する部分を取り除き、現在時点で実行可能な部分を比較する、という流れである。\par
Consensus function は、複数シナリオで得られた解の中から代表的な解を選ぶための類似度尺度である。生産スケジューリングなら次に処理するジョブ、VRPなら同じ顧客ペアが同じルートにあるか、TSPならアークや隣接関係が一致するか、施設配置なら選ばれた施設集合が一致するかで測る。\par
MSAの利点は、通常のMIPソルバーを使えることと、複数シナリオに対して極端に悪くない現在決定を選べることである。欠点は、stochasticityを解選択で間接的に扱うだけで、動的な再意思決定そのものを完全にモデル化しているわけではない点である。\par
試験では、MSAを『複数シナリオの平均を取る』だけと書かない。各シナリオの解を作り、その解の共通性や代表性を測り、現在実装すべき決定を選ぶ、という構造を説明する。Consensus functionは問題ごとに設計する必要がある。\par\NoteSection{確認問題と解答}\begin{enumerate}\item \textbf{L07-S022: 「Look-ahead Methods」の概念を、定義・具体例・モデル上の意味の順に説明せよ。}\\定義としては、MSAは複数シナリオを解き、解の共通性から現在決定を選ぶ。。具体例では、配送・在庫・フードデリバリーのような現実問題を用い、状態、決定、外生情報、報酬、または情報モデル/意思決定モデルのどこに対応するかを明示する。モデル上の意味として、この概念が目的関数、制約、状態表現、将来価値、または方策選択にどのような影響を与えるかを書く。試験では、用語だけを列挙せず、入力データから意思決定、さらに現実の実装へつながる流れを文章で示すことが重要である。\item \textbf{L07-S022: このスライドの考え方を、講義全体のDDDMプロセスの中に位置づけよ。}\\DDDMプロセスでは、現実世界のデータを集約して情報モデルを作り、意思決定モデルまたは方策で行動を選び、実装後に結果を観測して次の状態へ進む。このスライドの概念は、そのどこか一部だけではなく、データ表現、意思決定、将来不確実性、計算可能性のいずれかに関わる。答案では、まずこの概念がどの段階に属するかを述べ、次に他の段階へどう影響するかを書く。例えば価値関数なら将来価値を現在決定へ戻す役割、FIFOなら旅行時間情報モデルの整合性を保証する役割、CFAなら目的関数を調整して将来に良い行動を誘導する役割である。\end{enumerate}
\end{minipage}
\end{adjustbox}
\end{minipage}


\clearpage
\SlideHeader{Lecture 7 - Look-ahead Methods}{023}{Sampling}
\begin{minipage}[t]{0.405\textwidth}
\SlideImage{23}{../Materials/2026/3DM_07_Lookahead-Methods_SS26.pdf}
\begin{adjustbox}{max totalsize={\textwidth}{0.43\textheight},center}
\begin{minipage}{\textwidth}\scriptsize
\NoteSection{日本語訳}\begin{itemize}
\item Sampling（この用語は講義スライド上の専門語として扱う）
\item 状態 𝑆𝑘 and set of 決定s 𝑋(𝑆𝑘 ) xb xb
\item Sk+1(ω1) Sk+1(ω1)
\item ω1
\item Sxk1
\item ...
\item 𝑥 𝑥 𝑥
\item Post-決定 状態 candidates 𝑆𝑘 1 , 𝑆𝑘 2 , 𝑆𝑘 3 , … x1 ωL Sk+1ω(LωL) xb Sxk+1
\item b
\end{itemize}\NoteSection{試験対策ポイント}\begin{itemize}
\item Rolloutは候補決定後の状態から将来方策をシミュレーションし、現在決定を評価する。
\item simulation horizonは終端効果、計算時間、将来情報の信頼性に依存する。
\end{itemize}
\end{minipage}
\end{adjustbox}
\end{minipage}\hfill
\begin{minipage}[t]{0.58\textwidth}
\begin{adjustbox}{max totalsize={\textwidth}{0.89\textheight},center}
\begin{minipage}{\textwidth}\scriptsize
\NoteSection{解説}このページでは「Sampling」を理解する。スライド上の表現は短いが、試験では定義、具体例、モデル上の意味、手法の限界まで説明できる必要がある。\par
Rollout は、現在の候補決定を選んだ後、将来をある方策でシミュレーションして総報酬を見積もる方法である。単に現在報酬だけを見るのではなく、決定後の状態が将来どのような結果を生むかを評価する点が重要である。\par
Post-decision state rollout では、まず候補決定 x を実行した直後の後決定状態 S\textasciicircum{}x を作る。その後、外生情報をサンプリングし、ベース方策で未来を進める。これを複数回繰り返して平均すれば、その候補決定の期待的な将来性能を近似できる。\par
ホライズンを長くすれば終端効果や長期的影響を考えられるが、計算量が増え、遠い将来の予測は不確実になる。短いホライズンは速いが近視眼的になる可能性がある。在庫やダムのように長期的な蓄積が重要な問題では長めのホライズンが必要になりやすい。食品配送のように注文が短時間で完結する問題では短めのホライズンでも有効な場合がある。\par
Rolloutの欠点は、各候補決定について多くの未来シミュレーションが必要になり計算負荷が大きいことである。また、ベース方策が悪いと評価も悪くなり、サンプリング数が少ないと推定が不安定になる。\par
試験では、Rolloutを将来価値近似やlook-aheadと関連づけて説明する。候補決定を比べる、後決定状態から未来をサンプルする、ベース方策で将来を進める、平均性能で選ぶ、という流れを書くとよい。\par\NoteSection{確認問題と解答}\begin{enumerate}\item \textbf{L07-S023: 「Sampling」の概念を、定義・具体例・モデル上の意味の順に説明せよ。}\\定義としては、Rolloutは候補決定後の状態から将来方策をシミュレーションし、現在決定を評価する。。具体例では、配送・在庫・フードデリバリーのような現実問題を用い、状態、決定、外生情報、報酬、または情報モデル/意思決定モデルのどこに対応するかを明示する。モデル上の意味として、この概念が目的関数、制約、状態表現、将来価値、または方策選択にどのような影響を与えるかを書く。試験では、用語だけを列挙せず、入力データから意思決定、さらに現実の実装へつながる流れを文章で示すことが重要である。\item \textbf{L07-S023: このスライドの考え方を、講義全体のDDDMプロセスの中に位置づけよ。}\\DDDMプロセスでは、現実世界のデータを集約して情報モデルを作り、意思決定モデルまたは方策で行動を選び、実装後に結果を観測して次の状態へ進む。このスライドの概念は、そのどこか一部だけではなく、データ表現、意思決定、将来不確実性、計算可能性のいずれかに関わる。答案では、まずこの概念がどの段階に属するかを述べ、次に他の段階へどう影響するかを書く。例えば価値関数なら将来価値を現在決定へ戻す役割、FIFOなら旅行時間情報モデルの整合性を保証する役割、CFAなら目的関数を調整して将来に良い行動を誘導する役割である。\end{enumerate}
\end{minipage}
\end{adjustbox}
\end{minipage}


\clearpage
\SlideHeader{Lecture 7 - Look-ahead Methods}{024}{Example: Inventory Problem Samples}
\begin{minipage}[t]{0.405\textwidth}
\SlideImage{24}{../Materials/2026/3DM_07_Lookahead-Methods_SS26.pdf}
\begin{adjustbox}{max totalsize={\textwidth}{0.43\textheight},center}
\begin{minipage}{\textwidth}\scriptsize
\NoteSection{日本語訳}\begin{itemize}
\item 例: Inventory Problem Samples
\item 決定s
\item 𝐿
\item ෍ 𝑅(𝑆 𝑥1 , 𝜔𝑙 )
\item 𝑙=1
\item 状態
\item ෍ 𝑅(𝑆 𝑥2 , 𝜔𝑙 )
\item ෍ 𝑅(𝑆 𝑥3 , 𝜔𝑙 )
\item Samples may be generated independently of 決定 x and tested with hindsight
\end{itemize}\NoteSection{試験対策ポイント}\begin{itemize}
\item Policy Function Approximation, Rolling Horizon, CFA, Look-ahead, VFAを何を近似するかで区別する。
\item 各手法の強みと弱みを1セットで説明する。
\end{itemize}
\end{minipage}
\end{adjustbox}
\end{minipage}\hfill
\begin{minipage}[t]{0.58\textwidth}
\begin{adjustbox}{max totalsize={\textwidth}{0.89\textheight},center}
\begin{minipage}{\textwidth}\scriptsize
\NoteSection{解説}このページでは「Example: Inventory Problem Samples」を理解する。スライド上の表現は短いが、試験では定義、具体例、モデル上の意味、手法の限界まで説明できる必要がある。\par
Policy Function Approximation は、経験則やルールを直接方策として使う方法である。例えば『在庫が閾値を下回ったら注文する』『最も近いドライバーへ割り当てる』などである。計算は速く解釈しやすいが、複雑な将来影響を十分に扱えない可能性がある。\par
Rolling Horizon Re-Optimization は、現在時点で見えている情報を使って一定期間の最適化問題を解き、最初の一部だけ実行し、次時点でまた解き直す方法である。現実変化へ適応できるが、見えている範囲だけを最適化するため近視眼的になることがある。\par
Cost Function Approximation は、現在の目的関数や制約に補正項を加えて、将来に良い影響を持つ決定を誘導する方法である。例えば、将来需要が多い地域から車両を離さないようにペナルティを入れる。設計は問題依存で、補正項が適切でないと逆効果になる。\par
Look-ahead Methods は、将来シナリオを明示的に考え、現在決定の良し悪しを未来までシミュレーションまたは最適化して評価する方法である。将来を直接見るため直感的だが、計算負荷が大きく、サンプリングやホライズンの選び方に依存する。\par
Value Function Approximation は、状態または後決定状態の将来価値を特徴量とモデルで近似する方法である。学習した価値を使えば、各決定について遠い未来まで毎回シミュレーションしなくても将来影響を考慮できる。しかし良い特徴量と学習データが必要である。\par\NoteSection{確認問題と解答}\begin{enumerate}\item \textbf{L07-S024: 「Example: Inventory Problem Samples」の概念を、定義・具体例・モデル上の意味の順に説明せよ。}\\定義としては、Policy Function Approximation, Rolling Horizon, CFA, Look-ahead, VFAを何を近似するかで区別する。。具体例では、配送・在庫・フードデリバリーのような現実問題を用い、状態、決定、外生情報、報酬、または情報モデル/意思決定モデルのどこに対応するかを明示する。モデル上の意味として、この概念が目的関数、制約、状態表現、将来価値、または方策選択にどのような影響を与えるかを書く。試験では、用語だけを列挙せず、入力データから意思決定、さらに現実の実装へつながる流れを文章で示すことが重要である。\item \textbf{L07-S024: このスライドの考え方を、講義全体のDDDMプロセスの中に位置づけよ。}\\DDDMプロセスでは、現実世界のデータを集約して情報モデルを作り、意思決定モデルまたは方策で行動を選び、実装後に結果を観測して次の状態へ進む。このスライドの概念は、そのどこか一部だけではなく、データ表現、意思決定、将来不確実性、計算可能性のいずれかに関わる。答案では、まずこの概念がどの段階に属するかを述べ、次に他の段階へどう影響するかを書く。例えば価値関数なら将来価値を現在決定へ戻す役割、FIFOなら旅行時間情報モデルの整合性を保証する役割、CFAなら目的関数を調整して将来に良い行動を誘導する役割である。\end{enumerate}
\end{minipage}
\end{adjustbox}
\end{minipage}


\clearpage
\SlideHeader{Lecture 7 - Look-ahead Methods}{025}{Suitable Application for Sampling Transitions}
\begin{minipage}[t]{0.405\textwidth}
\SlideImage{25}{../Materials/2026/3DM_07_Lookahead-Methods_SS26.pdf}
\begin{adjustbox}{max totalsize={\textwidth}{0.43\textheight},center}
\begin{minipage}{\textwidth}\scriptsize
\NoteSection{日本語訳}\begin{itemize}
\item Suitable Application for Sampling Transitions（この用語は講義スライド上の専門語として扱う）
\item Approximation of 報酬 function (demand, 旅行時間, etc.)
\item Short-term anticipation of 確率的 information
\item Very 動的 problems (short lookahead horizon)
\item Emergency services（この用語は講義スライド上の専門語として扱う）
\item Dial-a-ride（この用語は講義スライド上の専門語として扱う）
\item To evaluate 決定s that are not impacted much by future 決定s
\item Inventory routing in bike sharing systems（この用語は講義スライド上の専門語として扱う）
\item A station is served usually only a few times within the day（この用語は講義スライド上の専門語として扱う）
\end{itemize}\NoteSection{試験対策ポイント}\begin{itemize}
\item Policy Function Approximation, Rolling Horizon, CFA, Look-ahead, VFAを何を近似するかで区別する。
\item 各手法の強みと弱みを1セットで説明する。
\end{itemize}
\end{minipage}
\end{adjustbox}
\end{minipage}\hfill
\begin{minipage}[t]{0.58\textwidth}
\begin{adjustbox}{max totalsize={\textwidth}{0.89\textheight},center}
\begin{minipage}{\textwidth}\scriptsize
\NoteSection{解説}このページでは「Suitable Application for Sampling Transitions」を理解する。スライド上の表現は短いが、試験では定義、具体例、モデル上の意味、手法の限界まで説明できる必要がある。\par
Policy Function Approximation は、経験則やルールを直接方策として使う方法である。例えば『在庫が閾値を下回ったら注文する』『最も近いドライバーへ割り当てる』などである。計算は速く解釈しやすいが、複雑な将来影響を十分に扱えない可能性がある。\par
Rolling Horizon Re-Optimization は、現在時点で見えている情報を使って一定期間の最適化問題を解き、最初の一部だけ実行し、次時点でまた解き直す方法である。現実変化へ適応できるが、見えている範囲だけを最適化するため近視眼的になることがある。\par
Cost Function Approximation は、現在の目的関数や制約に補正項を加えて、将来に良い影響を持つ決定を誘導する方法である。例えば、将来需要が多い地域から車両を離さないようにペナルティを入れる。設計は問題依存で、補正項が適切でないと逆効果になる。\par
Look-ahead Methods は、将来シナリオを明示的に考え、現在決定の良し悪しを未来までシミュレーションまたは最適化して評価する方法である。将来を直接見るため直感的だが、計算負荷が大きく、サンプリングやホライズンの選び方に依存する。\par
Value Function Approximation は、状態または後決定状態の将来価値を特徴量とモデルで近似する方法である。学習した価値を使えば、各決定について遠い未来まで毎回シミュレーションしなくても将来影響を考慮できる。しかし良い特徴量と学習データが必要である。\par\NoteSection{確認問題と解答}\begin{enumerate}\item \textbf{L07-S025: 「Suitable Application for Sampling Transitions」の概念を、定義・具体例・モデル上の意味の順に説明せよ。}\\定義としては、Policy Function Approximation, Rolling Horizon, CFA, Look-ahead, VFAを何を近似するかで区別する。。具体例では、配送・在庫・フードデリバリーのような現実問題を用い、状態、決定、外生情報、報酬、または情報モデル/意思決定モデルのどこに対応するかを明示する。モデル上の意味として、この概念が目的関数、制約、状態表現、将来価値、または方策選択にどのような影響を与えるかを書く。試験では、用語だけを列挙せず、入力データから意思決定、さらに現実の実装へつながる流れを文章で示すことが重要である。\item \textbf{L07-S025: このスライドの考え方を、講義全体のDDDMプロセスの中に位置づけよ。}\\DDDMプロセスでは、現実世界のデータを集約して情報モデルを作り、意思決定モデルまたは方策で行動を選び、実装後に結果を観測して次の状態へ進む。このスライドの概念は、そのどこか一部だけではなく、データ表現、意思決定、将来不確実性、計算可能性のいずれかに関わる。答案では、まずこの概念がどの段階に属するかを述べ、次に他の段階へどう影響するかを書く。例えば価値関数なら将来価値を現在決定へ戻す役割、FIFOなら旅行時間情報モデルの整合性を保証する役割、CFAなら目的関数を調整して将来に良い行動を誘導する役割である。\end{enumerate}
\end{minipage}
\end{adjustbox}
\end{minipage}


\clearpage
\SlideHeader{Lecture 7 - Look-ahead Methods}{026}{Post-Decision Rollout Algorithm (Bertsekas, 1995)}
\begin{minipage}[t]{0.405\textwidth}
\SlideImage{26}{../Materials/2026/3DM_07_Lookahead-Methods_SS26.pdf}
\begin{adjustbox}{max totalsize={\textwidth}{0.43\textheight},center}
\begin{minipage}{\textwidth}\scriptsize
\NoteSection{日本語訳}\begin{itemize}
\item Post-決定 ロールアウト Algorithm (Bertsekas, 1995)
\item 状態 𝑆𝑘 and set of 決定s 𝑋(𝑆𝑘 )
\item 𝑥 𝑥 𝑥
\item Post-決定 状態 candidates 𝑆𝑘 1 , 𝑆𝑘 2 , 𝑆𝑘 3 , …
\item For every 後決定状態 𝑆𝑘𝑥 :
\item Simulate 𝐿 paths (information model realizations + 決定s)
\item 決定s within シミュレーション: (runtime-efficient) base 方策 𝜋𝑏
\item Receive 𝐿 results（この用語は講義スライド上の専門語として扱う）
\item ෠ 𝑘𝑥 ) as average over simulated results（この用語は講義スライド上の専門語として扱う）
\end{itemize}\NoteSection{試験対策ポイント}\begin{itemize}
\item Rolloutは候補決定後の状態から将来方策をシミュレーションし、現在決定を評価する。
\item simulation horizonは終端効果、計算時間、将来情報の信頼性に依存する。
\end{itemize}
\end{minipage}
\end{adjustbox}
\end{minipage}\hfill
\begin{minipage}[t]{0.58\textwidth}
\begin{adjustbox}{max totalsize={\textwidth}{0.89\textheight},center}
\begin{minipage}{\textwidth}\scriptsize
\NoteSection{解説}このページでは「Post-Decision Rollout Algorithm (Bertsekas, 1995)」を理解する。スライド上の表現は短いが、試験では定義、具体例、モデル上の意味、手法の限界まで説明できる必要がある。\par
Rollout は、現在の候補決定を選んだ後、将来をある方策でシミュレーションして総報酬を見積もる方法である。単に現在報酬だけを見るのではなく、決定後の状態が将来どのような結果を生むかを評価する点が重要である。\par
Post-decision state rollout では、まず候補決定 x を実行した直後の後決定状態 S\textasciicircum{}x を作る。その後、外生情報をサンプリングし、ベース方策で未来を進める。これを複数回繰り返して平均すれば、その候補決定の期待的な将来性能を近似できる。\par
ホライズンを長くすれば終端効果や長期的影響を考えられるが、計算量が増え、遠い将来の予測は不確実になる。短いホライズンは速いが近視眼的になる可能性がある。在庫やダムのように長期的な蓄積が重要な問題では長めのホライズンが必要になりやすい。食品配送のように注文が短時間で完結する問題では短めのホライズンでも有効な場合がある。\par
Rolloutの欠点は、各候補決定について多くの未来シミュレーションが必要になり計算負荷が大きいことである。また、ベース方策が悪いと評価も悪くなり、サンプリング数が少ないと推定が不安定になる。\par
試験では、Rolloutを将来価値近似やlook-aheadと関連づけて説明する。候補決定を比べる、後決定状態から未来をサンプルする、ベース方策で将来を進める、平均性能で選ぶ、という流れを書くとよい。\par\NoteSection{確認問題と解答}\begin{enumerate}\item \textbf{L07-S026: 「Post-Decision Rollout Algorithm (Bertsekas, 1995)」の概念を、定義・具体例・モデル上の意味の順に説明せよ。}\\定義としては、Rolloutは候補決定後の状態から将来方策をシミュレーションし、現在決定を評価する。。具体例では、配送・在庫・フードデリバリーのような現実問題を用い、状態、決定、外生情報、報酬、または情報モデル/意思決定モデルのどこに対応するかを明示する。モデル上の意味として、この概念が目的関数、制約、状態表現、将来価値、または方策選択にどのような影響を与えるかを書く。試験では、用語だけを列挙せず、入力データから意思決定、さらに現実の実装へつながる流れを文章で示すことが重要である。\item \textbf{L07-S026: このスライドの考え方を、講義全体のDDDMプロセスの中に位置づけよ。}\\DDDMプロセスでは、現実世界のデータを集約して情報モデルを作り、意思決定モデルまたは方策で行動を選び、実装後に結果を観測して次の状態へ進む。このスライドの概念は、そのどこか一部だけではなく、データ表現、意思決定、将来不確実性、計算可能性のいずれかに関わる。答案では、まずこの概念がどの段階に属するかを述べ、次に他の段階へどう影響するかを書く。例えば価値関数なら将来価値を現在決定へ戻す役割、FIFOなら旅行時間情報モデルの整合性を保証する役割、CFAなら目的関数を調整して将来に良い行動を誘導する役割である。\end{enumerate}
\end{minipage}
\end{adjustbox}
\end{minipage}


\clearpage
\SlideHeader{Lecture 7 - Look-ahead Methods}{027}{Post-Decision Rollout Algorithm: Visualization}
\begin{minipage}[t]{0.405\textwidth}
\SlideImage{27}{../Materials/2026/3DM_07_Lookahead-Methods_SS26.pdf}
\begin{adjustbox}{max totalsize={\textwidth}{0.43\textheight},center}
\begin{minipage}{\textwidth}\scriptsize
\NoteSection{日本語訳}\begin{itemize}
\item Post-決定 ロールアウト Algorithm: Visualization
\item xb xb ω1 xb ω1
\item Sk+1(ω1) Sk+1(ω1) Sk+2(ω1) ... SK(ω1)
\item ω1
\item Sxk1
\item ...
\item x1 ωL Sk+1ω(LωL) xb Sxk+1
\item b
\item (ωL) ωL
\end{itemize}\NoteSection{試験対策ポイント}\begin{itemize}
\item Rolloutは候補決定後の状態から将来方策をシミュレーションし、現在決定を評価する。
\item simulation horizonは終端効果、計算時間、将来情報の信頼性に依存する。
\end{itemize}
\end{minipage}
\end{adjustbox}
\end{minipage}\hfill
\begin{minipage}[t]{0.58\textwidth}
\begin{adjustbox}{max totalsize={\textwidth}{0.89\textheight},center}
\begin{minipage}{\textwidth}\scriptsize
\NoteSection{解説}このページでは「Post-Decision Rollout Algorithm: Visualization」を理解する。スライド上の表現は短いが、試験では定義、具体例、モデル上の意味、手法の限界まで説明できる必要がある。\par
Rollout は、現在の候補決定を選んだ後、将来をある方策でシミュレーションして総報酬を見積もる方法である。単に現在報酬だけを見るのではなく、決定後の状態が将来どのような結果を生むかを評価する点が重要である。\par
Post-decision state rollout では、まず候補決定 x を実行した直後の後決定状態 S\textasciicircum{}x を作る。その後、外生情報をサンプリングし、ベース方策で未来を進める。これを複数回繰り返して平均すれば、その候補決定の期待的な将来性能を近似できる。\par
ホライズンを長くすれば終端効果や長期的影響を考えられるが、計算量が増え、遠い将来の予測は不確実になる。短いホライズンは速いが近視眼的になる可能性がある。在庫やダムのように長期的な蓄積が重要な問題では長めのホライズンが必要になりやすい。食品配送のように注文が短時間で完結する問題では短めのホライズンでも有効な場合がある。\par
Rolloutの欠点は、各候補決定について多くの未来シミュレーションが必要になり計算負荷が大きいことである。また、ベース方策が悪いと評価も悪くなり、サンプリング数が少ないと推定が不安定になる。\par
試験では、Rolloutを将来価値近似やlook-aheadと関連づけて説明する。候補決定を比べる、後決定状態から未来をサンプルする、ベース方策で将来を進める、平均性能で選ぶ、という流れを書くとよい。\par\NoteSection{確認問題と解答}\begin{enumerate}\item \textbf{L07-S027: 「Post-Decision Rollout Algorithm: Visualization」の概念を、定義・具体例・モデル上の意味の順に説明せよ。}\\定義としては、Rolloutは候補決定後の状態から将来方策をシミュレーションし、現在決定を評価する。。具体例では、配送・在庫・フードデリバリーのような現実問題を用い、状態、決定、外生情報、報酬、または情報モデル/意思決定モデルのどこに対応するかを明示する。モデル上の意味として、この概念が目的関数、制約、状態表現、将来価値、または方策選択にどのような影響を与えるかを書く。試験では、用語だけを列挙せず、入力データから意思決定、さらに現実の実装へつながる流れを文章で示すことが重要である。\item \textbf{L07-S027: このスライドの考え方を、講義全体のDDDMプロセスの中に位置づけよ。}\\DDDMプロセスでは、現実世界のデータを集約して情報モデルを作り、意思決定モデルまたは方策で行動を選び、実装後に結果を観測して次の状態へ進む。このスライドの概念は、そのどこか一部だけではなく、データ表現、意思決定、将来不確実性、計算可能性のいずれかに関わる。答案では、まずこの概念がどの段階に属するかを述べ、次に他の段階へどう影響するかを書く。例えば価値関数なら将来価値を現在決定へ戻す役割、FIFOなら旅行時間情報モデルの整合性を保証する役割、CFAなら目的関数を調整して将来に良い行動を誘導する役割である。\end{enumerate}
\end{minipage}
\end{adjustbox}
\end{minipage}


\clearpage
\SlideHeader{Lecture 7 - Look-ahead Methods}{028}{𝑋𝑘𝜋 (𝑆𝑘 ) = arg min 𝑅(𝑆𝑘 , 𝑥) + 𝑉(𝑆𝑘𝑥 )}
\begin{minipage}[t]{0.405\textwidth}
\SlideImage{28}{../Materials/2026/3DM_07_Lookahead-Methods_SS26.pdf}
\begin{adjustbox}{max totalsize={\textwidth}{0.43\textheight},center}
\begin{minipage}{\textwidth}\scriptsize
\NoteSection{日本語訳}\begin{itemize}
\item ∗
\item 𝑋𝑘𝜋 (𝑆𝑘 ) = arg min 𝑅(𝑆𝑘 , 𝑥) + 𝑉(𝑆𝑘𝑥 )
\item 𝑥∈𝑋(𝑆𝑘 )
\item 例: Inventory Problem
\item “Refill up “Refill up（この用語は講義スライド上の専門語として扱う）
\item to 50\%” to 50\%”
\item 𝑅=0 𝑅1 𝑅=0 𝑅4
\item 決定s
\item 状態
\end{itemize}\NoteSection{試験対策ポイント}\begin{itemize}
\item Policy Function Approximation, Rolling Horizon, CFA, Look-ahead, VFAを何を近似するかで区別する。
\item 各手法の強みと弱みを1セットで説明する。
\end{itemize}
\end{minipage}
\end{adjustbox}
\end{minipage}\hfill
\begin{minipage}[t]{0.58\textwidth}
\begin{adjustbox}{max totalsize={\textwidth}{0.89\textheight},center}
\begin{minipage}{\textwidth}\scriptsize
\NoteSection{解説}このページでは「𝑋𝑘𝜋 (𝑆𝑘 ) = arg min 𝑅(𝑆𝑘 , 𝑥) + 𝑉(𝑆𝑘𝑥 )」を理解する。スライド上の表現は短いが、試験では定義、具体例、モデル上の意味、手法の限界まで説明できる必要がある。\par
Policy Function Approximation は、経験則やルールを直接方策として使う方法である。例えば『在庫が閾値を下回ったら注文する』『最も近いドライバーへ割り当てる』などである。計算は速く解釈しやすいが、複雑な将来影響を十分に扱えない可能性がある。\par
Rolling Horizon Re-Optimization は、現在時点で見えている情報を使って一定期間の最適化問題を解き、最初の一部だけ実行し、次時点でまた解き直す方法である。現実変化へ適応できるが、見えている範囲だけを最適化するため近視眼的になることがある。\par
Cost Function Approximation は、現在の目的関数や制約に補正項を加えて、将来に良い影響を持つ決定を誘導する方法である。例えば、将来需要が多い地域から車両を離さないようにペナルティを入れる。設計は問題依存で、補正項が適切でないと逆効果になる。\par
Look-ahead Methods は、将来シナリオを明示的に考え、現在決定の良し悪しを未来までシミュレーションまたは最適化して評価する方法である。将来を直接見るため直感的だが、計算負荷が大きく、サンプリングやホライズンの選び方に依存する。\par
Value Function Approximation は、状態または後決定状態の将来価値を特徴量とモデルで近似する方法である。学習した価値を使えば、各決定について遠い未来まで毎回シミュレーションしなくても将来影響を考慮できる。しかし良い特徴量と学習データが必要である。\par\NoteSection{確認問題と解答}\begin{enumerate}\item \textbf{L07-S028: 「𝑋𝑘𝜋 (𝑆𝑘 ) = arg min 𝑅(𝑆𝑘 , 𝑥) + 𝑉(𝑆𝑘𝑥 )」の概念を、定義・具体例・モデル上の意味の順に説明せよ。}\\定義としては、Policy Function Approximation, Rolling Horizon, CFA, Look-ahead, VFAを何を近似するかで区別する。。具体例では、配送・在庫・フードデリバリーのような現実問題を用い、状態、決定、外生情報、報酬、または情報モデル/意思決定モデルのどこに対応するかを明示する。モデル上の意味として、この概念が目的関数、制約、状態表現、将来価値、または方策選択にどのような影響を与えるかを書く。試験では、用語だけを列挙せず、入力データから意思決定、さらに現実の実装へつながる流れを文章で示すことが重要である。\item \textbf{L07-S028: このスライドの考え方を、講義全体のDDDMプロセスの中に位置づけよ。}\\DDDMプロセスでは、現実世界のデータを集約して情報モデルを作り、意思決定モデルまたは方策で行動を選び、実装後に結果を観測して次の状態へ進む。このスライドの概念は、そのどこか一部だけではなく、データ表現、意思決定、将来不確実性、計算可能性のいずれかに関わる。答案では、まずこの概念がどの段階に属するかを述べ、次に他の段階へどう影響するかを書く。例えば価値関数なら将来価値を現在決定へ戻す役割、FIFOなら旅行時間情報モデルの整合性を保証する役割、CFAなら目的関数を調整して将来に良い行動を誘導する役割である。\end{enumerate}
\end{minipage}
\end{adjustbox}
\end{minipage}


\clearpage
\SlideHeader{Lecture 7 - Look-ahead Methods}{029}{Summary}
\begin{minipage}[t]{0.405\textwidth}
\SlideImage{29}{../Materials/2026/3DM_07_Lookahead-Methods_SS26.pdf}
\begin{adjustbox}{max totalsize={\textwidth}{0.43\textheight},center}
\begin{minipage}{\textwidth}\scriptsize
\NoteSection{日本語訳}\begin{itemize}
\item まとめ
\item Can capture 動的 development of information and 決定s
\item 利点ous, for example, to capture 状態 developments over time (vehicle
\item movements, job arrivals, etc.)（この用語は講義スライド上の専門語として扱う）
\item Only a few 決定s can be evaluated
\item 決定s significantly impacted by base 方策
\item The longer the horizon, the further away from realistic developments（この用語は講義スライド上の専門語として扱う）
\item Can be combined with runtime efficient PFAs or CFAs as base policies (and value（この用語は講義スライド上の専門語として扱う）
\item function approximation, as we see in 講義 10)
\end{itemize}\NoteSection{試験対策ポイント}\begin{itemize}
\item Policy Function Approximation, Rolling Horizon, CFA, Look-ahead, VFAを何を近似するかで区別する。
\item 各手法の強みと弱みを1セットで説明する。
\end{itemize}
\end{minipage}
\end{adjustbox}
\end{minipage}\hfill
\begin{minipage}[t]{0.58\textwidth}
\begin{adjustbox}{max totalsize={\textwidth}{0.89\textheight},center}
\begin{minipage}{\textwidth}\scriptsize
\NoteSection{解説}このページでは「Summary」を理解する。スライド上の表現は短いが、試験では定義、具体例、モデル上の意味、手法の限界まで説明できる必要がある。\par
Policy Function Approximation は、経験則やルールを直接方策として使う方法である。例えば『在庫が閾値を下回ったら注文する』『最も近いドライバーへ割り当てる』などである。計算は速く解釈しやすいが、複雑な将来影響を十分に扱えない可能性がある。\par
Rolling Horizon Re-Optimization は、現在時点で見えている情報を使って一定期間の最適化問題を解き、最初の一部だけ実行し、次時点でまた解き直す方法である。現実変化へ適応できるが、見えている範囲だけを最適化するため近視眼的になることがある。\par
Cost Function Approximation は、現在の目的関数や制約に補正項を加えて、将来に良い影響を持つ決定を誘導する方法である。例えば、将来需要が多い地域から車両を離さないようにペナルティを入れる。設計は問題依存で、補正項が適切でないと逆効果になる。\par
Look-ahead Methods は、将来シナリオを明示的に考え、現在決定の良し悪しを未来までシミュレーションまたは最適化して評価する方法である。将来を直接見るため直感的だが、計算負荷が大きく、サンプリングやホライズンの選び方に依存する。\par
Value Function Approximation は、状態または後決定状態の将来価値を特徴量とモデルで近似する方法である。学習した価値を使えば、各決定について遠い未来まで毎回シミュレーションしなくても将来影響を考慮できる。しかし良い特徴量と学習データが必要である。\par\NoteSection{確認問題と解答}\begin{enumerate}\item \textbf{L07-S029: 「Summary」の概念を、定義・具体例・モデル上の意味の順に説明せよ。}\\定義としては、Policy Function Approximation, Rolling Horizon, CFA, Look-ahead, VFAを何を近似するかで区別する。。具体例では、配送・在庫・フードデリバリーのような現実問題を用い、状態、決定、外生情報、報酬、または情報モデル/意思決定モデルのどこに対応するかを明示する。モデル上の意味として、この概念が目的関数、制約、状態表現、将来価値、または方策選択にどのような影響を与えるかを書く。試験では、用語だけを列挙せず、入力データから意思決定、さらに現実の実装へつながる流れを文章で示すことが重要である。\item \textbf{L07-S029: このスライドの考え方を、講義全体のDDDMプロセスの中に位置づけよ。}\\DDDMプロセスでは、現実世界のデータを集約して情報モデルを作り、意思決定モデルまたは方策で行動を選び、実装後に結果を観測して次の状態へ進む。このスライドの概念は、そのどこか一部だけではなく、データ表現、意思決定、将来不確実性、計算可能性のいずれかに関わる。答案では、まずこの概念がどの段階に属するかを述べ、次に他の段階へどう影響するかを書く。例えば価値関数なら将来価値を現在決定へ戻す役割、FIFOなら旅行時間情報モデルの整合性を保証する役割、CFAなら目的関数を調整して将来に良い行動を誘導する役割である。\end{enumerate}
\end{minipage}
\end{adjustbox}
\end{minipage}


\clearpage
\SlideHeader{Lecture 7 - Look-ahead Methods}{030}{Agenda}
\begin{minipage}[t]{0.405\textwidth}
\SlideImage{30}{../Materials/2026/3DM_07_Lookahead-Methods_SS26.pdf}
\begin{adjustbox}{max totalsize={\textwidth}{0.43\textheight},center}
\begin{minipage}{\textwidth}\scriptsize
\NoteSection{日本語訳}\begin{itemize}
\item アジェンダ
\item 先読み Methods
\item Case Study I: Post-決定 状態 ロールアウト Handelsblatt.de
\item Case Study II: 複数シナリオアプローチ
\end{itemize}
\end{minipage}
\end{adjustbox}
\end{minipage}\hfill
\begin{minipage}[t]{0.58\textwidth}
\begin{adjustbox}{max totalsize={\textwidth}{0.89\textheight},center}
\begin{minipage}{\textwidth}\scriptsize
\NoteSection{注記}このページは表紙・目次・事務情報・導入画像が中心であるため、無理に確認問題や過去問を付けない。
\end{minipage}
\end{adjustbox}
\end{minipage}


\clearpage
\SlideHeader{Lecture 7 - Look-ahead Methods}{031}{Case Study: Courier Service}
\begin{minipage}[t]{0.405\textwidth}
\SlideImage{31}{../Materials/2026/3DM_07_Lookahead-Methods_SS26.pdf}
\begin{adjustbox}{max totalsize={\textwidth}{0.43\textheight},center}
\begin{minipage}{\textwidth}\scriptsize
\NoteSection{日本語訳}\begin{itemize}
\item Case Study: Courier Service（この用語は講義スライド上の専門語として扱う）
\item Single vehicle（この用語は講義スライド上の専門語として扱う）
\item Depot (start and finish)（この用語は講義スライド上の専門語として扱う）
\item Limited 旅行時間 (shift)
\item Customer request pickup of parcels（この用語は講義スライド上の専門語として扱う）
\item Part of the customers is known in advance（この用語は講義スライド上の専門語として扱う）
\item Part of the customer request 動的ally
\item throughout the day（この用語は講義スライド上の専門語として扱う）
\item 決定s:
\end{itemize}\NoteSection{試験対策ポイント}\begin{itemize}
\item Policy Function Approximation, Rolling Horizon, CFA, Look-ahead, VFAを何を近似するかで区別する。
\item 各手法の強みと弱みを1セットで説明する。
\end{itemize}
\end{minipage}
\end{adjustbox}
\end{minipage}\hfill
\begin{minipage}[t]{0.58\textwidth}
\begin{adjustbox}{max totalsize={\textwidth}{0.89\textheight},center}
\begin{minipage}{\textwidth}\scriptsize
\NoteSection{解説}このページでは「Case Study: Courier Service」を理解する。スライド上の表現は短いが、試験では定義、具体例、モデル上の意味、手法の限界まで説明できる必要がある。\par
Policy Function Approximation は、経験則やルールを直接方策として使う方法である。例えば『在庫が閾値を下回ったら注文する』『最も近いドライバーへ割り当てる』などである。計算は速く解釈しやすいが、複雑な将来影響を十分に扱えない可能性がある。\par
Rolling Horizon Re-Optimization は、現在時点で見えている情報を使って一定期間の最適化問題を解き、最初の一部だけ実行し、次時点でまた解き直す方法である。現実変化へ適応できるが、見えている範囲だけを最適化するため近視眼的になることがある。\par
Cost Function Approximation は、現在の目的関数や制約に補正項を加えて、将来に良い影響を持つ決定を誘導する方法である。例えば、将来需要が多い地域から車両を離さないようにペナルティを入れる。設計は問題依存で、補正項が適切でないと逆効果になる。\par
Look-ahead Methods は、将来シナリオを明示的に考え、現在決定の良し悪しを未来までシミュレーションまたは最適化して評価する方法である。将来を直接見るため直感的だが、計算負荷が大きく、サンプリングやホライズンの選び方に依存する。\par
Value Function Approximation は、状態または後決定状態の将来価値を特徴量とモデルで近似する方法である。学習した価値を使えば、各決定について遠い未来まで毎回シミュレーションしなくても将来影響を考慮できる。しかし良い特徴量と学習データが必要である。\par\NoteSection{確認問題と解答}\begin{enumerate}\item \textbf{L07-S031: 「Case Study: Courier Service」の概念を、定義・具体例・モデル上の意味の順に説明せよ。}\\定義としては、Policy Function Approximation, Rolling Horizon, CFA, Look-ahead, VFAを何を近似するかで区別する。。具体例では、配送・在庫・フードデリバリーのような現実問題を用い、状態、決定、外生情報、報酬、または情報モデル/意思決定モデルのどこに対応するかを明示する。モデル上の意味として、この概念が目的関数、制約、状態表現、将来価値、または方策選択にどのような影響を与えるかを書く。試験では、用語だけを列挙せず、入力データから意思決定、さらに現実の実装へつながる流れを文章で示すことが重要である。\item \textbf{L07-S031: このスライドの考え方を、講義全体のDDDMプロセスの中に位置づけよ。}\\DDDMプロセスでは、現実世界のデータを集約して情報モデルを作り、意思決定モデルまたは方策で行動を選び、実装後に結果を観測して次の状態へ進む。このスライドの概念は、そのどこか一部だけではなく、データ表現、意思決定、将来不確実性、計算可能性のいずれかに関わる。答案では、まずこの概念がどの段階に属するかを述べ、次に他の段階へどう影響するかを書く。例えば価値関数なら将来価値を現在決定へ戻す役割、FIFOなら旅行時間情報モデルの整合性を保証する役割、CFAなら目的関数を調整して将来に良い行動を誘導する役割である。\end{enumerate}
\end{minipage}
\end{adjustbox}
\end{minipage}


\clearpage
\SlideHeader{Lecture 7 - Look-ahead Methods}{032}{Markov Decision Process}
\begin{minipage}[t]{0.405\textwidth}
\SlideImage{32}{../Materials/2026/3DM_07_Lookahead-Methods_SS26.pdf}
\begin{adjustbox}{max totalsize={\textwidth}{0.43\textheight},center}
\begin{minipage}{\textwidth}\scriptsize
\NoteSection{日本語訳}\begin{itemize}
\item マルコフ意思決定過程
\item 𝑡 = 60 𝑡 = 60 𝑡 = 70 ?
\item ?
\item 決定 Point: New Request
\item 状態: time, vehicle position, customers to serve, new request, planned tour
\item 決定s: Acceptance/Rejection, Adaption of the tour
\item 報酬: 1 (acceptance) or 0 (rejection)
\item Transition: Travel and new customer request（この用語は講義スライド上の専門語として扱う）
\end{itemize}\NoteSection{試験対策ポイント}\begin{itemize}
\item Policy Function Approximation, Rolling Horizon, CFA, Look-ahead, VFAを何を近似するかで区別する。
\item 各手法の強みと弱みを1セットで説明する。
\end{itemize}
\end{minipage}
\end{adjustbox}
\end{minipage}\hfill
\begin{minipage}[t]{0.58\textwidth}
\begin{adjustbox}{max totalsize={\textwidth}{0.89\textheight},center}
\begin{minipage}{\textwidth}\scriptsize
\NoteSection{解説}このページでは「Markov Decision Process」を理解する。スライド上の表現は短いが、試験では定義、具体例、モデル上の意味、手法の限界まで説明できる必要がある。\par
Policy Function Approximation は、経験則やルールを直接方策として使う方法である。例えば『在庫が閾値を下回ったら注文する』『最も近いドライバーへ割り当てる』などである。計算は速く解釈しやすいが、複雑な将来影響を十分に扱えない可能性がある。\par
Rolling Horizon Re-Optimization は、現在時点で見えている情報を使って一定期間の最適化問題を解き、最初の一部だけ実行し、次時点でまた解き直す方法である。現実変化へ適応できるが、見えている範囲だけを最適化するため近視眼的になることがある。\par
Cost Function Approximation は、現在の目的関数や制約に補正項を加えて、将来に良い影響を持つ決定を誘導する方法である。例えば、将来需要が多い地域から車両を離さないようにペナルティを入れる。設計は問題依存で、補正項が適切でないと逆効果になる。\par
Look-ahead Methods は、将来シナリオを明示的に考え、現在決定の良し悪しを未来までシミュレーションまたは最適化して評価する方法である。将来を直接見るため直感的だが、計算負荷が大きく、サンプリングやホライズンの選び方に依存する。\par
Value Function Approximation は、状態または後決定状態の将来価値を特徴量とモデルで近似する方法である。学習した価値を使えば、各決定について遠い未来まで毎回シミュレーションしなくても将来影響を考慮できる。しかし良い特徴量と学習データが必要である。\par\NoteSection{確認問題と解答}\begin{enumerate}\item \textbf{L07-S032: 「Markov Decision Process」の概念を、定義・具体例・モデル上の意味の順に説明せよ。}\\定義としては、Policy Function Approximation, Rolling Horizon, CFA, Look-ahead, VFAを何を近似するかで区別する。。具体例では、配送・在庫・フードデリバリーのような現実問題を用い、状態、決定、外生情報、報酬、または情報モデル/意思決定モデルのどこに対応するかを明示する。モデル上の意味として、この概念が目的関数、制約、状態表現、将来価値、または方策選択にどのような影響を与えるかを書く。試験では、用語だけを列挙せず、入力データから意思決定、さらに現実の実装へつながる流れを文章で示すことが重要である。\item \textbf{L07-S032: このスライドの考え方を、講義全体のDDDMプロセスの中に位置づけよ。}\\DDDMプロセスでは、現実世界のデータを集約して情報モデルを作り、意思決定モデルまたは方策で行動を選び、実装後に結果を観測して次の状態へ進む。このスライドの概念は、そのどこか一部だけではなく、データ表現、意思決定、将来不確実性、計算可能性のいずれかに関わる。答案では、まずこの概念がどの段階に属するかを述べ、次に他の段階へどう影響するかを書く。例えば価値関数なら将来価値を現在決定へ戻す役割、FIFOなら旅行時間情報モデルの整合性を保証する役割、CFAなら目的関数を調整して将来に良い行動を誘導する役割である。\end{enumerate}
\end{minipage}
\end{adjustbox}
\end{minipage}


\clearpage
\SlideHeader{Lecture 7 - Look-ahead Methods}{033}{Rollout Algorithm: Decision Space Reduction}
\begin{minipage}[t]{0.405\textwidth}
\SlideImage{33}{../Materials/2026/3DM_07_Lookahead-Methods_SS26.pdf}
\begin{adjustbox}{max totalsize={\textwidth}{0.43\textheight},center}
\begin{minipage}{\textwidth}\scriptsize
\NoteSection{日本語訳}\begin{itemize}
\item ロールアウト Algorithm: 決定 Space Reduction
\item 𝑡 = 60 𝑡 = 60 𝑡 = 70 ?
\item ?
\item ロールアウト Algorithm can evaluate only a limited number of 決定s
\item Acceptance: 2 決定s
\item Routing: n customers in the tour, 1 request: Up to (n+1)! 決定s
\item Reduction:（この用語は講義スライド上の専門語として扱う）
\item Explicit 決定 on acceptance
\item Routing heuristic: Keep the current tour and insert in the “cheapest” position（この用語は講義スライド上の専門語として扱う）
\end{itemize}\NoteSection{試験対策ポイント}\begin{itemize}
\item Rolloutは候補決定後の状態から将来方策をシミュレーションし、現在決定を評価する。
\item simulation horizonは終端効果、計算時間、将来情報の信頼性に依存する。
\end{itemize}
\end{minipage}
\end{adjustbox}
\end{minipage}\hfill
\begin{minipage}[t]{0.58\textwidth}
\begin{adjustbox}{max totalsize={\textwidth}{0.89\textheight},center}
\begin{minipage}{\textwidth}\scriptsize
\NoteSection{解説}このページでは「Rollout Algorithm: Decision Space Reduction」を理解する。スライド上の表現は短いが、試験では定義、具体例、モデル上の意味、手法の限界まで説明できる必要がある。\par
Rollout は、現在の候補決定を選んだ後、将来をある方策でシミュレーションして総報酬を見積もる方法である。単に現在報酬だけを見るのではなく、決定後の状態が将来どのような結果を生むかを評価する点が重要である。\par
Post-decision state rollout では、まず候補決定 x を実行した直後の後決定状態 S\textasciicircum{}x を作る。その後、外生情報をサンプリングし、ベース方策で未来を進める。これを複数回繰り返して平均すれば、その候補決定の期待的な将来性能を近似できる。\par
ホライズンを長くすれば終端効果や長期的影響を考えられるが、計算量が増え、遠い将来の予測は不確実になる。短いホライズンは速いが近視眼的になる可能性がある。在庫やダムのように長期的な蓄積が重要な問題では長めのホライズンが必要になりやすい。食品配送のように注文が短時間で完結する問題では短めのホライズンでも有効な場合がある。\par
Rolloutの欠点は、各候補決定について多くの未来シミュレーションが必要になり計算負荷が大きいことである。また、ベース方策が悪いと評価も悪くなり、サンプリング数が少ないと推定が不安定になる。\par
試験では、Rolloutを将来価値近似やlook-aheadと関連づけて説明する。候補決定を比べる、後決定状態から未来をサンプルする、ベース方策で将来を進める、平均性能で選ぶ、という流れを書くとよい。\par\NoteSection{確認問題と解答}\begin{enumerate}\item \textbf{L07-S033: 「Rollout Algorithm: Decision Space Reduction」の概念を、定義・具体例・モデル上の意味の順に説明せよ。}\\定義としては、Rolloutは候補決定後の状態から将来方策をシミュレーションし、現在決定を評価する。。具体例では、配送・在庫・フードデリバリーのような現実問題を用い、状態、決定、外生情報、報酬、または情報モデル/意思決定モデルのどこに対応するかを明示する。モデル上の意味として、この概念が目的関数、制約、状態表現、将来価値、または方策選択にどのような影響を与えるかを書く。試験では、用語だけを列挙せず、入力データから意思決定、さらに現実の実装へつながる流れを文章で示すことが重要である。\item \textbf{L07-S033: このスライドの考え方を、講義全体のDDDMプロセスの中に位置づけよ。}\\DDDMプロセスでは、現実世界のデータを集約して情報モデルを作り、意思決定モデルまたは方策で行動を選び、実装後に結果を観測して次の状態へ進む。このスライドの概念は、そのどこか一部だけではなく、データ表現、意思決定、将来不確実性、計算可能性のいずれかに関わる。答案では、まずこの概念がどの段階に属するかを述べ、次に他の段階へどう影響するかを書く。例えば価値関数なら将来価値を現在決定へ戻す役割、FIFOなら旅行時間情報モデルの整合性を保証する役割、CFAなら目的関数を調整して将来に良い行動を誘導する役割である。\end{enumerate}
\end{minipage}
\end{adjustbox}
\end{minipage}


\clearpage
\SlideHeader{Lecture 7 - Look-ahead Methods}{034}{Acceptance and the Value Function}
\begin{minipage}[t]{0.405\textwidth}
\SlideImage{34}{../Materials/2026/3DM_07_Lookahead-Methods_SS26.pdf}
\begin{adjustbox}{max totalsize={\textwidth}{0.43\textheight},center}
\begin{minipage}{\textwidth}\scriptsize
\NoteSection{日本語訳}\begin{itemize}
\item Acceptance and the 価値関数
\item Acceptance of a request reduces the available resources to serve future requests. (value（この用語は講義スライド上の専門語として扱う）
\item decreases)（この用語は講義スライド上の専門語として扱う）
\item Question: Is it worth accepting a request?（この用語は講義スライド上の専門語として扱う）
\item Valuable（この用語は講義スライド上の専門語として扱う）
\item If resource reduction is low,（この用語は講義スライド上の専門語として扱う）
\item If, through resource investment, future（この用語は講義スライド上の専門語として扱う）
\item requests are more efficient to use (coverage)（この用語は講義スライド上の専門語として扱う）
\item ベルマン Equation:
\end{itemize}\NoteSection{試験対策ポイント}\begin{itemize}
\item Policy Function Approximation, Rolling Horizon, CFA, Look-ahead, VFAを何を近似するかで区別する。
\item 各手法の強みと弱みを1セットで説明する。
\end{itemize}
\end{minipage}
\end{adjustbox}
\end{minipage}\hfill
\begin{minipage}[t]{0.58\textwidth}
\begin{adjustbox}{max totalsize={\textwidth}{0.89\textheight},center}
\begin{minipage}{\textwidth}\scriptsize
\NoteSection{解説}このページでは「Acceptance and the Value Function」を理解する。スライド上の表現は短いが、試験では定義、具体例、モデル上の意味、手法の限界まで説明できる必要がある。\par
Policy Function Approximation は、経験則やルールを直接方策として使う方法である。例えば『在庫が閾値を下回ったら注文する』『最も近いドライバーへ割り当てる』などである。計算は速く解釈しやすいが、複雑な将来影響を十分に扱えない可能性がある。\par
Rolling Horizon Re-Optimization は、現在時点で見えている情報を使って一定期間の最適化問題を解き、最初の一部だけ実行し、次時点でまた解き直す方法である。現実変化へ適応できるが、見えている範囲だけを最適化するため近視眼的になることがある。\par
Cost Function Approximation は、現在の目的関数や制約に補正項を加えて、将来に良い影響を持つ決定を誘導する方法である。例えば、将来需要が多い地域から車両を離さないようにペナルティを入れる。設計は問題依存で、補正項が適切でないと逆効果になる。\par
Look-ahead Methods は、将来シナリオを明示的に考え、現在決定の良し悪しを未来までシミュレーションまたは最適化して評価する方法である。将来を直接見るため直感的だが、計算負荷が大きく、サンプリングやホライズンの選び方に依存する。\par
Value Function Approximation は、状態または後決定状態の将来価値を特徴量とモデルで近似する方法である。学習した価値を使えば、各決定について遠い未来まで毎回シミュレーションしなくても将来影響を考慮できる。しかし良い特徴量と学習データが必要である。\par\NoteSection{確認問題と解答}\begin{enumerate}\item \textbf{L07-S034: 「Acceptance and the Value Function」の概念を、定義・具体例・モデル上の意味の順に説明せよ。}\\定義としては、Policy Function Approximation, Rolling Horizon, CFA, Look-ahead, VFAを何を近似するかで区別する。。具体例では、配送・在庫・フードデリバリーのような現実問題を用い、状態、決定、外生情報、報酬、または情報モデル/意思決定モデルのどこに対応するかを明示する。モデル上の意味として、この概念が目的関数、制約、状態表現、将来価値、または方策選択にどのような影響を与えるかを書く。試験では、用語だけを列挙せず、入力データから意思決定、さらに現実の実装へつながる流れを文章で示すことが重要である。\item \textbf{L07-S034: このスライドの考え方を、講義全体のDDDMプロセスの中に位置づけよ。}\\DDDMプロセスでは、現実世界のデータを集約して情報モデルを作り、意思決定モデルまたは方策で行動を選び、実装後に結果を観測して次の状態へ進む。このスライドの概念は、そのどこか一部だけではなく、データ表現、意思決定、将来不確実性、計算可能性のいずれかに関わる。答案では、まずこの概念がどの段階に属するかを述べ、次に他の段階へどう影響するかを書く。例えば価値関数なら将来価値を現在決定へ戻す役割、FIFOなら旅行時間情報モデルの整合性を保証する役割、CFAなら目的関数を調整して将来に良い行動を誘導する役割である。\end{enumerate}
\end{minipage}
\end{adjustbox}
\end{minipage}


\clearpage
\SlideHeader{Lecture 7 - Look-ahead Methods}{035}{𝑋𝑘𝜋 𝑆𝑘 = arg maxx∈𝑋 𝑆𝑘 𝑅 𝑆𝑘 , 𝑥 + 𝑉(𝑆𝑘𝑥 )}
\begin{minipage}[t]{0.405\textwidth}
\SlideImage{35}{../Materials/2026/3DM_07_Lookahead-Methods_SS26.pdf}
\begin{adjustbox}{max totalsize={\textwidth}{0.43\textheight},center}
\begin{minipage}{\textwidth}\scriptsize
\NoteSection{日本語訳}\begin{itemize}
\item ∗
\item 𝑋𝑘𝜋 𝑆𝑘 = arg maxx∈𝑋 𝑆𝑘 𝑅 𝑆𝑘 , 𝑥 + 𝑉(𝑆𝑘𝑥 )（この用語は講義スライド上の専門語として扱う）
\item Acceptance 決定s
\item 2 Policies（この用語は講義スライド上の専門語として扱う）
\item Myopic:（この用語は講義スライド上の専門語として扱う）
\item Value function is set to 0（この用語は講義スライド上の専門語として扱う）
\item Maximize immediate 報酬
\item Customers are always accepted（この用語は講義スライド上の専門語として扱う）
\item as long as they can be served（この用語は講義スライド上の専門語として扱う）
\end{itemize}\NoteSection{試験対策ポイント}\begin{itemize}
\item Rolloutは候補決定後の状態から将来方策をシミュレーションし、現在決定を評価する。
\item simulation horizonは終端効果、計算時間、将来情報の信頼性に依存する。
\end{itemize}
\end{minipage}
\end{adjustbox}
\end{minipage}\hfill
\begin{minipage}[t]{0.58\textwidth}
\begin{adjustbox}{max totalsize={\textwidth}{0.89\textheight},center}
\begin{minipage}{\textwidth}\scriptsize
\NoteSection{解説}このページでは「𝑋𝑘𝜋 𝑆𝑘 = arg maxx∈𝑋 𝑆𝑘 𝑅 𝑆𝑘 , 𝑥 + 𝑉(𝑆𝑘𝑥 )」を理解する。スライド上の表現は短いが、試験では定義、具体例、モデル上の意味、手法の限界まで説明できる必要がある。\par
Rollout は、現在の候補決定を選んだ後、将来をある方策でシミュレーションして総報酬を見積もる方法である。単に現在報酬だけを見るのではなく、決定後の状態が将来どのような結果を生むかを評価する点が重要である。\par
Post-decision state rollout では、まず候補決定 x を実行した直後の後決定状態 S\textasciicircum{}x を作る。その後、外生情報をサンプリングし、ベース方策で未来を進める。これを複数回繰り返して平均すれば、その候補決定の期待的な将来性能を近似できる。\par
ホライズンを長くすれば終端効果や長期的影響を考えられるが、計算量が増え、遠い将来の予測は不確実になる。短いホライズンは速いが近視眼的になる可能性がある。在庫やダムのように長期的な蓄積が重要な問題では長めのホライズンが必要になりやすい。食品配送のように注文が短時間で完結する問題では短めのホライズンでも有効な場合がある。\par
Rolloutの欠点は、各候補決定について多くの未来シミュレーションが必要になり計算負荷が大きいことである。また、ベース方策が悪いと評価も悪くなり、サンプリング数が少ないと推定が不安定になる。\par
試験では、Rolloutを将来価値近似やlook-aheadと関連づけて説明する。候補決定を比べる、後決定状態から未来をサンプルする、ベース方策で将来を進める、平均性能で選ぶ、という流れを書くとよい。\par\NoteSection{確認問題と解答}\begin{enumerate}\item \textbf{L07-S035: 「𝑋𝑘𝜋 𝑆𝑘 = arg maxx∈𝑋 𝑆𝑘 𝑅 𝑆𝑘 , 𝑥 + 𝑉(𝑆𝑘𝑥 )」の概念を、定義・具体例・モデル上の意味の順に説明せよ。}\\定義としては、Rolloutは候補決定後の状態から将来方策をシミュレーションし、現在決定を評価する。。具体例では、配送・在庫・フードデリバリーのような現実問題を用い、状態、決定、外生情報、報酬、または情報モデル/意思決定モデルのどこに対応するかを明示する。モデル上の意味として、この概念が目的関数、制約、状態表現、将来価値、または方策選択にどのような影響を与えるかを書く。試験では、用語だけを列挙せず、入力データから意思決定、さらに現実の実装へつながる流れを文章で示すことが重要である。\item \textbf{L07-S035: このスライドの考え方を、講義全体のDDDMプロセスの中に位置づけよ。}\\DDDMプロセスでは、現実世界のデータを集約して情報モデルを作り、意思決定モデルまたは方策で行動を選び、実装後に結果を観測して次の状態へ進む。このスライドの概念は、そのどこか一部だけではなく、データ表現、意思決定、将来不確実性、計算可能性のいずれかに関わる。答案では、まずこの概念がどの段階に属するかを述べ、次に他の段階へどう影響するかを書く。例えば価値関数なら将来価値を現在決定へ戻す役割、FIFOなら旅行時間情報モデルの整合性を保証する役割、CFAなら目的関数を調整して将来に良い行動を誘導する役割である。\end{enumerate}
\end{minipage}
\end{adjustbox}
\end{minipage}


\clearpage
\SlideHeader{Lecture 7 - Look-ahead Methods}{036}{𝑋𝑘𝜋 𝑆𝑘 = arg maxx∈𝑋 𝑆𝑘 𝑅 𝑆𝑘 , 𝑥 + 𝑉(𝑆𝑘𝑥 )}
\begin{minipage}[t]{0.405\textwidth}
\SlideImage{36}{../Materials/2026/3DM_07_Lookahead-Methods_SS26.pdf}
\begin{adjustbox}{max totalsize={\textwidth}{0.43\textheight},center}
\begin{minipage}{\textwidth}\scriptsize
\NoteSection{日本語訳}\begin{itemize}
\item ∗
\item 𝑋𝑘𝜋 𝑆𝑘 = arg maxx∈𝑋 𝑆𝑘 𝑅 𝑆𝑘 , 𝑥 + 𝑉(𝑆𝑘𝑥 )（この用語は講義スライド上の専門語として扱う）
\item 例
\item 𝑡 = 60
\item 𝑡 = 60 シミュレーション 𝑉෠ = 3
\item シミュレーション 𝑉෠ = 3.5
\end{itemize}\NoteSection{試験対策ポイント}\begin{itemize}
\item Policy Function Approximation, Rolling Horizon, CFA, Look-ahead, VFAを何を近似するかで区別する。
\item 各手法の強みと弱みを1セットで説明する。
\end{itemize}
\end{minipage}
\end{adjustbox}
\end{minipage}\hfill
\begin{minipage}[t]{0.58\textwidth}
\begin{adjustbox}{max totalsize={\textwidth}{0.89\textheight},center}
\begin{minipage}{\textwidth}\scriptsize
\NoteSection{解説}このページでは「𝑋𝑘𝜋 𝑆𝑘 = arg maxx∈𝑋 𝑆𝑘 𝑅 𝑆𝑘 , 𝑥 + 𝑉(𝑆𝑘𝑥 )」を理解する。スライド上の表現は短いが、試験では定義、具体例、モデル上の意味、手法の限界まで説明できる必要がある。\par
Policy Function Approximation は、経験則やルールを直接方策として使う方法である。例えば『在庫が閾値を下回ったら注文する』『最も近いドライバーへ割り当てる』などである。計算は速く解釈しやすいが、複雑な将来影響を十分に扱えない可能性がある。\par
Rolling Horizon Re-Optimization は、現在時点で見えている情報を使って一定期間の最適化問題を解き、最初の一部だけ実行し、次時点でまた解き直す方法である。現実変化へ適応できるが、見えている範囲だけを最適化するため近視眼的になることがある。\par
Cost Function Approximation は、現在の目的関数や制約に補正項を加えて、将来に良い影響を持つ決定を誘導する方法である。例えば、将来需要が多い地域から車両を離さないようにペナルティを入れる。設計は問題依存で、補正項が適切でないと逆効果になる。\par
Look-ahead Methods は、将来シナリオを明示的に考え、現在決定の良し悪しを未来までシミュレーションまたは最適化して評価する方法である。将来を直接見るため直感的だが、計算負荷が大きく、サンプリングやホライズンの選び方に依存する。\par
Value Function Approximation は、状態または後決定状態の将来価値を特徴量とモデルで近似する方法である。学習した価値を使えば、各決定について遠い未来まで毎回シミュレーションしなくても将来影響を考慮できる。しかし良い特徴量と学習データが必要である。\par\NoteSection{確認問題と解答}\begin{enumerate}\item \textbf{L07-S036: 「𝑋𝑘𝜋 𝑆𝑘 = arg maxx∈𝑋 𝑆𝑘 𝑅 𝑆𝑘 , 𝑥 + 𝑉(𝑆𝑘𝑥 )」の概念を、定義・具体例・モデル上の意味の順に説明せよ。}\\定義としては、Policy Function Approximation, Rolling Horizon, CFA, Look-ahead, VFAを何を近似するかで区別する。。具体例では、配送・在庫・フードデリバリーのような現実問題を用い、状態、決定、外生情報、報酬、または情報モデル/意思決定モデルのどこに対応するかを明示する。モデル上の意味として、この概念が目的関数、制約、状態表現、将来価値、または方策選択にどのような影響を与えるかを書く。試験では、用語だけを列挙せず、入力データから意思決定、さらに現実の実装へつながる流れを文章で示すことが重要である。\item \textbf{L07-S036: このスライドの考え方を、講義全体のDDDMプロセスの中に位置づけよ。}\\DDDMプロセスでは、現実世界のデータを集約して情報モデルを作り、意思決定モデルまたは方策で行動を選び、実装後に結果を観測して次の状態へ進む。このスライドの概念は、そのどこか一部だけではなく、データ表現、意思決定、将来不確実性、計算可能性のいずれかに関わる。答案では、まずこの概念がどの段階に属するかを述べ、次に他の段階へどう影響するかを書く。例えば価値関数なら将来価値を現在決定へ戻す役割、FIFOなら旅行時間情報モデルの整合性を保証する役割、CFAなら目的関数を調整して将来に良い行動を誘導する役割である。\end{enumerate}
\end{minipage}
\end{adjustbox}
\end{minipage}


\clearpage
\SlideHeader{Lecture 7 - Look-ahead Methods}{037}{Results}
\begin{minipage}[t]{0.405\textwidth}
\SlideImage{37}{../Materials/2026/3DM_07_Lookahead-Methods_SS26.pdf}
\begin{adjustbox}{max totalsize={\textwidth}{0.43\textheight},center}
\begin{minipage}{\textwidth}\scriptsize
\NoteSection{日本語訳}\begin{itemize}
\item Results（この用語は講義スライド上の専門語として扱う）
\item Acceptances (in \%)（この用語は講義スライド上の専門語として扱う）
\end{itemize}\NoteSection{試験対策ポイント}\begin{itemize}
\item Policy Function Approximation, Rolling Horizon, CFA, Look-ahead, VFAを何を近似するかで区別する。
\item 各手法の強みと弱みを1セットで説明する。
\end{itemize}
\end{minipage}
\end{adjustbox}
\end{minipage}\hfill
\begin{minipage}[t]{0.58\textwidth}
\begin{adjustbox}{max totalsize={\textwidth}{0.89\textheight},center}
\begin{minipage}{\textwidth}\scriptsize
\NoteSection{解説}このページでは「Results」を理解する。スライド上の表現は短いが、試験では定義、具体例、モデル上の意味、手法の限界まで説明できる必要がある。\par
Policy Function Approximation は、経験則やルールを直接方策として使う方法である。例えば『在庫が閾値を下回ったら注文する』『最も近いドライバーへ割り当てる』などである。計算は速く解釈しやすいが、複雑な将来影響を十分に扱えない可能性がある。\par
Rolling Horizon Re-Optimization は、現在時点で見えている情報を使って一定期間の最適化問題を解き、最初の一部だけ実行し、次時点でまた解き直す方法である。現実変化へ適応できるが、見えている範囲だけを最適化するため近視眼的になることがある。\par
Cost Function Approximation は、現在の目的関数や制約に補正項を加えて、将来に良い影響を持つ決定を誘導する方法である。例えば、将来需要が多い地域から車両を離さないようにペナルティを入れる。設計は問題依存で、補正項が適切でないと逆効果になる。\par
Look-ahead Methods は、将来シナリオを明示的に考え、現在決定の良し悪しを未来までシミュレーションまたは最適化して評価する方法である。将来を直接見るため直感的だが、計算負荷が大きく、サンプリングやホライズンの選び方に依存する。\par
Value Function Approximation は、状態または後決定状態の将来価値を特徴量とモデルで近似する方法である。学習した価値を使えば、各決定について遠い未来まで毎回シミュレーションしなくても将来影響を考慮できる。しかし良い特徴量と学習データが必要である。\par\NoteSection{確認問題と解答}\begin{enumerate}\item \textbf{L07-S037: 「Results」の概念を、定義・具体例・モデル上の意味の順に説明せよ。}\\定義としては、Policy Function Approximation, Rolling Horizon, CFA, Look-ahead, VFAを何を近似するかで区別する。。具体例では、配送・在庫・フードデリバリーのような現実問題を用い、状態、決定、外生情報、報酬、または情報モデル/意思決定モデルのどこに対応するかを明示する。モデル上の意味として、この概念が目的関数、制約、状態表現、将来価値、または方策選択にどのような影響を与えるかを書く。試験では、用語だけを列挙せず、入力データから意思決定、さらに現実の実装へつながる流れを文章で示すことが重要である。\item \textbf{L07-S037: このスライドの考え方を、講義全体のDDDMプロセスの中に位置づけよ。}\\DDDMプロセスでは、現実世界のデータを集約して情報モデルを作り、意思決定モデルまたは方策で行動を選び、実装後に結果を観測して次の状態へ進む。このスライドの概念は、そのどこか一部だけではなく、データ表現、意思決定、将来不確実性、計算可能性のいずれかに関わる。答案では、まずこの概念がどの段階に属するかを述べ、次に他の段階へどう影響するかを書く。例えば価値関数なら将来価値を現在決定へ戻す役割、FIFOなら旅行時間情報モデルの整合性を保証する役割、CFAなら目的関数を調整して将来に良い行動を誘導する役割である。\end{enumerate}
\end{minipage}
\end{adjustbox}
\end{minipage}


\clearpage
\SlideHeader{Lecture 7 - Look-ahead Methods}{038}{Tradeoff: Simulation runs vs. quality}
\begin{minipage}[t]{0.405\textwidth}
\SlideImage{38}{../Materials/2026/3DM_07_Lookahead-Methods_SS26.pdf}
\begin{adjustbox}{max totalsize={\textwidth}{0.43\textheight},center}
\begin{minipage}{\textwidth}\scriptsize
\NoteSection{日本語訳}\begin{itemize}
\item Tradeoff: シミュレーション runs vs. quality
\item Runtime and solution quality increase with number of runs（この用語は講義スライド上の専門語として扱う）
\item Accept Runtim（この用語は講義スライド上の専門語として扱う）
\item e
\item Runtime (in seconds)（この用語は講義スライド上の専門語として扱う）
\item Acceptances (in \%)（この用語は講義スライド上の専門語として扱う）
\item シミュレーション Runs
\end{itemize}\NoteSection{試験対策ポイント}\begin{itemize}
\item Policy Function Approximation, Rolling Horizon, CFA, Look-ahead, VFAを何を近似するかで区別する。
\item 各手法の強みと弱みを1セットで説明する。
\end{itemize}
\end{minipage}
\end{adjustbox}
\end{minipage}\hfill
\begin{minipage}[t]{0.58\textwidth}
\begin{adjustbox}{max totalsize={\textwidth}{0.89\textheight},center}
\begin{minipage}{\textwidth}\scriptsize
\NoteSection{解説}このページでは「Tradeoff: Simulation runs vs. quality」を理解する。スライド上の表現は短いが、試験では定義、具体例、モデル上の意味、手法の限界まで説明できる必要がある。\par
Policy Function Approximation は、経験則やルールを直接方策として使う方法である。例えば『在庫が閾値を下回ったら注文する』『最も近いドライバーへ割り当てる』などである。計算は速く解釈しやすいが、複雑な将来影響を十分に扱えない可能性がある。\par
Rolling Horizon Re-Optimization は、現在時点で見えている情報を使って一定期間の最適化問題を解き、最初の一部だけ実行し、次時点でまた解き直す方法である。現実変化へ適応できるが、見えている範囲だけを最適化するため近視眼的になることがある。\par
Cost Function Approximation は、現在の目的関数や制約に補正項を加えて、将来に良い影響を持つ決定を誘導する方法である。例えば、将来需要が多い地域から車両を離さないようにペナルティを入れる。設計は問題依存で、補正項が適切でないと逆効果になる。\par
Look-ahead Methods は、将来シナリオを明示的に考え、現在決定の良し悪しを未来までシミュレーションまたは最適化して評価する方法である。将来を直接見るため直感的だが、計算負荷が大きく、サンプリングやホライズンの選び方に依存する。\par
Value Function Approximation は、状態または後決定状態の将来価値を特徴量とモデルで近似する方法である。学習した価値を使えば、各決定について遠い未来まで毎回シミュレーションしなくても将来影響を考慮できる。しかし良い特徴量と学習データが必要である。\par\NoteSection{確認問題と解答}\begin{enumerate}\item \textbf{L07-S038: 「Tradeoff: Simulation runs vs. quality」の概念を、定義・具体例・モデル上の意味の順に説明せよ。}\\定義としては、Policy Function Approximation, Rolling Horizon, CFA, Look-ahead, VFAを何を近似するかで区別する。。具体例では、配送・在庫・フードデリバリーのような現実問題を用い、状態、決定、外生情報、報酬、または情報モデル/意思決定モデルのどこに対応するかを明示する。モデル上の意味として、この概念が目的関数、制約、状態表現、将来価値、または方策選択にどのような影響を与えるかを書く。試験では、用語だけを列挙せず、入力データから意思決定、さらに現実の実装へつながる流れを文章で示すことが重要である。\item \textbf{L07-S038: このスライドの考え方を、講義全体のDDDMプロセスの中に位置づけよ。}\\DDDMプロセスでは、現実世界のデータを集約して情報モデルを作り、意思決定モデルまたは方策で行動を選び、実装後に結果を観測して次の状態へ進む。このスライドの概念は、そのどこか一部だけではなく、データ表現、意思決定、将来不確実性、計算可能性のいずれかに関わる。答案では、まずこの概念がどの段階に属するかを述べ、次に他の段階へどう影響するかを書く。例えば価値関数なら将来価値を現在決定へ戻す役割、FIFOなら旅行時間情報モデルの整合性を保証する役割、CFAなら目的関数を調整して将来に良い行動を誘導する役割である。\end{enumerate}
\end{minipage}
\end{adjustbox}
\end{minipage}


\clearpage
\SlideHeader{Lecture 7 - Look-ahead Methods}{039}{Limiting the simulation horizon}
\begin{minipage}[t]{0.405\textwidth}
\SlideImage{39}{../Materials/2026/3DM_07_Lookahead-Methods_SS26.pdf}
\begin{adjustbox}{max totalsize={\textwidth}{0.43\textheight},center}
\begin{minipage}{\textwidth}\scriptsize
\NoteSection{日本語訳}\begin{itemize}
\item Limiting the シミュレーション horizon
\item 考え方: limiting the シミュレーション according to H 決定 points
\item Expected: tradeoff between anticipation and reliability（この用語は講義スライド上の専門語として扱う）
\item The longer the シミュレーション horizon, the more anticipation
\item The shorter the シミュレーション horizon, the more reliable the シミュレーション
\item 𝐻
\item sk
\item x0
\item x1
\end{itemize}\NoteSection{試験対策ポイント}\begin{itemize}
\item Policy Function Approximation, Rolling Horizon, CFA, Look-ahead, VFAを何を近似するかで区別する。
\item 各手法の強みと弱みを1セットで説明する。
\end{itemize}
\end{minipage}
\end{adjustbox}
\end{minipage}\hfill
\begin{minipage}[t]{0.58\textwidth}
\begin{adjustbox}{max totalsize={\textwidth}{0.89\textheight},center}
\begin{minipage}{\textwidth}\scriptsize
\NoteSection{解説}このページでは「Limiting the simulation horizon」を理解する。スライド上の表現は短いが、試験では定義、具体例、モデル上の意味、手法の限界まで説明できる必要がある。\par
Policy Function Approximation は、経験則やルールを直接方策として使う方法である。例えば『在庫が閾値を下回ったら注文する』『最も近いドライバーへ割り当てる』などである。計算は速く解釈しやすいが、複雑な将来影響を十分に扱えない可能性がある。\par
Rolling Horizon Re-Optimization は、現在時点で見えている情報を使って一定期間の最適化問題を解き、最初の一部だけ実行し、次時点でまた解き直す方法である。現実変化へ適応できるが、見えている範囲だけを最適化するため近視眼的になることがある。\par
Cost Function Approximation は、現在の目的関数や制約に補正項を加えて、将来に良い影響を持つ決定を誘導する方法である。例えば、将来需要が多い地域から車両を離さないようにペナルティを入れる。設計は問題依存で、補正項が適切でないと逆効果になる。\par
Look-ahead Methods は、将来シナリオを明示的に考え、現在決定の良し悪しを未来までシミュレーションまたは最適化して評価する方法である。将来を直接見るため直感的だが、計算負荷が大きく、サンプリングやホライズンの選び方に依存する。\par
Value Function Approximation は、状態または後決定状態の将来価値を特徴量とモデルで近似する方法である。学習した価値を使えば、各決定について遠い未来まで毎回シミュレーションしなくても将来影響を考慮できる。しかし良い特徴量と学習データが必要である。\par\NoteSection{確認問題と解答}\begin{enumerate}\item \textbf{L07-S039: 「Limiting the simulation horizon」の概念を、定義・具体例・モデル上の意味の順に説明せよ。}\\定義としては、Policy Function Approximation, Rolling Horizon, CFA, Look-ahead, VFAを何を近似するかで区別する。。具体例では、配送・在庫・フードデリバリーのような現実問題を用い、状態、決定、外生情報、報酬、または情報モデル/意思決定モデルのどこに対応するかを明示する。モデル上の意味として、この概念が目的関数、制約、状態表現、将来価値、または方策選択にどのような影響を与えるかを書く。試験では、用語だけを列挙せず、入力データから意思決定、さらに現実の実装へつながる流れを文章で示すことが重要である。\item \textbf{L07-S039: このスライドの考え方を、講義全体のDDDMプロセスの中に位置づけよ。}\\DDDMプロセスでは、現実世界のデータを集約して情報モデルを作り、意思決定モデルまたは方策で行動を選び、実装後に結果を観測して次の状態へ進む。このスライドの概念は、そのどこか一部だけではなく、データ表現、意思決定、将来不確実性、計算可能性のいずれかに関わる。答案では、まずこの概念がどの段階に属するかを述べ、次に他の段階へどう影響するかを書く。例えば価値関数なら将来価値を現在決定へ戻す役割、FIFOなら旅行時間情報モデルの整合性を保証する役割、CFAなら目的関数を調整して将来に良い行動を誘導する役割である。\end{enumerate}
\end{minipage}
\end{adjustbox}
\end{minipage}


\clearpage
\SlideHeader{Lecture 7 - Look-ahead Methods}{040}{Limited Horizon}
\begin{minipage}[t]{0.405\textwidth}
\SlideImage{40}{../Materials/2026/3DM_07_Lookahead-Methods_SS26.pdf}
\begin{adjustbox}{max totalsize={\textwidth}{0.43\textheight},center}
\begin{minipage}{\textwidth}\scriptsize
\NoteSection{日本語訳}\begin{itemize}
\item Limited Horizon（この用語は講義スライド上の専門語として扱う）
\item Adaption of ベルマン Equation:
\item 𝐿 min(𝐾,𝑘+𝐻)
\item 𝜋
\item 𝑉෠ 𝑆𝑘𝑥 = 𝐿−1 ෍ ෍ 𝑅(𝑆𝑖 , 𝑋𝑖 𝑏 𝑆𝑖 ) |𝑆𝑘𝑥 , 𝜔𝑙
\item 𝑙=1 𝑖=𝑘+1
\item シミュレーション of maximally 𝐻 決定 points
\item 𝐻 = 0: myopic (fast), 𝐻 = ∞: full シミュレーション to the final 状態 (slow)
\item Alternative (not considered here): Artificially discounting 報酬s by
\end{itemize}\NoteSection{試験対策ポイント}\begin{itemize}
\item Policy Function Approximation, Rolling Horizon, CFA, Look-ahead, VFAを何を近似するかで区別する。
\item 各手法の強みと弱みを1セットで説明する。
\end{itemize}
\end{minipage}
\end{adjustbox}
\end{minipage}\hfill
\begin{minipage}[t]{0.58\textwidth}
\begin{adjustbox}{max totalsize={\textwidth}{0.89\textheight},center}
\begin{minipage}{\textwidth}\scriptsize
\NoteSection{解説}このページでは「Limited Horizon」を理解する。スライド上の表現は短いが、試験では定義、具体例、モデル上の意味、手法の限界まで説明できる必要がある。\par
Policy Function Approximation は、経験則やルールを直接方策として使う方法である。例えば『在庫が閾値を下回ったら注文する』『最も近いドライバーへ割り当てる』などである。計算は速く解釈しやすいが、複雑な将来影響を十分に扱えない可能性がある。\par
Rolling Horizon Re-Optimization は、現在時点で見えている情報を使って一定期間の最適化問題を解き、最初の一部だけ実行し、次時点でまた解き直す方法である。現実変化へ適応できるが、見えている範囲だけを最適化するため近視眼的になることがある。\par
Cost Function Approximation は、現在の目的関数や制約に補正項を加えて、将来に良い影響を持つ決定を誘導する方法である。例えば、将来需要が多い地域から車両を離さないようにペナルティを入れる。設計は問題依存で、補正項が適切でないと逆効果になる。\par
Look-ahead Methods は、将来シナリオを明示的に考え、現在決定の良し悪しを未来までシミュレーションまたは最適化して評価する方法である。将来を直接見るため直感的だが、計算負荷が大きく、サンプリングやホライズンの選び方に依存する。\par
Value Function Approximation は、状態または後決定状態の将来価値を特徴量とモデルで近似する方法である。学習した価値を使えば、各決定について遠い未来まで毎回シミュレーションしなくても将来影響を考慮できる。しかし良い特徴量と学習データが必要である。\par\NoteSection{確認問題と解答}\begin{enumerate}\item \textbf{L07-S040: 「Limited Horizon」の概念を、定義・具体例・モデル上の意味の順に説明せよ。}\\定義としては、Policy Function Approximation, Rolling Horizon, CFA, Look-ahead, VFAを何を近似するかで区別する。。具体例では、配送・在庫・フードデリバリーのような現実問題を用い、状態、決定、外生情報、報酬、または情報モデル/意思決定モデルのどこに対応するかを明示する。モデル上の意味として、この概念が目的関数、制約、状態表現、将来価値、または方策選択にどのような影響を与えるかを書く。試験では、用語だけを列挙せず、入力データから意思決定、さらに現実の実装へつながる流れを文章で示すことが重要である。\item \textbf{L07-S040: このスライドの考え方を、講義全体のDDDMプロセスの中に位置づけよ。}\\DDDMプロセスでは、現実世界のデータを集約して情報モデルを作り、意思決定モデルまたは方策で行動を選び、実装後に結果を観測して次の状態へ進む。このスライドの概念は、そのどこか一部だけではなく、データ表現、意思決定、将来不確実性、計算可能性のいずれかに関わる。答案では、まずこの概念がどの段階に属するかを述べ、次に他の段階へどう影響するかを書く。例えば価値関数なら将来価値を現在決定へ戻す役割、FIFOなら旅行時間情報モデルの整合性を保証する役割、CFAなら目的関数を調整して将来に良い行動を誘導する役割である。\end{enumerate}
\end{minipage}
\end{adjustbox}
\end{minipage}


\clearpage
\SlideHeader{Lecture 7 - Look-ahead Methods}{041}{Simulation Horizon}
\begin{minipage}[t]{0.405\textwidth}
\SlideImage{41}{../Materials/2026/3DM_07_Lookahead-Methods_SS26.pdf}
\begin{adjustbox}{max totalsize={\textwidth}{0.43\textheight},center}
\begin{minipage}{\textwidth}\scriptsize
\NoteSection{日本語訳}\begin{itemize}
\item シミュレーション Horizon
\item Results for 𝐿 = 2,16,128 and 𝐻 = 0,5,10 … , 80, ∞（この用語は講義スライド上の専門語として扱う）
\item Initial increase in solution quality シミュレーション Runs
\item Slight decrease for longer horizons（この用語は講義スライド上の専門語として扱う）
\item Behavior depends on the number of シミュレーション runs
\item (decrease in runtime for small H)（この用語は講義スライド上の専門語として扱う）
\item Acceptances (in \%) シミュレーション Horizon
\end{itemize}\NoteSection{試験対策ポイント}\begin{itemize}
\item Policy Function Approximation, Rolling Horizon, CFA, Look-ahead, VFAを何を近似するかで区別する。
\item 各手法の強みと弱みを1セットで説明する。
\end{itemize}
\end{minipage}
\end{adjustbox}
\end{minipage}\hfill
\begin{minipage}[t]{0.58\textwidth}
\begin{adjustbox}{max totalsize={\textwidth}{0.89\textheight},center}
\begin{minipage}{\textwidth}\scriptsize
\NoteSection{解説}このページでは「Simulation Horizon」を理解する。スライド上の表現は短いが、試験では定義、具体例、モデル上の意味、手法の限界まで説明できる必要がある。\par
Policy Function Approximation は、経験則やルールを直接方策として使う方法である。例えば『在庫が閾値を下回ったら注文する』『最も近いドライバーへ割り当てる』などである。計算は速く解釈しやすいが、複雑な将来影響を十分に扱えない可能性がある。\par
Rolling Horizon Re-Optimization は、現在時点で見えている情報を使って一定期間の最適化問題を解き、最初の一部だけ実行し、次時点でまた解き直す方法である。現実変化へ適応できるが、見えている範囲だけを最適化するため近視眼的になることがある。\par
Cost Function Approximation は、現在の目的関数や制約に補正項を加えて、将来に良い影響を持つ決定を誘導する方法である。例えば、将来需要が多い地域から車両を離さないようにペナルティを入れる。設計は問題依存で、補正項が適切でないと逆効果になる。\par
Look-ahead Methods は、将来シナリオを明示的に考え、現在決定の良し悪しを未来までシミュレーションまたは最適化して評価する方法である。将来を直接見るため直感的だが、計算負荷が大きく、サンプリングやホライズンの選び方に依存する。\par
Value Function Approximation は、状態または後決定状態の将来価値を特徴量とモデルで近似する方法である。学習した価値を使えば、各決定について遠い未来まで毎回シミュレーションしなくても将来影響を考慮できる。しかし良い特徴量と学習データが必要である。\par\NoteSection{確認問題と解答}\begin{enumerate}\item \textbf{L07-S041: 「Simulation Horizon」の概念を、定義・具体例・モデル上の意味の順に説明せよ。}\\定義としては、Policy Function Approximation, Rolling Horizon, CFA, Look-ahead, VFAを何を近似するかで区別する。。具体例では、配送・在庫・フードデリバリーのような現実問題を用い、状態、決定、外生情報、報酬、または情報モデル/意思決定モデルのどこに対応するかを明示する。モデル上の意味として、この概念が目的関数、制約、状態表現、将来価値、または方策選択にどのような影響を与えるかを書く。試験では、用語だけを列挙せず、入力データから意思決定、さらに現実の実装へつながる流れを文章で示すことが重要である。\item \textbf{L07-S041: このスライドの考え方を、講義全体のDDDMプロセスの中に位置づけよ。}\\DDDMプロセスでは、現実世界のデータを集約して情報モデルを作り、意思決定モデルまたは方策で行動を選び、実装後に結果を観測して次の状態へ進む。このスライドの概念は、そのどこか一部だけではなく、データ表現、意思決定、将来不確実性、計算可能性のいずれかに関わる。答案では、まずこの概念がどの段階に属するかを述べ、次に他の段階へどう影響するかを書く。例えば価値関数なら将来価値を現在決定へ戻す役割、FIFOなら旅行時間情報モデルの整合性を保証する役割、CFAなら目的関数を調整して将来に良い行動を誘導する役割である。\end{enumerate}
\end{minipage}
\end{adjustbox}
\end{minipage}


\clearpage
\SlideHeader{Lecture 7 - Look-ahead Methods}{042}{Summary}
\begin{minipage}[t]{0.405\textwidth}
\SlideImage{42}{../Materials/2026/3DM_07_Lookahead-Methods_SS26.pdf}
\begin{adjustbox}{max totalsize={\textwidth}{0.43\textheight},center}
\begin{minipage}{\textwidth}\scriptsize
\NoteSection{日本語訳}\begin{itemize}
\item まとめ
\item MDP are often too complex to solve exactly（この用語は講義スライド上の専門語として扱う）
\item 考え方: Use of online シミュレーション to estimate the value function
\item 𝐾 𝐿 min(𝐾,𝑘+𝐻)
\item 𝜋∗ 𝜋
\item 𝑉 𝑆𝑘𝑥 = 𝔼 ෍ 𝑅(𝑆𝑖 , 𝑋𝑖 (𝑆𝑖 ))|𝑆𝑘𝑥 𝑉෠ 𝑆𝑘𝑥 = 𝐿−1 ෍ ෍ 𝑅(𝑆𝑖 , 𝑋𝑖 𝑏 𝑆𝑖 ) |𝑆𝑘𝑥 , 𝜔𝑙
\item 𝑖=𝑘+1 𝑙=1 𝑖=𝑘+1
\item シミュレーション is based on realizations and a basic 方策
\item 欠点 of Rollout Algorithm
\end{itemize}\NoteSection{試験対策ポイント}\begin{itemize}
\item Rolloutは候補決定後の状態から将来方策をシミュレーションし、現在決定を評価する。
\item simulation horizonは終端効果、計算時間、将来情報の信頼性に依存する。
\end{itemize}
\end{minipage}
\end{adjustbox}
\end{minipage}\hfill
\begin{minipage}[t]{0.58\textwidth}
\begin{adjustbox}{max totalsize={\textwidth}{0.89\textheight},center}
\begin{minipage}{\textwidth}\scriptsize
\NoteSection{解説}このページでは「Summary」を理解する。スライド上の表現は短いが、試験では定義、具体例、モデル上の意味、手法の限界まで説明できる必要がある。\par
Rollout は、現在の候補決定を選んだ後、将来をある方策でシミュレーションして総報酬を見積もる方法である。単に現在報酬だけを見るのではなく、決定後の状態が将来どのような結果を生むかを評価する点が重要である。\par
Post-decision state rollout では、まず候補決定 x を実行した直後の後決定状態 S\textasciicircum{}x を作る。その後、外生情報をサンプリングし、ベース方策で未来を進める。これを複数回繰り返して平均すれば、その候補決定の期待的な将来性能を近似できる。\par
ホライズンを長くすれば終端効果や長期的影響を考えられるが、計算量が増え、遠い将来の予測は不確実になる。短いホライズンは速いが近視眼的になる可能性がある。在庫やダムのように長期的な蓄積が重要な問題では長めのホライズンが必要になりやすい。食品配送のように注文が短時間で完結する問題では短めのホライズンでも有効な場合がある。\par
Rolloutの欠点は、各候補決定について多くの未来シミュレーションが必要になり計算負荷が大きいことである。また、ベース方策が悪いと評価も悪くなり、サンプリング数が少ないと推定が不安定になる。\par
試験では、Rolloutを将来価値近似やlook-aheadと関連づけて説明する。候補決定を比べる、後決定状態から未来をサンプルする、ベース方策で将来を進める、平均性能で選ぶ、という流れを書くとよい。\par\NoteSection{確認問題と解答}\begin{enumerate}\item \textbf{L07-S042: 「Summary」の概念を、定義・具体例・モデル上の意味の順に説明せよ。}\\定義としては、Rolloutは候補決定後の状態から将来方策をシミュレーションし、現在決定を評価する。。具体例では、配送・在庫・フードデリバリーのような現実問題を用い、状態、決定、外生情報、報酬、または情報モデル/意思決定モデルのどこに対応するかを明示する。モデル上の意味として、この概念が目的関数、制約、状態表現、将来価値、または方策選択にどのような影響を与えるかを書く。試験では、用語だけを列挙せず、入力データから意思決定、さらに現実の実装へつながる流れを文章で示すことが重要である。\item \textbf{L07-S042: このスライドの考え方を、講義全体のDDDMプロセスの中に位置づけよ。}\\DDDMプロセスでは、現実世界のデータを集約して情報モデルを作り、意思決定モデルまたは方策で行動を選び、実装後に結果を観測して次の状態へ進む。このスライドの概念は、そのどこか一部だけではなく、データ表現、意思決定、将来不確実性、計算可能性のいずれかに関わる。答案では、まずこの概念がどの段階に属するかを述べ、次に他の段階へどう影響するかを書く。例えば価値関数なら将来価値を現在決定へ戻す役割、FIFOなら旅行時間情報モデルの整合性を保証する役割、CFAなら目的関数を調整して将来に良い行動を誘導する役割である。\end{enumerate}
\end{minipage}
\end{adjustbox}
\end{minipage}


\clearpage
\SlideHeader{Lecture 7 - Look-ahead Methods}{043}{Look-ahead: Advantages and Disadvantages}
\begin{minipage}[t]{0.405\textwidth}
\SlideImage{43}{../Materials/2026/3DM_07_Lookahead-Methods_SS26.pdf}
\begin{adjustbox}{max totalsize={\textwidth}{0.43\textheight},center}
\begin{minipage}{\textwidth}\scriptsize
\NoteSection{日本語訳}\begin{itemize}
\item 先読み: 利点s and 欠点s
\item 利点:
\item Full 状態 information is considered
\item 確率性 is considered explicitly
\item Full 決定 space XOR
\item Dynamic developments can be considered（この用語は講義スライド上の専門語として扱う）
\item 欠点:
\item Runtime（この用語は講義スライド上の専門語として扱う）
\item Information model: only a few realizations can be considered（この用語は講義スライド上の専門語として扱う）
\end{itemize}\NoteSection{試験対策ポイント}\begin{itemize}
\item Rolloutは候補決定後の状態から将来方策をシミュレーションし、現在決定を評価する。
\item simulation horizonは終端効果、計算時間、将来情報の信頼性に依存する。
\end{itemize}
\end{minipage}
\end{adjustbox}
\end{minipage}\hfill
\begin{minipage}[t]{0.58\textwidth}
\begin{adjustbox}{max totalsize={\textwidth}{0.89\textheight},center}
\begin{minipage}{\textwidth}\scriptsize
\NoteSection{解説}このページでは「Look-ahead: Advantages and Disadvantages」を理解する。スライド上の表現は短いが、試験では定義、具体例、モデル上の意味、手法の限界まで説明できる必要がある。\par
Rollout は、現在の候補決定を選んだ後、将来をある方策でシミュレーションして総報酬を見積もる方法である。単に現在報酬だけを見るのではなく、決定後の状態が将来どのような結果を生むかを評価する点が重要である。\par
Post-decision state rollout では、まず候補決定 x を実行した直後の後決定状態 S\textasciicircum{}x を作る。その後、外生情報をサンプリングし、ベース方策で未来を進める。これを複数回繰り返して平均すれば、その候補決定の期待的な将来性能を近似できる。\par
ホライズンを長くすれば終端効果や長期的影響を考えられるが、計算量が増え、遠い将来の予測は不確実になる。短いホライズンは速いが近視眼的になる可能性がある。在庫やダムのように長期的な蓄積が重要な問題では長めのホライズンが必要になりやすい。食品配送のように注文が短時間で完結する問題では短めのホライズンでも有効な場合がある。\par
Rolloutの欠点は、各候補決定について多くの未来シミュレーションが必要になり計算負荷が大きいことである。また、ベース方策が悪いと評価も悪くなり、サンプリング数が少ないと推定が不安定になる。\par
試験では、Rolloutを将来価値近似やlook-aheadと関連づけて説明する。候補決定を比べる、後決定状態から未来をサンプルする、ベース方策で将来を進める、平均性能で選ぶ、という流れを書くとよい。\par\NoteSection{確認問題と解答}\begin{enumerate}\item \textbf{L07-S043: 「Look-ahead: Advantages and Disadvantages」の概念を、定義・具体例・モデル上の意味の順に説明せよ。}\\定義としては、Rolloutは候補決定後の状態から将来方策をシミュレーションし、現在決定を評価する。。具体例では、配送・在庫・フードデリバリーのような現実問題を用い、状態、決定、外生情報、報酬、または情報モデル/意思決定モデルのどこに対応するかを明示する。モデル上の意味として、この概念が目的関数、制約、状態表現、将来価値、または方策選択にどのような影響を与えるかを書く。試験では、用語だけを列挙せず、入力データから意思決定、さらに現実の実装へつながる流れを文章で示すことが重要である。\item \textbf{L07-S043: このスライドの考え方を、講義全体のDDDMプロセスの中に位置づけよ。}\\DDDMプロセスでは、現実世界のデータを集約して情報モデルを作り、意思決定モデルまたは方策で行動を選び、実装後に結果を観測して次の状態へ進む。このスライドの概念は、そのどこか一部だけではなく、データ表現、意思決定、将来不確実性、計算可能性のいずれかに関わる。答案では、まずこの概念がどの段階に属するかを述べ、次に他の段階へどう影響するかを書く。例えば価値関数なら将来価値を現在決定へ戻す役割、FIFOなら旅行時間情報モデルの整合性を保証する役割、CFAなら目的関数を調整して将来に良い行動を誘導する役割である。\end{enumerate}
\end{minipage}
\end{adjustbox}
\end{minipage}


\clearpage
\SlideHeader{Lecture 7 - Look-ahead Methods}{044}{Back to the Lecture: Outlook}
\begin{minipage}[t]{0.405\textwidth}
\SlideImage{44}{../Materials/2026/3DM_07_Lookahead-Methods_SS26.pdf}
\begin{adjustbox}{max totalsize={\textwidth}{0.43\textheight},center}
\begin{minipage}{\textwidth}\scriptsize
\NoteSection{日本語訳}\begin{itemize}
\item Back to the 講義: Outlook
\item Can we store シミュレーション results?
\item Machine Learning（この用語は講義スライド上の専門語として扱う）
\item 講義 8: Reinforcement Learning
\item 講義 9: Analytics Issues in RL
\end{itemize}\NoteSection{試験対策ポイント}\begin{itemize}
\item Policy Function Approximation, Rolling Horizon, CFA, Look-ahead, VFAを何を近似するかで区別する。
\item 各手法の強みと弱みを1セットで説明する。
\end{itemize}
\end{minipage}
\end{adjustbox}
\end{minipage}\hfill
\begin{minipage}[t]{0.58\textwidth}
\begin{adjustbox}{max totalsize={\textwidth}{0.89\textheight},center}
\begin{minipage}{\textwidth}\scriptsize
\NoteSection{解説}このページでは「Back to the Lecture: Outlook」を理解する。スライド上の表現は短いが、試験では定義、具体例、モデル上の意味、手法の限界まで説明できる必要がある。\par
Policy Function Approximation は、経験則やルールを直接方策として使う方法である。例えば『在庫が閾値を下回ったら注文する』『最も近いドライバーへ割り当てる』などである。計算は速く解釈しやすいが、複雑な将来影響を十分に扱えない可能性がある。\par
Rolling Horizon Re-Optimization は、現在時点で見えている情報を使って一定期間の最適化問題を解き、最初の一部だけ実行し、次時点でまた解き直す方法である。現実変化へ適応できるが、見えている範囲だけを最適化するため近視眼的になることがある。\par
Cost Function Approximation は、現在の目的関数や制約に補正項を加えて、将来に良い影響を持つ決定を誘導する方法である。例えば、将来需要が多い地域から車両を離さないようにペナルティを入れる。設計は問題依存で、補正項が適切でないと逆効果になる。\par
Look-ahead Methods は、将来シナリオを明示的に考え、現在決定の良し悪しを未来までシミュレーションまたは最適化して評価する方法である。将来を直接見るため直感的だが、計算負荷が大きく、サンプリングやホライズンの選び方に依存する。\par
Value Function Approximation は、状態または後決定状態の将来価値を特徴量とモデルで近似する方法である。学習した価値を使えば、各決定について遠い未来まで毎回シミュレーションしなくても将来影響を考慮できる。しかし良い特徴量と学習データが必要である。\par\NoteSection{確認問題と解答}\begin{enumerate}\item \textbf{L07-S044: 「Back to the Lecture: Outlook」の概念を、定義・具体例・モデル上の意味の順に説明せよ。}\\定義としては、Policy Function Approximation, Rolling Horizon, CFA, Look-ahead, VFAを何を近似するかで区別する。。具体例では、配送・在庫・フードデリバリーのような現実問題を用い、状態、決定、外生情報、報酬、または情報モデル/意思決定モデルのどこに対応するかを明示する。モデル上の意味として、この概念が目的関数、制約、状態表現、将来価値、または方策選択にどのような影響を与えるかを書く。試験では、用語だけを列挙せず、入力データから意思決定、さらに現実の実装へつながる流れを文章で示すことが重要である。\item \textbf{L07-S044: このスライドの考え方を、講義全体のDDDMプロセスの中に位置づけよ。}\\DDDMプロセスでは、現実世界のデータを集約して情報モデルを作り、意思決定モデルまたは方策で行動を選び、実装後に結果を観測して次の状態へ進む。このスライドの概念は、そのどこか一部だけではなく、データ表現、意思決定、将来不確実性、計算可能性のいずれかに関わる。答案では、まずこの概念がどの段階に属するかを述べ、次に他の段階へどう影響するかを書く。例えば価値関数なら将来価値を現在決定へ戻す役割、FIFOなら旅行時間情報モデルの整合性を保証する役割、CFAなら目的関数を調整して将来に良い行動を誘導する役割である。\end{enumerate}
\end{minipage}
\end{adjustbox}
\end{minipage}

\clearpage\ExamHeader{Lecture 7 - Look-ahead Methods}
\ExamQuestion{WS22/23 Aufgabe 9 (5 点)}
\ExamSubsection{問題内容}
Powellのpolicy classes。設問では、PFA, CFA, VFA, direct lookahead, policy search等の分類を説明すること。 ことが求められる。\par
\ExamSubsection{満点答案}
満点答案では、PFA、Rolling Horizon、CFA、Look-ahead、VFAを列挙し、それぞれが何を近似するかを説明する。PFAは方策そのもの、CFAは目的関数や制約、Look-aheadは将来シナリオ評価、VFAは将来価値を近似する。各手法の強みと弱みも必ず書く。\par
\ExamSubsection{具体例による解説}
具体例として、配送問題なら、顧客注文・車両位置・旅行時間を情報モデルにし、訪問順や割当を意思決定モデルで決める。在庫問題なら、在庫水準を状態、注文量を決定、需要を外生情報、在庫更新式を遷移、欠品費・保管費を費用として整理する。問題文の文脈に合わせて、抽象用語を必ず具体的な変数や行動に置き換える。\par
\ExamSubsection{採点で落とさない要素}
\begin{itemize}
\item 設問が求める概念の定義を最初に明確に書く。
\item 講義の専門語を列挙するだけでなく、現実例とモデル要素へ対応付ける。
\item 利点だけでなく限界・注意点も最低一つ述べる。
\item 式が出る問題では、記号の意味を文章で説明する。
\item 新しい事例問題では、状態・決定・外生情報・遷移・報酬/費用の順に整理する。
\end{itemize}
\vspace{2mm}\hrule\vspace{2mm}
\ExamQuestion{WS22/23 Aufgabe 10 (3 点)}
\ExamSubsection{問題内容}
Post-Decision Rolloutの欠点。設問では、計算負荷・モデル依存・サンプリング誤差などを説明すること。 ことが求められる。\par
\ExamSubsection{満点答案}
満点答案では、Rolloutが候補決定後の状態から未来をシミュレーションし、現在決定を評価する方法だと説明する。simulation horizonは、終端効果、計算負荷、将来情報の信頼性、問題の時間スケールに依存する。Post-decision rolloutの欠点として計算負荷、ベース方策依存、サンプリング誤差を挙げる。\par
\ExamSubsection{具体例による解説}
具体例として、配送問題なら、顧客注文・車両位置・旅行時間を情報モデルにし、訪問順や割当を意思決定モデルで決める。在庫問題なら、在庫水準を状態、注文量を決定、需要を外生情報、在庫更新式を遷移、欠品費・保管費を費用として整理する。問題文の文脈に合わせて、抽象用語を必ず具体的な変数や行動に置き換える。\par
\ExamSubsection{採点で落とさない要素}
\begin{itemize}
\item 設問が求める概念の定義を最初に明確に書く。
\item 講義の専門語を列挙するだけでなく、現実例とモデル要素へ対応付ける。
\item 利点だけでなく限界・注意点も最低一つ述べる。
\item 式が出る問題では、記号の意味を文章で説明する。
\item 新しい事例問題では、状態・決定・外生情報・遷移・報酬/費用の順に整理する。
\end{itemize}
\vspace{2mm}\hrule\vspace{2mm}
\ExamQuestion{WS22/23 Aufgabe 12 (3 点)}
\ExamSubsection{問題内容}
シミュレーションホライズン設定の考え方。設問では、全計画期間まで見る場合と短期で十分な場合を例で説明すること。 ことが求められる。\par
\ExamSubsection{満点答案}
満点答案では、Real Worldから観測データを集約してInformation Modelを作り、そこからDecision Modelを作り、解をImplementation/Disaggregationで現実の行動へ戻す流れを説明する。Information Modelは現実の圧縮表現であり、Decision Modelは決定変数、目的関数、制約により行動選択を評価する構造である。\par
\ExamSubsection{具体例による解説}
具体例として、配送問題なら、顧客注文・車両位置・旅行時間を情報モデルにし、訪問順や割当を意思決定モデルで決める。在庫問題なら、在庫水準を状態、注文量を決定、需要を外生情報、在庫更新式を遷移、欠品費・保管費を費用として整理する。問題文の文脈に合わせて、抽象用語を必ず具体的な変数や行動に置き換える。\par
\ExamSubsection{採点で落とさない要素}
\begin{itemize}
\item 設問が求める概念の定義を最初に明確に書く。
\item 講義の専門語を列挙するだけでなく、現実例とモデル要素へ対応付ける。
\item 利点だけでなく限界・注意点も最低一つ述べる。
\item 式が出る問題では、記号の意味を文章で説明する。
\item 新しい事例問題では、状態・決定・外生情報・遷移・報酬/費用の順に整理する。
\end{itemize}
\vspace{2mm}\hrule\vspace{2mm}
\ExamQuestion{SS23 Aufgabe 1 (6 点)}
\ExamSubsection{問題内容}
Rollout/VFAのシミュレーションホライズン。設問では、未来シミュレーションをどこまで伸ばすべきか、問題依存で論じること。 ことが求められる。\par
\ExamSubsection{満点答案}
満点答案では、Rolloutが候補決定後の状態から未来をシミュレーションし、現在決定を評価する方法だと説明する。simulation horizonは、終端効果、計算負荷、将来情報の信頼性、問題の時間スケールに依存する。Post-decision rolloutの欠点として計算負荷、ベース方策依存、サンプリング誤差を挙げる。\par
\ExamSubsection{具体例による解説}
具体例として、配送問題なら、顧客注文・車両位置・旅行時間を情報モデルにし、訪問順や割当を意思決定モデルで決める。在庫問題なら、在庫水準を状態、注文量を決定、需要を外生情報、在庫更新式を遷移、欠品費・保管費を費用として整理する。問題文の文脈に合わせて、抽象用語を必ず具体的な変数や行動に置き換える。\par
\ExamSubsection{採点で落とさない要素}
\begin{itemize}
\item 設問が求める概念の定義を最初に明確に書く。
\item 講義の専門語を列挙するだけでなく、現実例とモデル要素へ対応付ける。
\item 利点だけでなく限界・注意点も最低一つ述べる。
\item 式が出る問題では、記号の意味を文章で説明する。
\item 新しい事例問題では、状態・決定・外生情報・遷移・報酬/費用の順に整理する。
\end{itemize}
\vspace{2mm}\hrule\vspace{2mm}
\ExamQuestion{SS23 Aufgabe 2 (5 点)}
\ExamSubsection{問題内容}
Powellのpolicy classes。設問では、政策クラスを列挙し、rule-basedとoptimization-basedの違いも説明すること。 ことが求められる。\par
\ExamSubsection{満点答案}
満点答案では、PFA、Rolling Horizon、CFA、Look-ahead、VFAを列挙し、それぞれが何を近似するかを説明する。PFAは方策そのもの、CFAは目的関数や制約、Look-aheadは将来シナリオ評価、VFAは将来価値を近似する。各手法の強みと弱みも必ず書く。\par
\ExamSubsection{具体例による解説}
具体例として、配送問題なら、顧客注文・車両位置・旅行時間を情報モデルにし、訪問順や割当を意思決定モデルで決める。在庫問題なら、在庫水準を状態、注文量を決定、需要を外生情報、在庫更新式を遷移、欠品費・保管費を費用として整理する。問題文の文脈に合わせて、抽象用語を必ず具体的な変数や行動に置き換える。\par
\ExamSubsection{採点で落とさない要素}
\begin{itemize}
\item 設問が求める概念の定義を最初に明確に書く。
\item 講義の専門語を列挙するだけでなく、現実例とモデル要素へ対応付ける。
\item 利点だけでなく限界・注意点も最低一つ述べる。
\item 式が出る問題では、記号の意味を文章で説明する。
\item 新しい事例問題では、状態・決定・外生情報・遷移・報酬/費用の順に整理する。
\end{itemize}
\vspace{2mm}\hrule\vspace{2mm}
\ExamQuestion{SS23 Aufgabe 4 (3 点)}
\ExamSubsection{問題内容}
Post-Decision Rolloutの欠点。設問では、後決定状態からのシミュレーションの計算・近似上の弱点を説明すること。 ことが求められる。\par
\ExamSubsection{満点答案}
満点答案では、Rolloutが候補決定後の状態から未来をシミュレーションし、現在決定を評価する方法だと説明する。simulation horizonは、終端効果、計算負荷、将来情報の信頼性、問題の時間スケールに依存する。Post-decision rolloutの欠点として計算負荷、ベース方策依存、サンプリング誤差を挙げる。\par
\ExamSubsection{具体例による解説}
具体例として、配送問題なら、顧客注文・車両位置・旅行時間を情報モデルにし、訪問順や割当を意思決定モデルで決める。在庫問題なら、在庫水準を状態、注文量を決定、需要を外生情報、在庫更新式を遷移、欠品費・保管費を費用として整理する。問題文の文脈に合わせて、抽象用語を必ず具体的な変数や行動に置き換える。\par
\ExamSubsection{採点で落とさない要素}
\begin{itemize}
\item 設問が求める概念の定義を最初に明確に書く。
\item 講義の専門語を列挙するだけでなく、現実例とモデル要素へ対応付ける。
\item 利点だけでなく限界・注意点も最低一つ述べる。
\item 式が出る問題では、記号の意味を文章で説明する。
\item 新しい事例問題では、状態・決定・外生情報・遷移・報酬/費用の順に整理する。
\end{itemize}
\vspace{2mm}\hrule\vspace{2mm}
\ExamQuestion{SS23 Aufgabe 5 (8 点)}
\ExamSubsection{問題内容}
レストラン配送のCFA/PFA。設問では、遅延リスクを費用近似または方策近似で扱い、注文統合の判断を説明すること。 ことが求められる。\par
\ExamSubsection{満点答案}
満点答案では、Real Worldから観測データを集約してInformation Modelを作り、そこからDecision Modelを作り、解をImplementation/Disaggregationで現実の行動へ戻す流れを説明する。Information Modelは現実の圧縮表現であり、Decision Modelは決定変数、目的関数、制約により行動選択を評価する構造である。\par
\ExamSubsection{具体例による解説}
具体例として、配送問題なら、顧客注文・車両位置・旅行時間を情報モデルにし、訪問順や割当を意思決定モデルで決める。在庫問題なら、在庫水準を状態、注文量を決定、需要を外生情報、在庫更新式を遷移、欠品費・保管費を費用として整理する。問題文の文脈に合わせて、抽象用語を必ず具体的な変数や行動に置き換える。\par
\ExamSubsection{採点で落とさない要素}
\begin{itemize}
\item 設問が求める概念の定義を最初に明確に書く。
\item 講義の専門語を列挙するだけでなく、現実例とモデル要素へ対応付ける。
\item 利点だけでなく限界・注意点も最低一つ述べる。
\item 式が出る問題では、記号の意味を文章で説明する。
\item 新しい事例問題では、状態・決定・外生情報・遷移・報酬/費用の順に整理する。
\end{itemize}
\vspace{2mm}\hrule\vspace{2mm}
\ExamQuestion{WS23/24 Aufgabe 10 (4 点)}
\ExamSubsection{問題内容}
予測なし動的確率在庫問題で注文を発生させる方法。設問では、終端費用/在庫切れペナルティ等のCFA的な調整として説明すること。 ことが求められる。\par
\ExamSubsection{満点答案}
満点答案では、Real Worldから観測データを集約してInformation Modelを作り、そこからDecision Modelを作り、解をImplementation/Disaggregationで現実の行動へ戻す流れを説明する。Information Modelは現実の圧縮表現であり、Decision Modelは決定変数、目的関数、制約により行動選択を評価する構造である。\par
\ExamSubsection{具体例による解説}
具体例として、配送問題なら、顧客注文・車両位置・旅行時間を情報モデルにし、訪問順や割当を意思決定モデルで決める。在庫問題なら、在庫水準を状態、注文量を決定、需要を外生情報、在庫更新式を遷移、欠品費・保管費を費用として整理する。問題文の文脈に合わせて、抽象用語を必ず具体的な変数や行動に置き換える。\par
\ExamSubsection{採点で落とさない要素}
\begin{itemize}
\item 設問が求める概念の定義を最初に明確に書く。
\item 講義の専門語を列挙するだけでなく、現実例とモデル要素へ対応付ける。
\item 利点だけでなく限界・注意点も最低一つ述べる。
\item 式が出る問題では、記号の意味を文章で説明する。
\item 新しい事例問題では、状態・決定・外生情報・遷移・報酬/費用の順に整理する。
\end{itemize}
\vspace{2mm}\hrule\vspace{2mm}
\ExamQuestion{SS24 Aufgabe 4 (8 点)}
\ExamSubsection{問題内容}
Multiple Scenario ApproachとConsensus Function。設問では、複数シナリオの独立解から一致度を測り、現在決定を抽出する考え方を説明すること。 ことが求められる。\par
\ExamSubsection{満点答案}
満点答案では、Real Worldから観測データを集約してInformation Modelを作り、そこからDecision Modelを作り、解をImplementation/Disaggregationで現実の行動へ戻す流れを説明する。Information Modelは現実の圧縮表現であり、Decision Modelは決定変数、目的関数、制約により行動選択を評価する構造である。\par
\ExamSubsection{具体例による解説}
具体例として、配送問題なら、顧客注文・車両位置・旅行時間を情報モデルにし、訪問順や割当を意思決定モデルで決める。在庫問題なら、在庫水準を状態、注文量を決定、需要を外生情報、在庫更新式を遷移、欠品費・保管費を費用として整理する。問題文の文脈に合わせて、抽象用語を必ず具体的な変数や行動に置き換える。\par
\ExamSubsection{採点で落とさない要素}
\begin{itemize}
\item 設問が求める概念の定義を最初に明確に書く。
\item 講義の専門語を列挙するだけでなく、現実例とモデル要素へ対応付ける。
\item 利点だけでなく限界・注意点も最低一つ述べる。
\item 式が出る問題では、記号の意味を文章で説明する。
\item 新しい事例問題では、状態・決定・外生情報・遷移・報酬/費用の順に整理する。
\end{itemize}
\vspace{2mm}\hrule\vspace{2mm}
\ExamQuestion{WS24/25 Aufgabe 4 (6 点)}
\ExamSubsection{問題内容}
Rollout/VFAのシミュレーションホライズン。設問では、SS23 Q1と同じ論点。ホライズン設定を例とともに説明すること。 ことが求められる。\par
\ExamSubsection{満点答案}
満点答案では、Rolloutが候補決定後の状態から未来をシミュレーションし、現在決定を評価する方法だと説明する。simulation horizonは、終端効果、計算負荷、将来情報の信頼性、問題の時間スケールに依存する。Post-decision rolloutの欠点として計算負荷、ベース方策依存、サンプリング誤差を挙げる。\par
\ExamSubsection{具体例による解説}
具体例として、配送問題なら、顧客注文・車両位置・旅行時間を情報モデルにし、訪問順や割当を意思決定モデルで決める。在庫問題なら、在庫水準を状態、注文量を決定、需要を外生情報、在庫更新式を遷移、欠品費・保管費を費用として整理する。問題文の文脈に合わせて、抽象用語を必ず具体的な変数や行動に置き換える。\par
\ExamSubsection{採点で落とさない要素}
\begin{itemize}
\item 設問が求める概念の定義を最初に明確に書く。
\item 講義の専門語を列挙するだけでなく、現実例とモデル要素へ対応付ける。
\item 利点だけでなく限界・注意点も最低一つ述べる。
\item 式が出る問題では、記号の意味を文章で説明する。
\item 新しい事例問題では、状態・決定・外生情報・遷移・報酬/費用の順に整理する。
\end{itemize}
\vspace{2mm}\hrule\vspace{2mm}
\ExamQuestion{SS25 Exercise 3 (4 点)}
\ExamSubsection{問題内容}
予測なし動的在庫問題で注文を発生させる方法。設問では、WS23/24 Q10と同じ論点。CFA/ペナルティで行動を誘導すること。 ことが求められる。\par
\ExamSubsection{満点答案}
満点答案では、Real Worldから観測データを集約してInformation Modelを作り、そこからDecision Modelを作り、解をImplementation/Disaggregationで現実の行動へ戻す流れを説明する。Information Modelは現実の圧縮表現であり、Decision Modelは決定変数、目的関数、制約により行動選択を評価する構造である。\par
\ExamSubsection{具体例による解説}
具体例として、配送問題なら、顧客注文・車両位置・旅行時間を情報モデルにし、訪問順や割当を意思決定モデルで決める。在庫問題なら、在庫水準を状態、注文量を決定、需要を外生情報、在庫更新式を遷移、欠品費・保管費を費用として整理する。問題文の文脈に合わせて、抽象用語を必ず具体的な変数や行動に置き換える。\par
\ExamSubsection{採点で落とさない要素}
\begin{itemize}
\item 設問が求める概念の定義を最初に明確に書く。
\item 講義の専門語を列挙するだけでなく、現実例とモデル要素へ対応付ける。
\item 利点だけでなく限界・注意点も最低一つ述べる。
\item 式が出る問題では、記号の意味を文章で説明する。
\item 新しい事例問題では、状態・決定・外生情報・遷移・報酬/費用の順に整理する。
\end{itemize}
\vspace{2mm}\hrule\vspace{2mm}
\ExamQuestion{SS25 Exercise 9 (10 点)}
\ExamSubsection{問題内容}
5つのADP method classesと各欠点。設問では、各クラスの機能だけでなく、欠点を1つずつ説明すること。 ことが求められる。\par
\ExamSubsection{満点答案}
満点答案では、PFA、Rolling Horizon、CFA、Look-ahead、VFAを列挙し、それぞれが何を近似するかを説明する。PFAは方策そのもの、CFAは目的関数や制約、Look-aheadは将来シナリオ評価、VFAは将来価値を近似する。各手法の強みと弱みも必ず書く。\par
\ExamSubsection{具体例による解説}
具体例として、配送問題なら、顧客注文・車両位置・旅行時間を情報モデルにし、訪問順や割当を意思決定モデルで決める。在庫問題なら、在庫水準を状態、注文量を決定、需要を外生情報、在庫更新式を遷移、欠品費・保管費を費用として整理する。問題文の文脈に合わせて、抽象用語を必ず具体的な変数や行動に置き換える。\par
\ExamSubsection{採点で落とさない要素}
\begin{itemize}
\item 設問が求める概念の定義を最初に明確に書く。
\item 講義の専門語を列挙するだけでなく、現実例とモデル要素へ対応付ける。
\item 利点だけでなく限界・注意点も最低一つ述べる。
\item 式が出る問題では、記号の意味を文章で説明する。
\item 新しい事例問題では、状態・決定・外生情報・遷移・報酬/費用の順に整理する。
\end{itemize}
\vspace{2mm}\hrule\vspace{2mm}